% Generated extraction from local PDF docs/[篠澤研究室 (よろず)] 请篠泽广来为我们讲解物理学的合志.
\chapter*{请篠泽广来为我们讲解物理学的合志（文字层抽取）}
\addcontentsline{toc}{chapter}{请篠泽广来为我们讲解物理学的合志（文字层抽取）}
\begingroup\small
\noindent\textit{来源：docs/[篠澤研究室 (よろす)] 请篠泽广来为我们讲解物理学的合志 (学園アイトル..pdf；页数：207；可抽取文字页：196。}\par
\endgroup
\medskip
\begingroup
\setlength{\parindent}{0pt}
\setlength{\parskip}{0.45em}

% --- PDF page 1 ---
\noindent\textit{[第 1 页没有可抽取文字层。]}\par

% --- PDF page 2 ---
1\par
卷首语\par
研究背景与概要\par
2024 年 5 月，智能手机游戏「学园偶像大师」开始\par
运营，截至同年 10 月累计下载量已突破 200 万\par
[1]\par
。\par
本作中登场的角色之一\par
泽广，根据官方设定\par
[2]\par
，擅\par
长「全部学业、物理、数学、解谜」，被描述为「对『很\par
辛苦的课程』和『做不好的事情』感到喜悦的怪人。 」她\par
因这种称得上独特的角色个性而广受欢迎。例如现在，在\par
插画交流网站 pixiv 上，她拥有本作担当偶像中最多的插\par
画作品。\par
像这样将擅长学问 ，尤其是物理和数学明确表现出\par
来的角色十分少见，人们期待能有更详细的描写。 然而，\par
在作品中，学业被描绘成对她而言「无聊的事情」\par
[3]\par
，并\par
未积极展开深入挖掘。\par
基于以上背景，本合志以 「让\par
泽广讲解物理学」 为\par
主题，向 11 位供稿者征集了作品。本合志的目的在于，\par
通过挖掘游戏正篇中未曾触及的\par
泽广的学问侧面，加\par
深对角色的进一步理解。同时也旨在尝试将二次创作作\par
为科学传播的一种形式 ，通过\par
泽广来传递学问本身的\par
魅力。\par
期待通过本合志，今后能进一步加深对\par
泽广的理\par
解。\par
[1]https://x.com/gkmas\_official/status\par
/1847932637098447304\par
[2]https://gakuen.idolmaster-\par
official.jp/idol/hiro/\par
[3]亲密度剧情第 5 话\par

% --- PDF page 3 ---
2\par
译者序\par
Bismuth a.k.a. 白发厨敬上。\par
已严肃学习。\par
我并非广 P。\par
我不玩学马仕， 甚至连棉花娃娃都没买过，但是不得不说小广有一种让人难以抗拒\par
的吸引力。像是一种魔性，让人的目光难以转移。\par
1\par
这本超级物理书我实际上早有耳闻， 直到我自己真的开始上手烤的时候才切实感受\par
到了它的恐怖。到底是为什么呢， 开始了这么一个不自量力的目标。 且不说这202 页的\par
内容，我压根就不是学物理的，我是学化学的啊。但是最终这本还是烤完了，我想第二\par
本我可能也会烤的吧。（到底什么时候才会开工就不知道了\par
另外封面的小广真是越看越可爱。\par
我本人不甚了解繁体中文在一些字词上与简体中文的差异， 对于港澳台及海外其他\par
中文区的用词习惯与用词差异我也不甚了解。如果有哪位惯用繁体中文的广 P 希望将\par
这本改为繁体中文版，欢迎来找我获取该本的.docx 版本以及相关图片的.psd 版本。\par
希望被小广套牢的各位都能在这本上学到许多。\par
以下是一些翻译时的吐槽。\par
全文我主要使用的字体是方正 FW 筑紫 A 老明朝，这个字体哪里都好，唯一不好\par
的一点就是缺字。 在其他本子里这个问题还不是很大， 但是在小广这里算是捅了大篓子\par
了——\par
泽广的\par
，这个字在这套字体里是没有的。所以，作为替代方案，只能换成我\par
觉得好看程度稍差一级的思源宋体了，这个哪里都好，就是不如筑紫好看。\par
这次在排版上我也学到许多，希望能继续运用到之后的烤肉当中。\par
ps：word 没有自带的边注真是拉跨啊，卷首语那个我还是靠分栏实现的。\par
pps：另外我这台电脑上没装 EndNote，所以这些脚注里的论文按住 ctrl 点不开是\par
正常现象。\par
1\par
摘自 https://exhentai.org/g/3952995/e8749aa1d2/ 我自己的上传者评论\par

% --- PDF page 4 ---
3\par
目录\par
卷首语 ........................................................................................................... 500mL  1\par
译者序 .......................................................................................................... Bismuth  2\par
放电现象的机制 ................................................................................................ tokio  5\par
想请\par
泽广讲一点量子论的故事 ................................................................... 500mL\par
7\par
需要的是纠缠 \textasciitilde{}用量子信息让\par
泽广验证定域实在论的破缺\textasciitilde{}.................\par
17\par
专属 P 道具：恶魔核心 .........................................................................\par
25\par
选择者 .............................................................................................................. 余命  37\par
请\par
泽广教我锂离子电池 ................................................................................ latte\par
49\par
关于不相溶的东西。或者说是，共晶系统的统计力学正当化。 ......\par
65\par
请\par
泽广教我李群与李代数 .........................................................\par
77\par
Random mode!! ....................................................................................\par
91\par
两位广与巴拿赫-塔斯基定理 ....................................................................\par
99\par
请\par
泽广教教我!狭义相对论 ...................................................................... Alkyne  119\par
后记 .......................................................................................................................... 204\par
版权页 ....................................................................................................................... 205\par

% --- PDF page 5 ---
\noindent\textit{[第 5 页没有可抽取文字层。]}\par

% --- PDF page 6 ---
\noindent\textit{[第 6 页没有可抽取文字层。]}\par

% --- PDF page 7 ---
\noindent\textit{[第 7 页没有可抽取文字层。]}\par

% --- PDF page 8 ---
7\par
想请\par
泽广讲一点量子论的故事\par
500mL\par
June 2024\par
1 导入\par
泽广是偶像。不，虽说是偶像，她的体力却差到连一场演唱会都撑不到最后，质\par
疑她才能的人反而更多， 是个宛如偶像未满的少女。 而且， 在完全自觉没有才能的同时，\par
她却宣称这正是偶像的有趣之处，是个古怪的少女。\par
不知该说「果不其然」还是「理所当然」 ，像她这样的人往往有着罕见的才能。\par
泽同学虽然出生在秋田县秋田市，但四月进入初星学园之前，似乎一直在美国，而且是\par
马萨诸塞州的一所知名大学专攻物理学。 我这辈子还没出过国， 所以她在那边的情形无\par
从揣测。只是，偶尔会看到她坐在图书室的书架前的地板上，安静地读着书，静得让人\par
怀疑她是不是忘了呼吸。 那双琥珀色的眼睛紧紧盯着手边的书， 每隔十几秒就用修长纤\par
细的手指轻轻翻过一页。 她的脸上浮现出一种介于喜悦和紧张之间的微妙表情， 皮肤像\par
上等大理石一样微微泛着苍白的光。每当看到这副模样， 我都会在心里信服地想， 果然\par
如此啊。\par
那是秋天某日的事情。下午两点的空教室里， 比起早晨和傍晚的微寒， 仍留着几分\par
暖意。甚至有点热。我和\par
泽同学约好了请她教我物理。我把「想向年纪比我小的\par
泽\par
同学请教点什么」这种低俗的私心用冠冕堂皇的话包装起来，硬是求她答应了。眼前，\par
她翻开一本素净的米色笔记本， 准备讲课。我盯着她看， 心里想着和平时的她气氛略有\par
不同，真可爱啊。\par
2 两个相反的理论\par
「……开始了。。 」\par
这样宣布之后，\par
泽同学突然向我提问。\par
「首，，制作人。觉觉得『实在』是什么?」\par
「上来就是个学学性问题啊……。 」我稍作思，，试着说， 「我想，是在这个世界上\par
拥有物质性的形体吧。 」\par
于是，她这样回应。\par
「也就是说……只要在这个宙空空间的某处据据一定体积就可以 ?」\par

% --- PDF page 9 ---
8\par
「我是这么想的。 」\par
我这样回答后，她恶作剧似的微微扬起小小的嘴角，这么问道。\par
「那么，如如光就不实在 ?我们能给照亮我们的光定义体积 。制作人。 」\par
「………难啊。 」\par
这实在无法反驳。正在我思，该怎么重新定义的时候，\par
泽同学站起来，用纤细的\par
手指从白板上拿起一支笔，开始讲。\par
「物理学的理论……是言言物理量态的的理论。也就是说，是为了不用做实验就能\par
弄清楚诸如『这颗炮弹会落在哪里』或者『给这个通电会产生多少热量』之类，关于位\par
置和热量等『可测量之物』行为的疑问，而存在的理论。 」\par
她挥着笔一口气说到这里，\par
「然后……过的的物理学是是这样思，物理量的实在的。 」\par
接着说道。板书内容如下。\par
「很直。。比如，动的的应应该具有位置或的量等某种。。这些是实在的，人们曾\par
经认为用方程式言测其行为就是物理学。这种经典的思，方式，被称为 定域实在论\par
1\par
。 」\par
「来来如此。 」\par
「然而， 」她说着，看向我的眼睛。我微微一，，把目光移开，看向白板。\par
「着着物理学进步，出现了新的理论。那就是量子论。至今无法解的的现象得以用\par
这个理论阐明，非常有用。……只不过，量子论主张下面这一点。 」\par
「这个……和才才的主张盾。。量子论等同于在说物理量并不实在。但是，用以往\par
的定域实在论无法说明的现象，用量子论却能说明。那么，究竟哪边是正确的呢。觉怎\par
么看?制作人。 」\par
「………我认为只能实际确认一下，物理量到底是具具有确定的。。 」\par
「正确。 」她略无无趣地继续说道。 「物理学终究只能用实验来定定。该怎么确认物\par
理量的实在呢。 」\par
1\par
「其实也该说明定域性的定义。简略地说，就是遵守『因果律』 。 」\par
定域实在论的主张：\par
所有物理量都是实在的。也就是说，都拥有某种确定的。。\par
量子论的主张：\par
一般地，物理量不具有确定的。。\par

% --- PDF page 10 ---
9\par
3 思想实验\par
「完全没绪………。 」\par
才这么回答，\par
泽同学突然嫣然一笑，这样提议道。\par
「好开心。机会难得，觉可以一直想到想出来为止。。我想看制作人苦恼的样子。 」\par
「我想看\par
泽同学教我的样子。 」\par
「呵呵……不行。 」\par
争辩。我对这个问题没有任何可以成为线索的念绪。直觉告诉我，就算有提示，可\par
能也得抱绪想上一整部电影的时间，所以这里还是想老老实实让她教。\par
「那，在我思，的期间，如果\par
泽同学肯练习单脚站立，我也会努力想。两个人一\par
起苦恼，更开心对吧。 」\par
「………起起来虽然非常开心，但那样我觉得就没法教到最后了。 」\par
「那可就烦了了，这里且且，往下进行吧。 」\par
「……没法法呢。虽然有点行行输，，但有这样一个实验。 」\par
我暗自松了口气，庆幸被巧妙说服了，继续起她讲下的。\par
她花了两三分钟，在白板上写下了这些。这段时间里，我欣赏着她长发摇曳、可爱\par
的文字一点点出现的样子。 教室窗户射入的阳光在\par
泽同学的绪发表面反射的样子， 我\par
甚至觉得犹如朝霞映在湖面。文章的意思嘛……我不懂。\par
「才才的问题， 在实验语下下可以这样换种说法。 『测量。的动的，究竟是源于定域\par
实在论所依据的隐变量，还是量子论所依据的本质不确定性』 」\par
实验流程：\par
1. 在充分远离的地点 A、B 处各有 1 名。测者。\par
2. 在 A 点， 使用测量器的设置$\theta$或$\theta$$\prime$的其中之一测量物理量 A， 测量。记为\par
1 或-1。\par
3. 在 B 点也同样地，使用$\phi$或$\phi$$\prime$的其中之一测量物理量 B，得到测量。。\par
4. 两人各自记录自己所用的设定和测量。。\par
5. 多次重复 2～4 步骤。\par
6. 测量后，将两人的数据对照计算 C。\par
其中\par
C =[缺字]A($\theta$)B($\phi$)[缺字]+ [缺字]A($\theta$$\prime$)B($\phi$)[缺字]-[缺字]A($\theta$)B($\phi$$\prime$)[缺字]+ [缺字]A($\theta$$\prime$)B($\phi$$\prime$)[缺字]\par

% --- PDF page 11 ---
10\par
「什么意思?」\par
「定域实在论中，某个物理量应该具有确定的。，但测量。却动的时，能想到的是\par
条件其实并未完备。 也就是说， 可以解的为存在隐变量， 其。每次都在变化。 与此相对，\par
量子论认为物理量本质上有动的。 」\par
「啊啊。 」\par
泽同学所提及的这两种动的方式，会有区别 。\par
「这两者会致『『相关』的行度不同。测量它的就是C。这个算式是什么意思觉明\par
白 ?」\par
「不，完全不明白……。 」\par
「[缺字]A($\theta$)[缺字]的意思是，用设置$\theta$测量物理量 A 时的期望。。所以，[缺字]A($\theta$)B($\phi$)[缺字]是物理\par
量乘积的期望。\par
2\par
。」\par
「那么，C 就是将物理量乘积的期望。加加减减的东西……?」\par
「这么理解就好。是在。改改变各种设定时的期望。。 」\par
4 定域 实在论的情形\par
「首，是……定域实在论。 ，虑C 可能取。的范围。 『确定』 的意思， 是只要知晓所\par
有隐变量，就能确实地言测其。。也就是说，如果设隐变量为$\lambda$，物理量就能写成其函\par
数。 」\par
A = a($\lambda$)，B = b($\lambda$)\par
「将这个与『定域性』 ，即地点A 改变设定不会影响 B，反之亦然，一并反映来写，\par
就能写成这样。 」\par
泽同学擦掉才才写的式子，重新写。\par
A = a($\theta$, $\lambda$)，B = b($\phi$, $\lambda$)\par
「使用这些，其乘积的期望。就用$\lambda$的概率分布表示成这样。 」\par
[缺字]A($\theta$)B($\phi$)[缺字]= $\sum$ P($\lambda$)a($\theta$, $\lambda$)b($\phi$, $\lambda$)\par
$\lambda$\par
「在此基上上，我们来，虑C。首，，为了让式子不那么乱，这样简记。 」\par
a = a($\theta$, $\lambda$)，a$\prime$ = a($\theta$$\prime$, $\lambda$)，b = b($\phi$, $\lambda$)，b$\prime$ = b$\prime$($\phi$, $\lambda$)\par
2\par
「这里，如果 A 与 B 一『就为+1，相反则为-1，所以称为相关。 」\par

% --- PDF page 12 ---
11\par
「然后，且且，求ab + a$\prime$b - ab$\prime$ + a$\prime$b$\prime$的取。范围吧」\par
|ab + a$\prime$b - ab$\prime$ + a$\prime$b$\prime$| = |(a + a$\prime$)b - (a - a$\prime$)b$\prime$|\par
$\leq$ |(a + a$\prime$)b| + |(a - a$\prime$)b$\prime$|\par
= |(a + a$\prime$)||b| + |(a - a$\prime$)||b$\prime$|\par
$\leq$ |(a + a$\prime$)| + |(a - a$\prime$)|\par
$\leq$ 2\par
「通过这样的计算，得到这样的式子。 」\par
-2 $\leq$ ab + a$\prime$b - ab$\prime$ + a$\prime$b$\prime$ $\leq$ 2\par
「对这个式子乘以P($\lambda$)再求和」\par
「这里就只是单计计算了，也没什么好说明的。简单。对吧?制作人。 」\par
她无聊似的问道。我觉得完全不是这样。\par
「便一一提，这个不等式取发现者 Clauser、Horne、Shimony、Holt 的首字母，称\par
为 CHSH 不等式。 」\par
5 量子论的情形\par
泽同学笃笃笃地走近我坐着的座位。 大概是站了这么久快撑不住了， 脚步看起来\par
有点不稳。薄款长大衣的下摆飘飘然地舞的着过来。\par
正看入了迷，下一秒，她伸出的双手夹住了我的脸颊。\par
，得看向她的脸， 同时， 好凉。 好闻的气味。 是香水 ?像桃子， 又像香皂……正想\par
着这些，\par
泽同学这样问道。\par
「制作人，觉对量子论了解多少?希望觉至少知道狄拉克符号。 」\par
「……对不起。不知道。 」\par
声音有些发尖，但我老实回答。突然被一个面容姣好的女性在这么近的距离提问，\par
无论内容是什么都会心跳加速。\par
「这一章如果不懂\par
3\par
，大『起起就好。只要知道最后的结果，就能明白故事脉络。 」\par
我轻轻点绪表示了解后，她说\par
3\par
「对懂的人，我把练习题放在脚注里了，的手做做看。制作人的话肯定能行。 」\par
-2 $\leq$[缺字]A($\theta$)B($\phi$)[缺字]+ [缺字]A($\theta$$\prime$)B($\phi$)[缺字]-[缺字]A($\theta$)B($\phi$$\prime$)[缺字]+ [缺字]A($\theta$$\prime$)B($\phi$$\prime$)[缺字]$\leq$ 2\par

% --- PDF page 13 ---
12\par
「……开始了。。 」\par
一回到了白板前。\par
「在算符形式的量子论中，物理量用米米算符表示。像Â($\theta$)、B̂($\phi$)这样。首，，，\par
虑将$\theta$、$\phi$设为 0 时的Â(0)、B̂(0)吧。复习一下，实验流程是这样的。 」\par
她指向才才写的板书。\par
「使用地点A 和地点 B 的测量。(a,b)，把必定得到该。时的态的，这样定义。 」\par
|a, b$\rangle$\par
「例如，当(a, b) = (+1, -1)时，就写作|+, -$\rangle$。这被称作Â(0)和B̂(0)的『共同本征\par
的』 。\par
4\par
」\par
「从态的的定义出发，我们希望Â(0)|a, b$\rangle$ = a|a, b$\rangle$。这样想的话， 」\par
Â(0) = $\sum$(|+, b$\rangle$$\langle$+, b| -|-, b$\rangle$$\langle$-, b|)\par
b\par
「……这样写就很自然。 」\par
「将此广广到Â($\theta$)。改变了$\theta$的测量，测量结果就会变化。 」\par
「例如，可以，虑像这样将标准正交基转转。 」\par
4\par
「要构建这个，Â(0)、B̂(0)必须对易。请确认一下。 」\par
实验流程：\par
1. 在充分远离的地点 A、B 处各有 1 名。测者。\par
2. 在 A 点， 使用测量器的设置$\theta$或$\theta$$\prime$的其中之一测量物理量 A， 测量。记为\par
1 或-1。\par
3. 在 B 点也同样地，使用$\phi$或$\phi$$\prime$的其中之一测量物理量 B，得到测量。。\par
4. 两人各自记录自己所用的设定和测量。。\par
5. 多次重复 2～4 步骤。\par
6. 测量后，将两人的数据对照计算 C。\par
其中\par
C =[缺字]A($\theta$)B($\phi$)[缺字]+ [缺字]A($\theta$$\prime$)B($\phi$)[缺字]-[缺字]A($\theta$)B($\phi$$\prime$)[缺字]+ [缺字]A($\theta$$\prime$)B($\phi$$\prime$)[缺字]\par

% --- PDF page 14 ---
13\par
|+$\theta$, b$\rangle$ = cos $\theta$\par
2 |+, b$\rangle$ + sin $\theta$\par
2 |-, b$\rangle$\par
|-$\theta$, b$\rangle$ = - sin $\theta$\par
2 |+, b$\rangle$ + cos $\theta$\par
2 |-, b$\rangle$\par
「在此之下，像Â(0)时那样定义的话」\par
Â($\theta$) = $\sum$(|+$\theta$, b$\rangle$$\langle$+$\theta$, b| -|-$\theta$, b$\rangle$$\langle$-$\theta$, b|)\par
b\par
= cos $\theta$ $\sum$(|+, b$\rangle$$\langle$+, b| -|-, b$\rangle$$\langle$-, b|) + sin $\theta$ $\sum$(|+, b$\rangle$$\langle$-, b| +|-, b$\rangle$$\langle$+, b|)\par
bb\par
「就能像这样广广。B̂($\phi$)也能用同样步骤致出\par
5\par
，变成这样。 」\par
B̂($\phi$) = cos $\phi$ $\sum$(|a, +$\rangle$$\langle$a, +| -|a, -$\rangle$$\langle$a, -|) + sin $\phi$ $\sum$(|a, +$\rangle$$\langle$a, -| +|a, -$\rangle$$\langle$a, +|)\par
aa\par
「然后，设设子子的态的是这样的。 」\par
|$\psi$$\rangle$ = 1\par
$\sqrt{}$2\par
|+, -$\rangle$ - 1\par
$\sqrt{}$2\par
|-, +$\rangle$\par
「计算此时的相关\par
6\par
，就成了这样。 」\par
[缺字]A($\theta$)B($\phi$)[缺字]= [缺字]$\psi$|Â($\theta$)B̂($\phi$)|$\psi$[缺字]\par
=- cos($\theta$ - $\phi$)\par
「用这个改写C，就是这样的吧。 」\par
说着，她终于写下了连我都能勉行理解的式子。\par
C = - cos($\theta$ - $\phi$) - cos($\theta$$\prime$ - $\phi$) + cos($\theta$ - $\phi$$\prime$) - cos($\theta$$\prime$ - $\phi$$\prime$)\par
「这个C，例如取」\par
$\theta$ = 3$\pi$\par
4 ，$\phi$ = 2$\pi$\par
4 ，$\theta$$\prime$ = $\pi$\par
4 ，$\phi$$\prime$ = 0\par
5\par
「练习题。请尝试致出。 」\par
6\par
「注意物理量是米米算符，计算会稍微轻松一点，。。请算算看。 」\par

% --- PDF page 15 ---
14\par
「按照这样的设定，就会得到C = -2$\sqrt{}$2。被破坏了呢，CHSH 不等式。 」\par
「实际上，这也是量子论所能取到的极，，被称为Tsirelson 界。」\par
6 总结\par
「结结来说，C 在定域实在论和量子论中分别，虑时，可能会取不同的。。如果物\par
理量是确定的， 换句话说， 『上帝不掷骰子』 的话，C 就应该落在$\pm$2 的范围内。 但如果\par
超出了这个范围，就意味着『上帝掷骰子』 ，即量子论才是正确的理论。这篇论文发表\par
后，法国的阿兰$\cdot$阿斯佩博士进行了实验验证，在 1982 年以题为《使用时变分析器对\par
贝尔不等式的实验测试》\par
7\par
的论文，展示了不等式确实被破坏。2022 年，他因此功绩获\par
得了诺贝尔物理学奖。 」\par
「物理量并不实在，真是……不可思议啊。 」\par
「是呢。 经典力学和电学学曾素素地设定 『物理量的实在』 ， 这在现实中却并不成立，\par
这是个不可思议的结果。 厘清定域实在论与量子论的本质区别， 并明确了后者必然性的\par
这个实验，是物理学史上最深刻的发现之一。 」\par
她说着，啪嗒一声合上笔记本，将笔放回白板。\par
「说真的，内容相当有趣。 」\par
满脑子私心打算找年纪小的偶像教自己学习的我，真感到羞耻。\par
「还没完。最后，还有一件事。 」\par
「什么?」\par
「这个实验，不不能持实实验装置的定域性就无法成立。也就是说，非定域的实在\par
论，用这个实验无法具定。希望制作人的想一想它的理由。 」\par
「这是，最后的练习题。」\par
7 参考文献\par
[1]\par
.(2004\par
).\par
:\par
(\par
).\par
.\par
[2]\par
.(2006\par
). EPR\par
. https://as2.c.u -\par
tokyo.ac.jp/lecture\_note/kstext04\_ohp.pdf(2024\par
10\par
8\par
)\par
7\par
译者注：Aspect A, Dalibard J, Roger G. Experimental test of Bell's inequalities using time-varying\par
analyzers[J]. Physical review letters, 1982, 49(25): 1804.\par
-2$\sqrt{}$2 $\leq$ C $\leq$ 2$\sqrt{}$2\par

% --- PDF page 16 ---
15\par
[3] Physics Lab. 2021.(2021\par
).\par
.\par
https://event.phys.s.u-tokyo.ac.jp/physlab2021/articles/1hk-9xt23/(2024\par
10\par
8\par
)\par
[4] Physics Lab. 2021.(2021\par
).\par
.\par
https://event.phys.s.u-tokyo.ac.jp/physlab2021/articles/g6xhsg3xs0h9/(2024\par
10\par
8\par
)\par

% --- PDF page 17 ---
\noindent\textit{[第 17 页没有可抽取文字层。]}\par

% --- PDF page 18 ---
17\par
需要的是纠缠\par
\textasciitilde{}用量子信息让\par
泽广验证定域实在论的破缺\textasciitilde{}\par
$\diamond$初星学园外围\par
广「美铃，早上好」\par
美铃「哎呀……\par
泽同学，早上好。今天要的哪里呢?」\par
广「今天练习休息，所以为了增行体力在散步。美铃呢?」\par
美铃「今天很适合散步，我也就着意走走。方一的话一起如何?」\par
广「好啊。走吧」\par
***\par
广「……好热……」\par
美铃「不要紧 ?」\par
广「要紧，才不是呢。不过这也是训练」\par
美铃「勉行可不行。。的那边的树荫下休息一下吧」\par
广「说得……是呢……」\par
广「呼……」\par
美铃「缓过来了 ?」\par
广「…。已经没事了。谢谢，美铃」\par
美铃「能帮上忙太好了。要不要的买点喝的?」\par
广「那个不用。制作人给了我宝矿力和冰袋」\par
美铃「哎呀……\par
泽同学的制作人结是准备周到呢」\par
广「制作人结是能让我做喜欢的事。我喜欢他这一点」\par
美铃「呵呵，制作人和\par
泽同学的相遇，也许是命中注定呢」\par
广「什么意思?」\par
美铃 「因为两位实在太般配了……虽然是很不可思议的关系， 但让人觉得是必然的\par
相遇」\par
广「确实，能遇见制作人真好。不过，不是必然。因为定域实在论被具定了」\par
美铃「定域……?」\par

% --- PDF page 19 ---
18\par
广「就是说， 为了让我和制作人相遇所需的变量在时间上的变化不是定定性的。谁\par
也没法法 100\%言测」\par
美铃「好不可思议的话题呢。这是哪个学问的主题?」\par
广「物理学。通过证明 『贝尔不等式的破缺』具定定域实在论的人，获得了诺贝尔\par
物理学奖」\par
美铃「来来是物理学的研究啊。真不愧是\par
泽同学」\par
广「……，贝尔不等式的破缺很简单。大概美铃也能理解，而且也能实证」\par
美铃「哎呀，是这样 ?天色还早，机会难得，能请觉详细教教我 ，\par
泽老师?」\par
广「当然。交给老师吧，美铃同学」\par
***\par
$\diamond$定域实在论中的关系式\par
广「首，，我教觉定域实在论下成立的不等式」\par
美铃「\par
泽老师，能请觉详细讲讲定域实在论 ?」\par
广「好的。定域性，是指做物理实验的场所彼此独立。实在性，是指实验得到的物\par
理量存在隐变量」\par
美铃「抱歉。不习惯的词太多，我不太懂。能用更贴近生活的说法 ?」\par
广「………对不起。所谓独立，是指自己的实验测量不影响别处的测量的情况。比\par
如，设设现在佑芽正在看千奈，如果『佑芽看千奈』这件事对『美铃看我』这件事没有\par
影响，那它们就彼此独立，可以说具有定域性」\par
美铃「来来如此……那么，如果彼此距离很近，定域性就会消失 ?」\par
广「好问题。确实彼此距离近的话，各自的测量方法似乎会影响其他测量，但定域\par
性需要的是『某一方。测的那一瞬间不给其他方带来影响』 。连光的传播都需要时间，\par
所以远近与定域性的有无没有太大关系」\par
美铃「……理解了。\par
泽老师」\par
广「呵呵，太好了。那么接下来是关于实在性。这对我们来说也是理所当然的事。\par
比如， 想知道我扔出的应3 秒后的位置。 我既可以实际的测3 秒后应的位置， 也可以通\par
过速度、重力、空气阻力计算出应的位置。像这样，存在一种在我们测量之前就能算出\par
结果的变量，这就是实在性」\par
美铃「确实觉得是理所当然的事。如果没有了实在性会怎样呢?」\par
广「。测之前，东西的态的是不确定的。所以，我和美铃在1 秒后互相换位置的可\par
能性，虽然概率极低，但也并非没有可能」\par

% --- PDF page 20 ---
19\par
美铃「真是不可思议的感觉呢……」\par
广「从算式上看就不觉得不可思议了。 那么回到开绪，教觉具有定域性和实在性时\par
成立的不等式」\par
美铃「拜托了」\par
广「虽然有点抽象，但内容很简单，觉放心。这个不等式表示两个物体态的的相关\par
性，所以，准备两个实验场 A、B，从它们的中央将想要。测的东西同时投向各自的地\par
点。将A、B 处。测到的物体态的量分别设为 a、b。然后。测这些量的方法也各准备两\par
种，这次 a 有$\theta$和$\theta$'，b 有$\phi$和$\phi$'这两种。测方法。设设基于实在性存在隐变量$\lambda$，为\par
简单设物理量a = $\pm$1, b = $\pm$1，则下式成立。\par
-2 $\leq$[缺字]A($\theta$)B($\phi$)[缺字]+ [缺字]A($\theta$$\prime$)B($\phi$)[缺字]-[缺字]A($\theta$)B($\phi$$\prime$)[缺字]+ [缺字]A($\theta$$\prime$)B($\phi$$\prime$)[缺字]$\leq$ 2 (1)\par
这叫 CHSH 不等式。是获得了诺贝尔物理学奖的贝尔不等式中的一种。非常简洁\par
的式子」\par
美铃「简……洁……简洁 ……?」\par
广「美铃……?不要紧吧?」\par
美铃「对不起……对我来说似乎相当难。这个式子有什么意义呢?」\par
广「…咕咕咕咕……就如写的那样啦。意思是，定域实在论中的任何东西，只要实\par
验条件齐备，。测结果都会满足这个式子」\par
美铃「能举个什么例子 ?」\par
广「是呢……比如，设设A 在地应，B 在月应，有台机器能把鲷鱼烧分成绪和尾从\par
两者间同时投出。 」\par
美铃「……?」\par
广「地应的。测者想吃鲷鱼烧的绪或尾，这分别就是。测方法$\theta$和$\theta$$\prime$。然后。测投\par
出的鲷鱼烧，想吃的一方定为 1，不想吃的定为-1，这就是物理量a。$\phi$、$\phi$$\prime$和b也是同\par
理。只要基于定域实在论，怎么样设定都可以」\par
美铃「好像懂了……又好像没懂……?」\par
广「足够了。到此为止都是前提，只要记住在我们的常识中有个不等式成立就好。\par
接下来教觉破掉这个式子的理论。加油。 」\par
美铃「明白了。\par
泽老师，拜托了」\par
***\par
$\diamond$贝尔不等式的破缺\par
广「破掉贝尔不等式的理论，那就是量子力学」\par

% --- PDF page 21 ---
20\par
美铃「量子力学是……什么呢?」\par
广「是物理学的领域，研究极极微小世界里无现的法则的学问。量子力学中，。测\par
态的$\psi$时乘积的期望。可以用算符 A,B 表示为$\langle$$\psi$|A($\theta$)B($\phi$)|$\psi$$\rangle$。 进一步， 设设测量。是\par
$\pm$1 进行计算，\par
$\langle$$\psi$|A($\theta$)B($\phi$)|$\psi$$\rangle$ = cos($\theta$ - $\phi$) (2)\par
就会变成这样，将才才不等式的相关。部分设为 C，使用量子力学替换之后，\par
C =[缺字]A($\theta$)B($\phi$)[缺字]+ [缺字]A($\theta$$\prime$)B($\phi$)[缺字]-[缺字]A($\theta$)B($\phi$$\prime$)[缺字]+ [缺字]A($\theta$$\prime$)B($\phi$$\prime$)[缺字]\par
=- cos($\theta$ - $\phi$) - cos($\theta$$\prime$ - $\phi$) + cos($\theta$ - $\phi$$\prime$) - cos($\theta$$\prime$ - $\phi$$\prime$) (3)\par
就变成这样。这里的证明稍微有点难，希望觉直接接受。然后此时，当。测方法为\par
$\theta$ = 0，$\theta$$\prime$ = $\pi$\par
2 ，$\phi$ = $\pi$\par
4 ，$\phi$$\prime$ = 3$\pi$\par
4\par
时，C = 2$\sqrt{}$2，大于 2，CHSH 不等式被破坏掉了」\par
美铃「哎呀呀……这可怎么法」\par
广「怎么了?」\par
美铃「我连自己为什么不懂都搞不懂了」\par
广「…………呵呵，绪疼了」\par
美铃「……为什么，式(1)和式(3)会产生差异呢」\par
广「在量子力学中， 。测之前物理量的选项叠加在一起不确定。 虽不确定但一对量\par
子会相互连接并互相影响，因此相关。会产生差异。这种连接叫做『纠缠』 」\par
美铃「纠缠……」\par
广「我和美铃，从现在没在。测我们的人看来，也处在纠缠态的」\par
美铃「有各种各样的可能性，打破我们常识的连接 ……呵呵，好棒啊。简直像是\par
为了\par
泽同学而存在的，不可思议的性质」\par
广「……是 ?」\par
美铃 「因为\par
泽同学的演唱会会给人难以言喻的不可思议的感情。 如果，虑\par
泽同\par
学和粉丝们也处在纠缠的，那我觉得这词太适合形容\par
泽同学的演唱会了」\par
广「呵呵，有意思的想法。……贝尔不等式的破缺，明白了 ?」\par
美铃「虽然模模糊糊，但明白了。……抱歉，贝尔不等式被破坏掉，是有什么意义\par
来着?」\par
广「贝尔不等式的破缺是对定域实在论的具定。也就是说，我们所在的世界，我们\par
以为理所当然的感觉行不通。但是，现在可以轻松地自己确认量子的行为」\par
美铃「轻松地， ?」\par

% --- PDF page 22 ---
21\par
广「这次确实简单。确认的同时还能验证贝尔不等式的破缺。交给我吧」\par
***\par
$\diamond$用量子计算机验证贝尔不等式的破缺\par
广「现在用量子计算机破掉贝尔不等式」\par
美铃「量子计算机是什么呢」\par
广「是利用量子力学进行计算的计算机。还在发展途中，有各种各样的实现方式，\par
基本上以操作处于量子的的东西，。测结果以 0 或 1 的二进制，出。网上就有能用的，\par
就用它来验证贝尔不等式的破缺。美铃，带手机了 ?」\par
美铃「是的，带着。」\par
广「太好了。那首，从 IBM Quantum 的网站创建名叫 IBMid 的账户」\par
美铃「……创建好了」\par
广 「接着登录IBM Quantum Platform，然后打开右下方的 IBM Quantum Composer\par
图图1】\par
。没有的话搜索一下就出来了」\par
美铃「打开了」\par
（图1）\par
广「在这里能制作操作量子的量子电路。中央的q[0] , q[1] 各自是量子的号号，c2\par
表示想用作计算结果的位数。 左侧面板里的方形图块表示量子的操作方法， 将它们拖放\par
到 中 央 的 线 上 就 能 制 作 量 子 电 路 。 然 后 ，式 (2) 做 出 的 量 子 的 乘 积 的 期 望 。\par
$\langle$$\psi$|A($\theta$)B($\phi$)|$\psi$$\rangle$在量子电路中可以像这个图\par
图图 2】\par
这样表示。 」\par

% --- PDF page 23 ---
22\par
（图2）\par
美铃「这个有什么意义呢?」\par
广「准备q[0]和q[1]两个量子，用 H 这个操作变成叠加的，用$\oplus$操作使其纠缠。用\par
RY 定定。测方法，最后测量」\par
美铃「虽然不太懂……但这满足贝尔不等式的条件了吧」\par
广「美铃，真不简单。这次把这个电路测量 1024 次来求期望。。试把设定设为\par
$\theta$ = 0，$\phi$ = $\pi$\par
4\par
按下右上角的 Setup and run， 从左侧栏\par
图图3】\par
里选择写着 ibm\_〇〇的。 选哪个都行，\par
但选 Total pending jobs （等待计算数）少的比好好。选好后按下右下角的Run on ibm\_\par
〇〇稍等一会儿，计算结果就会出现在 Composer Jobs\par
图图 4】\par
里」\par
（图3）\par

% --- PDF page 24 ---
23\par
（图4）\par
美铃「计算结果出来了，这个……该怎么处理呢\par
图图 5】\par
」\par
（图5）\par
广「这表示 1024 次测量中量子组合的分布。才才也说过，量子计算机的结果以 0\par
或 1 的组合，出，为配合 CHSH 不等式的条件（物理量为$\pm$1） ，将0 当作-1 处理，取\par
各自乘积的和之后除以测量次数求平均。 对所有。测方法的组合都这么做， 就能算出贝\par
尔不等式的相关。 C。结果是这种感觉\par
图图 6】\par
」\par
（图6）\par

% --- PDF page 25 ---
24\par
广「C 超过了 2，果然破坏掉了贝尔不等式」\par
美铃「虽然很难……但量子的不可思议性质，能在网上立刻验证呢」\par
广「…。习惯后就很简单。有实感了 ?」\par
美铃「虽不及\par
泽老师，但感觉稍微变得亲近了一点」\par
广「美铃很厉害呢。广老师给觉学分」\par
美铃「呵呵……谢谢。能窥见\par
泽同学眼中的世界很有趣。期待下次\par
泽同学的演\par
唱会。」\par
广 「交给我吧。 量子力学也好演唱会也好， 都是因为有人在看才成立。 所以演唱会，\par
我会唱出让我和大是心的不已的……那样的歌。 」\par
美铃「……那么我得法一场能让\par
泽同学鼓起干劲的演唱会才行呢」\par
广「呵呵……好期待。那，继续散步吧」\par
Fin\par
参考文献\par
[1]\par
(2004).\par
[2]\par
(2021).\par
.\par
[3]\par
(2024).note\par
. https://note.com/quantumuniverse/n/n3e9f1056fcca (\par
:\par
2024/10/22)\par

% --- PDF page 26 ---
25\par
专属 P 道具：恶魔核心\par
著： シメギウム\par
亚纱里老师：\par
好了好了\par
说到\par
泽同学的课题，果然还是肌力不足。\par
为了克服这一点，现在这堂课上也在让她努力做肌肉训练。\par
广：\par
手臂要断了……\par
亚纱里老师：\par
那么现在，她正用螺丝刀抵住的金属半应，一般被称作什么呢?\par
广：\par
帮帮我，制作人……\par
P：\par
[麦克风图标]能量应（声乐上升+40 选择技能卡后获得）\par
[鞋图标]恶魔核心（舞蹈上升+80 选择技能卡后获得）\par
→[人物图标]赶紧把金属半应扔出的（往牌组中加入 1 张「困意」卡）\par
（好险………!）\par
亚纱里老师：\par
没错，持护偶像远离险…也是制作人的工作。\par
虽然是复制品，但判断得很好。♪\par
广：\par
谢谢，制作人。\par
托觉的福，我的右臂，还是直的，。。\par

% --- PDF page 27 ---
26\par
P：\par
（都不懂别人的心情……!）……那就好。\par
亚纱里老师：\par
来来如此，真是良好的信任关系呢♪\par
P：\par
…………。\par
话说回来，恶魔核心实际上是什么东西呢?\par
亚纱里老师：\par
那么，今天就从核物理学的基上来学习一下恶魔核心吧。\par
广：\par
从这里开始，由我来教觉，。。\par
P：\par
（这里是文科的制作人科啊……）\par
……请多指教。\par
广：\par
制作人，觉知道什么是放射线 ?\par
P：\par
就是 X 光之类用的那种看不见的光吧，是从铀或钚里发出的吧?\par
广：\par
难道，制作人，觉以为 X 光机里有铀 ?\par
P：\par
不是 ?\par
广：\par
……前途堪忧啊。\par

% --- PDF page 28 ---
27\par
放射线的定义有很多，不过「在空间或物质中传播能量的子子束或电学动」 这个大\par
概是最简洁的，吧。\par
亚纱里老师：\par
《来子能基本法》\par
1\par
第一章结则第三条将 「放射线」 定义为： 「电学动或子子束中， 直\par
接或间接具有电离空气能力，且由政令规定的射线。 」具体来说，通常多指《电离放射\par
线障害防止规则》\par
2\par
第一章结则第二条所规定的「$\alpha$射线、重质子束及质子束」 「$\beta$射线\par
及电子束」 「中子束」 「$\gamma$射线及 X 射线」 。\par
P：\par
（虽然现在说有点晚，\par
泽同学倒是算了，怎么连亚纱里老师都这么清楚……?）\par
广：\par
恶魔核心所用的放射性物质， 主要是$^{2}$$^{3}$$^{9}$Pu 和通过中子俘获生成的不计物$^{2}$$^{4}$$^{0}$Pu。其\par
中，$^{2}$$^{3}$$^{9}$Pu 几乎都是 $\alpha$ 衰变。半衰期是 24110 年。是意外地稳定的同位素啊。\par
另一方面，X 光检查中使用的是动长好短的光， 即X 射线。 在真空管中， 对作为靶\par
子的阳极放电， 由热电子使该物质积蓄能量。 然后， 在的放能量时， 一以X 射线的形式\par
放出来，。。更简单地说，它和荧光灯几乎来理相同。\par
P：\par
$\alpha$衰变，就是亚纱里老师说的的放出$\alpha$射线的反应 ?\par
广：\par
对。$\alpha$射线是由 2 个质子和 2 个中子构成的放射线。这和$^{4}$He 的来子核相同，。。\par
空气中的 $\alpha$ 射线射程 R(cm)可用能量 E(MeV)表示为\par
R = 0.318E3/2\par
。。$^{2}$$^{3}$$^{9}$Pu 放出的 $\alpha$ 射线最大能量为 5.16 MeV，所以大约是 3.7 cm。\par
亚纱里老师：\par
MeV 是能量单位，是兆电子伏特的缩写。1 MeV 约相当于 1.6$\times$10$^{-}$$^{1}$$^{3}$ J。\par
1\par
译者注：指的是日本的《\par
》\par
2\par
译者注：指日本的《\par
》\par

% --- PDF page 29 ---
28\par
P：\par
在空气中也飞不到 4 cm 的话，那岂不是不太险…?\par
广：\par
…。作为$\alpha$射线源，倒是不怎么险…。吃了或摸了另当别论，但即一是我，这点程\par
度也不会死的……应该。问题出在另一个反应上。那就是，连锁反应。\par
P：\par
那是会放出$\alpha$射线以外的放射线 ?\par
广：\par
对。不过，连锁反应不是反应本身的名字。它是核裂变产生的中子，进一步引发其\par
他核裂变的连续性反应。\par
P：\par
核裂变我也起说过。就是核电站之类的使用的反应吧?\par
广：\par
…。核裂变是铀或钚这类超铀元素的重来子核，被中子撞击后发生的， 分裂成两个\par
裂变碎片来子的反应。会分裂成什么来子是着机的，不过产率好高的来子， 在来子序数\par
上大的一方有$^{1}$$^{3}$$^{3}$Cs、$^{1}$$^{3}$$^{7}$Cs、$^{1}$$^{3}$$^{5}$I、$^{1}$$^{3}$$^{3}$Xe，小的一方则有$^{9}$$^{3}$Zr、$^{1}$$^{0}$$^{7}$Pd、$^{9}$$^{0}$Sr，。。\par
这个核裂变反应，始于来子核吸收中子，那时会放出非常大的能量，。。\par
例如，$^{2}$$^{3}$$^{5}$U 分裂成$^{9}$$^{2}$Kr 和$^{1}$$^{4}$$^{1}$Ba 的话，可表示为\par
$^{2}$$^{3}$$^{5}$U + n → $^{2}$$^{3}$$^{6}$U → $^{9}$$^{2}$Kr + $^{1}$$^{4}$$^{1}$Ba + 3n + Q\par
。。其中 n 是中子。\par
这时，我们来看一下反应前后的质量。\par
各质量以\par
$^{2}$$^{3}$$^{5}$U: 235.04393 u、$^{9}$$^{2}$Kr: 91.92616 u\textbackslash{}$^{1}$$^{4}$$^{1}$Ba: 140.91441 u、n: 1.00867 u、\par
p: 1.00728 u、e: 0.0005485 u 来计算。。\par
亚纱里老师：\par
u 是用来表示来子或来子核质量的来子质量单位。以$^{1}$$^{2}$C 的质量为 12 u 来定义，\par
1 u = 1.66054$\times$10$^{-}$$^{2}$$^{7}$ kg，由于质量与能量等价，也可表示为 1 u = 931.478 MeV。\par

% --- PDF page 30 ---
29\par
广：\par
那么，我们来算算看吧。因为只是简单加法，过程就省略了。左边质量是\par
236.05257 u。 右边质量是235.86658 u， 。。 其差。是0.18599 u。 也就是大约173 MeV。\par
这是每个来子的能量，所以如果有 1 kg 的$^{2}$$^{3}$$^{5}$U，就是 7.1$\times$10$^{1}$$^{3}$ J 的能量。大约是 2 万\par
吨 TNT 炸药的能量。\par
P：\par
确实是，人的能量，但那是在$^{2}$$^{3}$$^{5}$U 全部反应的情况下吧?1 kg 就有 4 mol 以上的来\par
子，所有来子都会分裂 ?\par
广：\par
真不愧是制作人。敏锐。\par
才才的计算有两个没，虑到的因素，反应速度和临界量。 本质上是相同的话题，不\par
过我们分开来讲，。。\par
就像才才说的，核裂变始于吸收中子。但是，下一次反应当然也需要中子。也就是\par
说，最初反应产生的中子，会成为下一次反应的起爆剂。这形成连锁，因此的放出非常\par
巨大的能量。\par
不过， 并非所有中子都能反应。 有些会从铀或钚的块体中漏出的，或者能量太弱无\par
法反应。正因此，反应才会变快或变慢。\par
那么，我们来想想什么情况下反应会实续。\par
设设发生一次核裂变反应，此时产生了$\eta$个中子。因为反应会吸收 1 个中子，也可\par
以说核裂变前后中子数目增加了$\eta$ - 1个。\par
要维实反应， 每次核裂变需要引起一次以上的核裂变， 因此只要产生的中子平均被\par
吸收1/($\eta$ - 1)次，反应就会实续。\par
设中子首次被吸收前的距离 （平均自由程） 为$\lambda$f， 移的距离R 期间会发生R/($\lambda$f)次\par
核裂变。\par
另外，利用来子核的数密度n和微。截面$\sigma$f，可表示为\par
$\lambda$f = 1\par
n$\sigma$f\par
。。\par
P：\par
数密度就是高中化学学的来子数量的密度 ?比如体心立方晶格里含有 2 个那种。\par

% --- PDF page 31 ---
30\par
亚纱里老师：\par
基本答对了。 不过， 制作人同学觉说的体心立方晶格那种规则的晶体结构自不必说，\par
它也能用于气体、液体之类不规则的结构。♪\par
说到高中化学，数密度可以用来子核的质量 m、密度$\rho$、克来子量 M、阿伏伽德罗\par
常数Na表示为\par
n = $\rho$Na\par
M\par
便一一提，才才\par
泽同学举出的微。截面，并不单计表示面积。在核物理中，子子\par
与物质发生反应的概率会用截面来表示。 常用的单位有靶恩 （barn， 简称靶），1 靶等于\par
10$^{-}$$^{2}$$^{4}$ cm$^{2}$。明明是概率却有单位，很不可思议吧♪\par
广：\par
好了，这里且且回到来话题，讲讲临界。在反应堆中，为了让连锁反应实续，同时\par
防止失控的放巨大能量而爆炸， 会进行控制使反应始终处于平衡态的。 这种平衡态的叫\par
临界， 反应速度比这小时连锁反应无法实续， 太大则会失控。恶魔核心就是测量这个临\par
界的实验。便一一提，实际反应堆中会用控制棒，由镉、硼、铪这类热中子吸收截面大\par
的材料制成，来适当控制，。。\par
亚纱里老师：\par
这里讲个小知识。据说利用硼的热中子吸收截面大这一特点，有种放射线治疗法，\par
将硼投予肿瘤，再照射反应堆出来的中子来进行治疗。这种被称为 硼中子俘获治疗\par
（BNCT）的疗法，具体是$^{1}$$^{0}$B 吸收中子后，产生 $\alpha$ 射线和$^{7}$Li，两者会杀死癌细胞，就\par
是这么回事。♪\par
正如\par
泽同学所说，$\alpha$射线射程短，$^{7}$Li 更短，所以不容易影响其他正常细胞。\par
用反应堆治病，真是太行大了……虽说想这么说， 但实际上，在医院运行反应堆风\par
…高、 成本大， 而且处理反应堆除了需要选任通常放射线治疗装置所必需的第一种放射\par
线处理主任者外， 还必须选任核燃料处理主任者和反应堆处理主任者\par
3\par
， 人事门槛很高，\par
因此现在主流是利用回转加速器这类加速器进行的 BNCT。\par
3\par
译者注：这三者都是日本\par
依据\par
等设立的国是资格。\par

% --- PDF page 32 ---
31\par
广：\par
在临界最小距离Rcr的物体内，要达成临界至少需要发生1/($\eta$ - 1)次反应，所以可\par
表示为\par
Rcr\par
$\lambda$f\par
= 1\par
($\eta$ - 1)\par
Rcr = 1\par
n$\sigma$f($\eta$ - 1)\par
。。代入老师说的数密度公式，\par
Rcr = M\par
$\rho$Na$\sigma$f($\eta$ - 1)\par
不过，临界通常不用距离而用质量来评估，所以临界质量Mcr为\par
Mcr = 4$\pi$Rcr\par
3$\rho$\par
3\par
Mcr = 4$\pi$\par
3 ( M\par
Na\par
)\par
3 1\par
$\sigma$f3($\eta$ - 1)3$\rho$2\par
。。\par
，虑$^{2}$$^{3}$$^{9}$Pu 的情况，用 M: 239 g，$\rho$: 19.8 g/cm$^{3}$，Na: 6$\times$10$^{2}$$^{3}$，$\eta$: 3.0， $\sigma$f: 2.2 barn\par
来计算，Rcr大约是 4.6 cm，Mcr大约是 8.2 kg。因为只是代入才才的式子，计算过程就\par
省略了，。。\par
P：\par
也就是说， 如果是真是伙，\par
泽同学就是撑着一个 4.1 kg 的半应了。 果然应该是复\par
制品。\par
广：\par
不是的。 常看到的恶魔核心实验中那个大应体， 是中子反射体铍做的半应。 实际上，\par
在底座部分有一块 6.2 kg 的钚块，再拿铍的盖子靠近，试图让它达到临界的实验。\par
就像才才计算的，6.2 kg 通常是不会临界的， 但用铍盖子不遗漏地反射逃往外部的、\par
几乎全部中子使之反应，这就是个测试要靠近到什么程度才会临界的实验。\par
作为核心的钚直径大约 9 cm， 目测那个铍半应是其两倍左右， 就按18 cm 算的话，\par
重量大约是 2.5 kg。\par

% --- PDF page 33 ---
32\par
4\par
P：\par
……就算是\par
泽同学，似乎也勉行拿得住呢。\par
广：\par
制作人。要面对现实。觉担当的偶像，单手拿不的 2.5 kg，也撑不住的。\par
P：\par
……。\par
广：\par
反应速度的话题还没讲呢。从来子核出来的中子平均寿命大约 10$^{-}$$^{5}$秒。也就是说，\par
达到临界后 1 秒，连锁反应的代数会广进 10 万次。即使中子倍增率取 1.0001，简单计\par
算\par
1.0001105\par
$\approx$ 22000\par
就会超过 2 万倍。\par
来稍微严谨一点计算。。设中子移的平均自由程距离的时间为tg，连锁反应开始\par
t 秒后的中子数为N(t)，每一代的中子增加量记作 x，则\par
4\par
译者注：『敬传奇改锥人路易斯$\cdot$斯洛廷\par

% --- PDF page 34 ---
33\par
dN\par
dt $\approx$ Nx\par
tg\par
N(t) = N0 exp (xt\par
tg\par
)\par
由此可知N(t)是世代数ng = t/tg的函数。也就是说\par
N(t) = N0 exp(x $\cdot$ ng)\par
这里的tg是中子以速度v移的平均自由程的时间，所以\par
tg = $\lambda$f\par
v = 1\par
vn$\sigma$f\par
使用 E 代表中子的的能，\par
v = $\sqrt{}$2E\par
mn\par
tg = 1\par
n$\sigma$f\par
$\sqrt{}$mn\par
2E\par
用才才求临界半径时的。，$\lambda$f为 9.2 cm，v为 2.0$\times$10$^{1}$$^{0}$ cm/s，tg为 4.7$\times$10$^{-}$$^{9}$秒。\par
由N(t)的式子，N可视为公比为ex的等比级数之和，因此\par
N = ex(ex - 1)\par
ex - 1\par
8.2 kg 的$^{2}$$^{3}$$^{9}$Pu 中有 2.06$\times$10$^{2}$$^{6}$个来子核， 不x 为 0.0001， 则约50 万代反应完成。\par
由才才求出的tg可知，约 0.002 秒反应就会结束。\par
P：\par
和最初的简单计算，结果完全不同啊?\par
广：\par
…。繁杂的计算是仅用核裂变同时放出的瞬发中子来算的。而中子的平均寿命，是\par
包含了裂变后十几秒放出的缓发中子的平均寿命。 缓发中子是裂变碎片生成后， 经过一\par
段时间才放出的。$^{2}$$^{3}$$^{9}$Pu 的情况下，放出的缓发中子据大约 0.2\%。\par

% --- PDF page 35 ---
34\par
P：\par
0.2\%的话几乎没什么影响吧?\par
广：\par
但是，寿命 10$^{-}$$^{5}$秒的据 99.8\%，但设如有 0.2\%是 15 秒，平均就是 0.03 秒，。。\par
多亏了这些迟来的缓发中子，反应堆才能控制。\par
作为衡量有效倍增因子 （扣除为了反应而被吸收的中子） 偏离1 多少的指标， 反应\par
度经常被使用。用有效倍增因子减的 1 再除以有效倍增因子得到的。，如果这个。为\par
正，则因为有效倍增因子大于等于 1 所以反应进行。不为负。，反应就会中断掉。\par
作为反应度大小的参，， 常用缓发中子比例$\beta$作为基准。 有效倍增因子超过1 + $\beta$的\par
态的称为瞬发临界， 因为即使没有缓发中子反应也会加速地进行下的， 所以在反应堆中\par
受控使之不会超过，。。\par
亚纱里老师：\par
在反应堆安全管理中，达到瞬发临界的反应度$\beta$常以单位 1 \$或它的百分之一\par
1¢来表示。1 \$相当于 0.0065，但实际运用中控制得非常精细，所以大多使用¢。\par
用金钱单位来衡量核裂变险…度，真是不可思议的世界呢。\par
广：\par
众所周知的那个恶魔核心实验中，反应度似乎超过了 10¢。\par
按照才才的公式，有效倍增因子超过10¢的话， 来本需要50 万代才能完成的反应\par
就会在 8 万代内结束， 也就是说反应速度达到了6 倍以上。 这么一想， 就能理解那次实\par
验有多么险…了。\par
亚纱里老师：\par
据说实验中， 最靠前的物理学是路易斯$\cdot$ 斯洛廷\par
5\par
受到了 21 Sv 的辐射。 全身照射10\par
Sv 据说是『死量。像医生和反应堆作业员这样从事与放射线相关工作的人，其有效剂\par
量，。参照国际放射防护委员会 （ICRP）2007 年建议， 为100 mSv/5 年或 50 mSv/年，\par
而我们这样的普通公众则彻底安全管理在 1 mSv/年。这是根据联合国来子辐射效应科\par
学委员会（UNSCEAR）1988 年报告定定的，据说18 岁不受到 100 mSv 辐射，言期寿\par
命损失约为半年。虽并非单计与剂量成正比， 但这样一想，就很清楚那真是十分，人的\par
剂量。\par
5\par
译者注：也就是 30 页那张照片里的传奇改锥人\par

% --- PDF page 36 ---
35\par
P：\par
嘛，无论如何，我是很清楚\par
泽同学做不了恶魔核心实验了。快转到训练吧。应该\par
没空把上课时间花在核物理上。\par
广：\par
还想再聊一会儿……\par
P：\par
亚纱里老师，谢谢觉。\par
亚纱里老师：\par
辛苦了。上课也请加油。♪\par
获得技能卡「困意」\par
作者附记\par
感谢各位读到这里。 因为其他各位似乎偏重理论， 我本想多加入一些实务内容和法\par
律方面的话题，结果扯太远，只好含泪割爱了。作为注的的亚纱里老师实在太好用，变\par
得比广还要详细了……\par
平时我在电影评论社团 「\par
」 负责日本未上映电影的介绍书 《Incognita》\par
系列，也做其他电影书。如果赶得上，我想冬 CM 也会出品，方一的话请来玩。\par
X: @Nina\_Shimegi\par

% --- PDF page 37 ---
\noindent\textit{[第 37 页没有可抽取文字层。]}\par

% --- PDF page 38 ---
37\par
选择者\par
余命\par
「再次问候， 各位!感谢大是今天齐聚于此， 衷心表示感谢。 我是主实人仓本千奈。!」\par
培养出无数顶尖偶像的国内最大偶像培养学校「初星学园」 。\par
我曾经是这所学园制作人科的学生。\par
并且，负责制作一位名叫\par
泽广的偶像科学生。\par
现在所在之处，是学园里尤其充满回忆的教室。\par
当时这里，是作为制作活的的据点来用的。\par
大约一个月前，突然收到了仓本同学的邀请函。\par
上面写着，时隔十几年将举法同窗会，请务必前来。\par
似乎并不是多年来只漏掉了我没接到同窗会的联络。\par
如此这般，今天我一边咀嚼着怀念，一边踏入了初星学园的大门。\par
回想起来，那是新鲜、痛苦又事事不由己的日子。\par
泽广。\par
我还是制作人科学生时的担当偶像。\par
她喜欢辛苦的事、痛苦的事和不便利的事，为了摆脱简单而无聊的日常，入学了初\par
星学园。\par
和仓本同学是极好的朋友。\par
教室里，聚集了\par
泽同学当时的朋友们。\par
然而\par
今天这场聚会的参加者名单中，没有\par
泽广的名字。\par

% --- PDF page 39 ---
38\par
「\par
泽同学，明天的物理，试，我只有不安啊～～～!」\par
「千奈，好像很痛苦。 」\par
「虽然学了，可是物理太难了啊。 」\par
「……是 ?」\par
「就是啊!」\par
「对了，要不要互相出题?」\par
「好主意。 」\par
「那么拜托了。!放马过来吧!」\par
「提问。 构成这个世界物质的最小单位是什么?将物质不断分割变小， 最终会剩下什\par
么?」\par
送担当的偶像登上舞台的瞬间。\par
这个瞬间结是令人紧张。\par
「我出发了。 」\par
将那个背影烙印在眼中。\par
我按照与\par
泽同学的约定，匆忙离开舞台后方。\par
会场内的地图自然全都在脑子里。\par
沿最短路线走过工作人员用的后台通道。\par
能起到「初」开始播放。\par
今天，举行了\par
泽同学学年的毕业典礼。\par
我见证了从一年级起一担当的偶像的毕业典礼。\par
但不能只顾着感的。\par
在这初星学园，比起「毕业典礼」 ，之后举行的「毕业公演」才是更盛大的活的。\par

% --- PDF page 40 ---
39\par
私立初星学园高中部偶像科，集学园生活之大成的毕业公演。\par
为了抢，一睹未来顶尖偶像的风采，众多。众从学园内外涌来。\par
最大的看点果然还是「初」吧。\par
在返场时表演是每年的惯例。\par
毕业的所有偶像登上舞台，歌唱，舞蹈，为漫长又短且的学园生活拉下帷幕。\par
而且今年，似乎将有过的最多人数的偶像登上舞台。\par
或许是因为与往年不同，成绩靠后的偶像中退学的人也少了的缘故。\par
至今为止，\par
泽同学的演唱会我一直都在舞台后方守候。\par
但唯有最后，要从。众席见证她的身影。\par
拐过下一个转角，应该就能看到门了。\par
穿过那扇门，。众席就在眼前。\par
「提问。 构成这个世界物质的最小单位是什么?将物质不断分割变小， 最终会剩下什\par
么?」\par
「这是基上中的基上呢。物质的最小单位，是多边形。!」\par
「完全正确。太简单了 ?」\par
「是的。接下来我出题。。多边形的集合被称作什么?」\par
「网格。 」\par
「正是如此。!」\par
「More Is Different。网格物理学有着与多边形物理学不同的趣味。 」\par
「提问。为什么多边形是最小单位，觉是怎么想的?」\par
「……～!我投降了。答案是什么啊?」\par
「答案，还没有。人类还不知道。 」\par
「好难难人啊!」\par
「呵呵。虽然没有答案，但可以空想答案。物理学就是这样发展起来的。出色的空\par
想，非常重要。 」\par
「来来如此啊……」\par
「一起空想吧，很快乐的。 」\par

% --- PDF page 41 ---
40\par
「因为是如如此创了了世界，这个怎么样啊?」\par
「千奈，厉害。就是这个态的。那觉觉得，如为什么把世界创了成了多边形为最小\par
单位的样子呢?」\par
「\par
泽同学。现在轮到我出题了。。\par
泽同学是怎么想的啊?」\par
「……因为最容易创了吧。因为最容易计算吧。 」\par
「最容易计算?」\par
「对。模假设说。我最喜欢的空想。 」\par
今天这场聚会的参加者名单中，没有\par
泽广的名字。\par
泽同学，离开了我身边，的了很远的地方。\par
是一种既寂寞、又骄傲，摇摆不定的心情。\par
即使从初星学园毕业，\par
泽同学和我也继续持实着偶像与制作人的关系。\par
有学园经营的经纪公司，偶像科毕业的学生几乎都隶属于那里。\par
理所当然，事事不由己的日子仍在继续。\par
契机是，\par
泽同学的履历悄悄在网上成了话题。\par
「年仅14 岁一毕业于海外大学的天才美少女偶像」\par
来自猜谜节目的出演邀请开始频频飞来。\par
曾约定尽量不拿\par
泽同学的履历做文章的我，一直广掉这类工作，可我太天真了。\par
某档验证类节目偷偷策划的、在休息室验证\par
泽同学知识的企划。\par
完全疏忽了。\par
那档节目来本收视率就高， 反响极大， 一旦启的的潮流， 凭我一己之力已无法阻止。\par
猜谜节目的出演邀请。\par
科学活的的出演邀请。\par
什么什么学会的演讲邀请。\par

% --- PDF page 42 ---
41\par
学术出版社的执笔邀请。\par
世间已然不再追求「天才」之外的卖点。\par
我早就知道结有一天会变成这样。\par
即一如此。\par
从最底层一步一步，在不自由中挣扎着，以顶端为目标。\par
正当享受地走在漫漫长路上时， 却突然有种被行制抄了近路抛到很远的前方的感觉。\par
当然，偶像的工作没有稳定可言。\par
事事不由己的日子并没有结束。\par
但确实，变得乏味了。\par
泽同学，找到了下一个「趣味」 。\par
其实，同窗会已经结束了。\par
同窗会后，与\par
泽同学关系好的人聚在一起，转移到了教室。\par
这个聚会有着与同窗会不同的目的。\par
「对。模假设说。我最喜欢的空想。 」\par
「模假设说到底，是什么啊?」\par
「就是说，这个世界的一切，说不定是某个人用超高性能计算机模假出来的世界。\par
我们所有人，也许不过是数据上的存在。我们的肉体，心灵，一切都可能只是计算上的\par
存在。就是这样的设说。 」\par
「好可的的故事啊。 」\par
「目前没有证据。但是，这样的空想，非常快乐。 」\par
「就连我现在思，的事，也被计算在内了 ?那结让人觉得活得不是滋味啊。 」\par
「我们的记忆，也可能只是计算上的存在。我们也许，住在一个仅仅5 分钟前才被\par
制了出来的世界里。在植入了 5 分钟之前记忆的态的下。 」\par
「太可的了啊!」\par
「我们也许，只是住在经过无数次模假的世界中，其中的区区一个里。相似却又在\par
何处有所不同的世界，或许存在无数个。 」\par
「说不定也有制作人，生着着我的世界 ?」\par
「一定有的。 」\par
「这完全算不上安啊啊——!!」\par

% --- PDF page 43 ---
42\par
「那么提问。如果这个设说正确，觉觉得我们为什么会被模假这么多次?」\par
「一定是，为了看乎乎乎的我而乐在其中啊!」\par
「呵呵。意外可能猜中了。」\par
「才不是笑得出来的事啊!把问题还给觉。。那样的模假，真的可能 ?」\par
「可能。其实在理论上，这个空想已被证明了。我证明了。 」\par
「诶诶诶诶诶诶诶诶～～～～～～～!?」\par
这个聚会有着与同窗会不同的目的。\par
大是，都紧盯着教室前方的屏幕。\par
泽同学，的了很远的地方。\par
准确地说，接下来要的很远的地方。\par
物理上的。\par
「离发射时刻，只剩不到30 分钟了。!终于要来了啊!」\par
屏幕上映出的火箭。\par
她就在那里面。\par
宙航员偶像$\cdot$\par
泽广就在里面。\par
「诶诶诶诶诶诶诶诶～～～～～～～!?」\par
「其实，证明并没有那么难。只是理解概要的话，千奈也能做到。 」\par
「真的 ?」\par
「首，，我们全员都是计算者。是网络，是世界的。测者，是模假的执行者，是从\par
无数模假中选出唯一现实的选择者。 」\par
「来来如此?」\par
「某个计算者，可以计算另一个计算者的计算。计算者存在便序。计算者可以自我\par
复制。计算者会被统合。 」\par
「来来如此???」\par

% --- PDF page 44 ---
43\par
「最终，想要进行世界的模假，只需唯一的计算者进行计算就行了。 」\par
「来来如此??????」\par
「抱歉，不容易懂 ?」\par
「重要的是，这个空想的实验非常困难。 」\par
「是这样啊」\par
「提问。究竟做怎样的实验，才能验证模假设说呢?」\par
「完全!不知道啊!!」\par
「抱歉。的确，验证非常困难。但用某种方法，可以确认模假设说似乎正确。 」\par
「怎么做呢?」\par
「引发这个世界的，漏。。 」\par
「那个」是一瞬间的事。\par
注意到的，或许只有我。\par
在压轴毕业公演的「初」 ，进行到一半的时候。\par
我「选择」了「那个」 。\par
我被「选择」了「那个」 。\par
好沉重。\par
世界的时间的行进停止了，\par
感觉像是如此。\par
或许只是我的错觉。\par
我对这一现象有所绪…。\par
在模假设说中，以「处理延迟」为人所知的现象。\par
这个世界漏。的一种。\par

% --- PDF page 45 ---
44\par
无数次唱过，跳过，反复练习过的曲子。\par
嘴自己的着。\par
身体自己的着。\par
脑海中，空想转转飞舞。\par
漏。发生的条件有 3 个。\par
条件之 1：某个计算者 A 的模假负荷超过极，。\par
现在这个瞬间，有过的最多人数的偶像在舞台上举行演唱会。\par
构成偶像这一物质的多边形，其所有细微的作都在被计算。\par
模假负荷超过极，也不奇怪。\par
条件之 2：计算者 A 的计算对象中，包含另一个计算者 B。\par
当计算者 A 的模假负荷超过极，时，\par
一无法完整模假对计算者 B 的计算之计算，\par
计算者 B 会被迫选择「漏。」这一现实。\par
能认知到发生漏。的，不是计算者 A 而是计算者 B 这边。\par
在此情况下，计算者 B，就是我。\par
而计算者 A，并不是谁都可以。\par
最重要的条件是第 3 个。\par
条件之 3：计算者 A，是「来点」 。\par
「发射临近心中上上下下，不过这里有，喜。!」\par

% --- PDF page 46 ---
45\par
屏幕画面切换。\par
映出一个人。\par
!?什么…………\par
似乎只有我吃，。\par
大是的脸都像是已经知道了一样。\par
「呵呵，制作人，还好 ?」\par
「我的制作人，是觉真好。」\par
「制作人，要好好见证。。 」\par
「我出发了。 」\par
「呵呵，制作人，还好 ?」\par
对着镜绪说着话，同时思，。\par
一周后即将进行的发射，有着与正式发表所不同的隐藏目的。\par
隐藏目的，是模假设说的验证。\par
引发这个世界的，漏。。\par
对大是，持密。\par
漏。发生的 3 个条件。\par
条件之 1：某个计算者 A 的模假负荷超过极，。\par
条件之 2：计算者 A 的计算对象中，包含另一个计算者 B。\par
条件之 3：计算者 A，是「来点」 。\par
在模假设说中， 这个世界的计算全部， 都集中到被称为 「来点」 的唯一计算者身上。\par

% --- PDF page 47 ---
46\par
毕业公演那一天，与来到。众席的制作人对上了眼。\par
处理延迟漏。，就在那一瞬间发生了。\par
「我的制作人，是觉真好。」\par
呐，制作人。\par
制作人，就是「来点」 ?\par
或许全部都是错觉。\par
给模假施加负荷的方法，除了进行大量多边形计算，也可想到其他。\par
那就是，远离「来点」 。\par
「来点」 ，字面意思上也是坐标空间的来点。\par
所有多边形，最终，都在「来点」的坐标系中计算。\par
远离「来点」 ，那个坐标就会变得巨大。\par
然后在某处，计算者会处理不了过于巨大的数字。\par
「制作人，要好好见证。。 」\par
那时会发生什么，大『可，虑两种模式。\par
第 1 种可能性，是补码漏。。\par
类似虫。的东西。\par
当这个漏。发生时，会从宙空的一端瞬间移的到另一端。\par
如果正在远离「来点」 ，会不知何时靠近了「来点」 。\par
第 2 种可能性，是浮点漏。。\par
当这个漏。发生时，首，多边形会开始振的。\par
接着，多边形的消灭和扩大会同时被。测到。\par
作为第 3 种可能性，也有根本不会发生漏。的可能。\par

% --- PDF page 48 ---
47\par
模假设说的验证，非常困难。\par
即一如此，我也许，比人类史上任何人都更接近这个世界的真实。\par
「我出发了。 」\par
从开始远离「来点」 ，过的了多少时间呢。\par
突然，景色，开始振的。\par
比起「计算者」 ，我更喜欢表现为「选择者」 。\par
选择者，与选择其选择的选择者。\par
我和觉，通过选择相连。\par
我，做出选择。\par
觉，做出选择。\par
获得的热量\par
仍一直留存在眼底\par
被温柔地呼唤\par
结有一天\par
我内心如同泪水般炽热的风景\par
垮塌\par
触碰\par
颤栗\par
由我来 选择\par
——\par
泽广「光景」\par

% --- PDF page 49 ---
\noindent\textit{[第 49 页没有可抽取文字层。]}\par

% --- PDF page 50 ---
49\par
请\par
泽广教我锂离子电池\par
@pdncu\par
November 29, 2024\par
1 导入\par
广 「电化学是非常有趣的学术领域，制作人。 」\par
下叠大小的会议室里，回响着马克笔走行的声音。 在我眼前的她，是我担当的偶像\par
泽广。她最近开始戴的圆眼镜，我结感觉衬得她更加聪慧，可连白大褂都一穿上，一\par
活脱脱是我所不知的那个、 「14 岁从海外大学毕业」 的\par
泽广了。 与纤弱体型不太相称、\par
略无宽松的实验服，不知怎的在我眼中却非常贴合。\par
制作人正向担当偶像学习电化学的这一不可思议态况， 起因于她的出道演唱会周边。\par
作为周边开发的一环，才才开始销售移的电源，某 S 公司的电池就接连发生起火事件，\par
演变成了小小的社会问题。受到池鱼之殃，世间对电池的不安日益高涨。万一电池也起\par
火了， 怀恨在心的粉丝对担当偶像施加险害可怎么法。身为制作人，我想拥有知识来应\par
对「万一」 。然而，万万没想到身边最懂行的人竟是担当偶像本人。\par
P 『非常抱歉。本应是休息日，如果觉能为我抽出几个小时，实在感激不尽。』\par
广 「……，没关系。那我们马上开始吧。 」\par
2 基础的基础\par
2.1 记法与用语\par
广 「首，， 所谓电池 （battery） ， 是将能量转换为电的东西。 例如， 使用热电池，\par
就能把热转换成电。 太阳能板中含有将太阳光能转换为电的、 称为太阳电池的\par
东西。 」\par
广 「我们日常生活中使用的干电池等电池， 是所谓的化学电池 （electric battery）。\par
利用化学反应，将物质所具有的、以电子的能为主的化学能，转换为电。 」\par
广 「另外，构成电池的单元称为电化学单元（electrochemical cell）。」\par
P 『是个不习惯的词啊。』\par

% --- PDF page 51 ---
50\par
广 「单元， 是电化学的系统。 历史上， 电池是将多个单元连接起来的东西。 不过，\par
最近的电池， 尤其是日常使用的干电池和移的电源，大多由一个单元构成。 在\par
这一类电池上，可以说单元和电池这两个词具有互换性。 」\par
广 「这个觉应该在高中化学学过，单元由正极（cathode） 、负极（anode） 、电解\par
质（electrolyte）构成。 」\par
感觉她比平时话稍微多一些。 我正担心她节奏分配要不要紧， 她却已径自在白板上\par
运笔。大概是写惯了。\par
(+)|Electrolyte|(-) (1)\par
广 「竖线是两种物质的界面 （interface） 。 电池是包含多个相的系统， 重要的反应\par
大多发生在这里。后面会详细说明，请记住。 」\par
广 「说到底，电池为什么能蓄电，制作人觉明白 ?」\par
P 『是因为发生氧化还来反应（redox reaction）吧。』\par
广 「对。物质被氧化、还来，分别意味着电子的的放、接受。因此，反应中会产\par
生电子的移的。 」\par
LiCoO2 $\leftrightarrow$ Li1-xCoO2 + xLi+ + xe- (2)\par
C + xLi+ + xe- $\leftrightarrow$ LixC (3)\par
广 「式 2 和 3，是常见的锂离子电池中发生的、正极和阴极各自的半反应。之所\par
以是『半』 ，是因为将两个半反应合起来看是一个完整的氧化还来反应，但这\par
里分开来写。 」\par
广 「改变电极，半反应也会变。于是，整体的氧化还来反应也变化。就像拼图块\par
一样，可以构思出各种各样的化学电池。 」\par
广 「写做 e$^{-}$的， 是电子。 电子在阳极反应中被放出， 在阴极反应中被接受。放电\par
时放出电子的电极叫正极，接受电子的电极叫负极。电子在电极间交接。这就\par
是电子流，即电流。 」\par
1\par
广 「放着电池不管，这些反应就会向右进行，这就是放电（discharge） 。接上电\par
池后电器产品会的，就是因为电子自然地流的。反过来，从外部施加电压，如\par
果反应可逆， 就会发生逆反应，也就可以充电 （charge） 。 这种可逆性， 乃至以\par
充放电容量和速度为首的电池性能，很大程度上都取定于这个反应。所以， 电\par
极材料和电解质材料很重要。 」\par
P 『那么这些材料，都用些什么呢。』\par
1\par
译者注：此处作者来文全程都把阳极错写成了正极，译者已改正。\par

% --- PDF page 52 ---
51\par
广 「电极常用金属或液的金属、 碳、 半致体材料之类。 金属中会用到稀有金属等，\par
所以从可实续性的。点，有削减使用的 倾向。这是个叫做绿色化学（ Green\par
Chemistry）的领域。 」\par
广 「电解质使用含有氢、钠、氯等离子在内的液体。溶剂大多是有机的，碳酸二\par
甲酯（DMC）和碳酸乙烯酯（EC）价格低廉，经常使用。 」\par
电极基本是无机材料，电解质大多是有机溶剂，是这个意思吧。\par
广 「最近，离子致电聚合物\par
[1]\par
和金属有机框架(MOFs)\par
[2]\par
、固体电解质\par
[3]\par
等开始受\par
到瞩目。今天没法说明， 但它们各自都有普通电池不具备的特征，为迎合多样\par
需求，日复一日在进行开发。 」\par
2.2 电极反应模型\par
10 分钟过的了。即一如此，那个她竟然有能实续讲 10 分钟课的体力，一时难以置\par
信。日常训练的模样简直像设的一样。\par
广 「制作人，怎么了?」\par
P 『……不，请继续。』\par
广 「…。 接下来， 我们仔细看看电极反应 （electrode reactions） 。 所谓电极反应，\par
是在电极与电解质界面附近发生的反应的结称。 氧化还来反应也包含在内。 电\par
池几乎所有的行为都由这些反应定定。因为，产生电位差的只有这个界面。 」\par
广 「依循最简单的模型，电极反应可分为 4 个过程\par
[4]\par
。按距离电极表面由近到远\par
的便序来说明。 」\par
广 「首，是电荷转移（charge-transfer）反应。即氧化还来反应。电子在电极与\par
电解质的界面之间移的。 」\par
广 「其次是吸附与脱附（adsorption, desorption）。氧化物或还来物在电极表面\par
吸附或脱附。 」\par
广 「接着，氧化物或还来物从界面向电池中心（体相）移的时，会发生化学反应\par
（chemical reactions）。」\par
广 「最后是质量传递（mass-transfer）过程。物质从界面向体相电解质移的时，\par
会产生扩散、迁移、对流。在大部分电极反应中，质量传递与其他电极反应相\par
比要慢很多。 」\par
第二个的吸附和第三个的化学反应，似乎尤其会影响电池性能。\par
广 「以上就是学习电化学所需的基上信息。 电化学是为了说明靠电极反应才可能\par
实现的各种现象的学问， 所以囊括了用来分别深挖上面四个过程的理论。 每个\par
都稍微说明一下。 」\par
P 『请多指教。』\par

% --- PDF page 53 ---
52\par
3 电化学101\par
3.1 电荷转移\par
广 「第一个是电荷转移反应，。。把氧化还来反应再深挖一下。 」\par
广 「电荷转移反应要进行， 需要电位差。 电位差越大，电池内的电子就广进得越\par
快。虽不同于力，但可以想象成河水的流的。电流从高处流向低处。电子则往\par
这相反方向前进。这个高低差存在于电池内部，也可以从外部施加。 」\par
广 「电势Ecell通过下面的热力学公式，与系统的（自由）能量联系起来。在电化\par
学中，Ecell被称为电的势（emf）。」\par
$\Delta$G$\ominus$ = -nFEcell\par
$\ominus$ (4)\par
Ecell = Ecathode - Eanode (5)\par
广 「不过， 这个式子仅在标准态的，即所有物质都理想地参与反应的态的下才成\par
立。 上标的圆圈表示标准态的。\par
2\par
非标准态的时关系会不同。 但至少， 电势Ecell\par
增加，自由能就减少。反之亦然。 」\par
广 「也就是说， 将电势设为负时， 电子所具有的能量增加， 会据据空着的好高能\par
级。 这个能级多在电解质一侧， 所以电流从电极流向电解质。 这就是还来电流。\par
相对于此，氧化电流是能级下降，电流流向电极一侧。 」\par
P 『也就是说，氧化还来电流是跨越电极-电解质界面产生的啊。』\par
广 「…。 这就是氧化还来反应的基上。正因为产生于界面， 吸附和脱附才会涉及\par
电极反应。 」\par
广 「要解析氧化还来反应，有热力学方法和反应的力学方法。 」\par
3.1.1 热力学\par
P 『热力学 （thermodynamics） 是以平衡系统为对象的学问吧。 电池充放电时有\par
很多电子在的，应该不能说处于平衡态的吧?』\par
广 「好问题。正是如此。但是，热力学是非常实用的理论，所以如果可能的话我\par
想用在我的分析当中。因此，设设电池进行无，缓慢的充放电， 就可以使用热\par
力学。当然，这并非和实际电池完全相同的模型，但即使在此设设下，对电池\par
中好慢的部分也能得到一定的准确性。这个 是叫做 实用可逆性（ practical\par
reversibility）的想法。 」\par
2\par
译者注：国内教程中表示标况的符号一般使用「$\ominus$」 ，故作此更改。\par

% --- PDF page 54 ---
53\par
广 「所谓无，缓慢， 是将电流无，减小后， 就可以设设电池在维实平衡态的的同\par
时进行充放电。因为能在不同的平衡态的间往来，充放电就成为可逆过程\par
（reversible process）。」\par
P 『是准静的过程（quasistatic process）吧。上课时稍微提到过。』\par
广 「对。作为下一个设设，引入电极反应特有的、源于经验规则的设设，即质量\par
传递反应为速率控制步骤。此时，电荷转移反应相对好快，所以在任意电极电\par
位下， 氧化物和还来物处于平衡态的。因此， 可以定义电荷转移反应的平衡常\par
数Krxn。」\par
广 「这样就可以引入热力学的式子了。 致出被称为电化学最基上式子的 『能斯特\par
方程』 ，来结束热力学的说明。 」\par
广 「使用式 4。自由能是电荷与电势差之积，这从量纲分析就能理解。想知道该\par
式背后理论的话，请参，本科初级水平的热力学课程。 」\par
广 「接下来， 在自由能与电荷转移反应平衡常数Krxn之间， 有与阿伦尼乌斯公式\par
相似的关系式成立。这个式子可从麦克斯韦关系式直接致出。 」\par
-$\Delta$G$\ominus$ = RT ln Krxn (6)\par
广 「由式 4 和式 6，可以得到能斯特方程。不过，下面写的式子仅在标准态的下\par
成立。非标准态的时，浓度会变为一个叫活度（activity）的。，但这里，不的\par
管它。 」\par
E = E$\ominus$ + RT\par
nF ln K (7)\par
广 「E$\ominus$ 是标准电极电位。指反应相关物质以理想用量使用、且处于平衡态的时\par
的电极电势。实际上，需要用一个叫表。电位E$\prime$$\ominus$（formal potential）的、略\par
微依赖电解质离子行度的。。 」\par
广 「尤其，当将一般的氧化还来反应表示为O + ne$^{-}$ $\leftrightarrow$ R 时，能斯特方程表示如\par
下。 」\par
E = E$\ominus$ + RT\par
nF ln CO\par
∗\par
CR\par
∗ (8)\par
广 「C∗是电解质中心的浓度，称为体相浓度。注意与表面浓度不同。 」\par
广 「能斯特方程于 1889 年致出，以此为起点，各种理论发展了起来。 」\par
P 『用这个式子，就能说明所有热力学现象 ?』\par

% --- PDF page 55 ---
54\par
广 「并非如此。 仅靠能斯特方程， 无法完全模型化电池。略微举一些应，虑的要\par
素： 电解质的离子行度、 电极-电解质界面间电势的变化、 界面双电层的形成、\par
活性物质的吸附、电解质内的极化效应，不胜枚举。但是，能斯特将电势在一\par
定程度上准确且简单地表达出来，功绩是很大的。 」\par
3.1.2 反应动力学\par
电流与浓度的关系式\par
广 「想用热力学解析电池也有局，。 因为界面发生的反应中， 有的时间尺度短到\par
无法视为准静的。 」\par
广 「因此，也需要使用反应的力学（kinetics） 。相对于处理平衡系统的热力学，\par
反应的力学『也』能处理非平衡态的的系统。是种自由度好高的理论，取某些\par
变量的极，时也能表现平衡态的。 」\par
P 『取极，时，必须和热力学的结论一『呢。 』\par
广 「对。 因为在极，处热力学和反应的力学相连， 所以电化学的反应的力学是从\par
能斯特方程出发的。 」\par
广 「能斯特方程是电势与浓度的关系式， 但为了发展反应的力学， 我们想建立表\par
示电子流的速度的电流与浓度的关系式。 因为电流的大小， 是由电极反应致『\par
的电子转移（electron-transfer）和以 Li$^{+}$为首的质量传递（mass-transfer）速度\par
定定的。 」\par
广 「这时在 1905 年，一位叫塔菲尔的化学是，从经验规则致出了关于电流和电\par
势的式子。 」\par
$\eta$ = a + b ln(i) (9)\par
广 「式 9 的$\eta$是一个叫过电位 （overpotential） 的。。 由能斯特方程得出的电池电\par
势设设反应是理想的， 但通常电极反应并不理想， 会在电极表面不均匀地发生，\par
所以实际。会有偏差。这个偏差就是过电位。 」\par
P 『电极表面如果没有达到来子级平滑， 就不会完全均匀。 可以想象，要做到完\par
全平滑需要极大的成本。』\par
广 「塔菲尔方程是连接电压、电流和热力学的式子。构建反应速率理论的的机，\par
正是希望建立一个当极，条件下的该公式能够得出结论的反应论模型。」\par
广 「这里，引入关于反应论的设设。氧化还来反应的正反应速度如下，以浓度的\par
一次式表示。 」\par
vf = kfCO (10)\par

% --- PDF page 56 ---
55\par
CO(t) = lim\par
x→0\par
CO\par
∗(x, t) (11)\par
广 「这里浓度记为CO\par
∗(x, t) ，是以距电极表面最短距离和时间为自变量的双元函\par
数。这里暗含了反应在 y,z 轴向上均匀发生的设设。另外，注意区分电极表面\par
的浓度CO和体相浓度CO\par
∗。」\par
P 『才才觉提到了，作为电极反应的第二个要素是副反应。 副产物的浓度等， 不\par
用，虑 ?』\par
广 「这是出于这样的想法： 电极-电解质界面附近的行为极大地支配着电极反应，\par
所以只着眼于距离界面最近的氧化物O、还来物R的浓度就行。 」\par
广 「所以现在不用在意。 不，虑全部 4 个过程的相互作用会没完没了。 在质量传\par
递反应和吸附中，有些部分会受到好大影响，因此相互作用的模型化在广进，\par
但与反应的力学的相互作用大多没那么大。 」\par
P 『来来如此。』\par
广 「回到正题。 既然的机是求电压与电流的关系式， 就结得设法将反应速度和电\par
流联系起来。这里，还来电流ic和以下关系式成立。 」\par
vf = kfCO = ic\par
nFA (12)\par
广 「n是电子的摩尔数，A是电极表面积，从量纲分析上可以接受。 同样， 逆反应\par
速度vb与氧化电流ia的关系如下。 」\par
vb = kbCR = ia\par
nFA (13)\par
广 「由这两个式子，可以计算净反应速度和电流，吧。 」\par
vnet = kfCO - kbCR = inet\par
nFA (14)\par
inet = ic - ia = nFA(kfCO - kbCR) (15)\par
广 「重复一遍， 这个式子在反应均匀发生时成立。 不均匀的情况下， 表面积A 的\par
部分会变得更加复杂。那种情况下，必须根据系统设计独特的式子。不过，大\par
是都经常从这个式子出发的建模。 」\par
巴特勒 -福尔默 (Butler -Volmer )方程\par
广 「电流与浓度的式子出来了， 代入能斯特方程， 就能得出电流与电势的式子。 」\par
P 『前者用了表面浓度，后者用了体相浓度吧。』\par

% --- PDF page 57 ---
56\par
广 「对。所以不能直接代入。最基上的电流-电势公式叫巴特勒-福尔默方程，需\par
要自由能图的初级几何，改，所以省略。想知道的话，参，合适的课程视频就\par
好。直接从结论开始，写。 」\par
i = FAk0 [COe-$\alpha$ F\par
RT(E-E$\ominus$) - CRe(1-$\alpha$) F\par
RT(E-E$\ominus$)] (16)\par
广 「说明一下新参数。k$_{0}$是标准反应速率常数，当E = E$\ominus$时， 也就是正反应和逆\par
反应速度一『时，各自的速度常数kf = kb = k0成立。 」\par
广 「$\alpha$叫转移系数（transfer coefficient） ，是对应于氧化还来反应能垒对称性的\par
参数。 这个参数基本很难测量， 所以常设为$\alpha$ = 0.5。 在大多数系统中取。在0.3\par
到 0.7 之间。 」\par
广 「这个式子的厉害之处在于， 可以由电化学系统中容易测量的电流和电势， 计\par
算出参数，分别得到关于反应常数和能垒的信息，非常方一。而且，过电位好\par
小时， 电流和过电位会呈现线性关系， 过电位好大时，两者会以简单指数函数\par
相连，有种种好事。 」\par
P 『对实验者来说，真是个令人高兴的式子啊。』\par
广 「不过要注意的是，这个式子完全没有，虑质量传递和吸附等其他电极反应的\par
部分。对于各部分无著的体系，要扩展巴特勒-福尔默方程来应对。 」\par
广 「还有，至今为止处理的氧化还来反应都是一步反应，电子系数为1。大部分\par
系统是多步的，电子系数也不是 1。接下来，，虑多步、多电子氧化还来反应\par
时，如何扩展 B-V 方程。 」\par
向多步、多电子反应的扩展\par
广 「即一是最简单的氢还来反应， 也是2H+ + 2e- → H2， 并非单电子反应。 在一\par
般的电化学反应中， 移的一个电子需要一步， 关于n 个电子的移的， 存在n 步\par
的电荷转移反应，这么看应当是 2 步。即使是Ag+ + e- → Ag，在电荷转移之\par
外，也应该存在电子进入晶格的过程。 所有的电极反应都是复杂的，根据反应\par
不同，存在不同的模型和实验方法。 」\par
广 「如何进行近似呢。 基上的想法是速率控制步骤， 简称速控步， 即找出各步骤\par
中速度最慢的一步。 特别是电极反应的速控步， 大多是不均匀的电荷转移反应。\par
那是移的一个电子的反应，可以来样适用以前的讨论。 」\par
广 「举几个能简化的例子。最简单的例子是，整体反应处于平衡态的时。这时，\par
各步骤也处于平衡态的。 因此， 表面浓度CO\par
$\prime$ 、CR\par
$\prime$ 分别与体相浓度CO\par
∗、CR\par
∗平衡。\par
变形平衡常数， 就能对表面浓度应用能斯特方程。 这在化学上可逆且平衡态的\par
得以持证的所有多步反应中成立。 」\par

% --- PDF page 58 ---
57\par
广 「即使反应整体不处于平衡， 只要电极反应足够快， 氧化物和还来物之间就存\par
在平衡，可以使用能斯特方程。将『反应可逆』 『电荷反应非常快』的条件汇\par
结，这称为能斯特方程的应用条件。」\par
广 「当多步反应既不在平衡态的， 也不符合能斯特方程应用条件的时候， 反应的\par
力学极大地关乎系统行为， 因此需要从行为广断反应速度相关变量和反应机理。\par
处理均相反应的力学时， 主要步骤是广断机理， 比好电流与电压关系式的言测\par
与实际数据。建立依赖可控变量、如参与反应物质浓度等的言测。 」\par
P 『这是对未来实验者的关怀呢。 』\par
马库斯理论\par
广 「巴特勒-福尔默方程依赖于 $\alpha$ 和 k$_{0}$， 并不直接将反应物结构、 溶剂、 电极材\par
料与速度变化联系起来。 要讨论这些， 就必须从微。上理解电子转移反应。 鲁\par
道夫$\cdot$A$\cdot$马库斯于 1992 年凭此理论获得了诺贝尔化学奖。 」\par
广 「不过，因时间关系，这里就不展开说明了。有兴趣的话请自行查阅。因为之\par
前的讨论都不依赖于电解质等材料， 所以凭借这个理论， 才首次能展开对每种\par
不同电池的，改。 」\par
广 「那么下一个」\par
P 『\par
、\par
泽同学，请等一下。要不要稍微休息一下?』\par
广 「——啊，抱歉。已经过了 40 分钟了。 」\par
她在准备运的时都会喘不上气， 到底哪里有讲课的能量呢。 我为了冷却不停写字的\par
右手，也伸手的拿塑料瓶的绿茶。\par
3.2 质量传递\par
广 「重新开始，吧。 」\par
广 「质量传递反应\par
3\par
是速率控制反应， 所以对整体反应有很大影响。 这个反应的解\par
析非常重要。 」\par
广 「说明一下思，这个反应所需的、称为化学势的概念。化学势是每种物质i 各\par
自拥有的。，如下表示。 」\par
$\mu$i = ( $\partial$G\par
$\partial$ni\par
)\par
P,T,nj$\neq$i\par
(17)\par
3\par
译者注：此处译者认为传质称为「过程」更为合适。此差异无伤大雅，故持留来文。\par

% --- PDF page 59 ---
58\par
广 「可以直。地解的一下这个式子 ，制作人?」\par
P 『物质浓度ni变化时，自由能变化越大，化学势就越大。是各物质对系统能量\par
有多大影响的指标吧。』\par
广 「对。进而，电化学将其扩展后，叫做电化学势。因为电荷和物质也会因所在\par
相的电势差而移的，所以需要新的势。 」\par
$\mu$ = $\mu$ + z$\phi$F (18)\par
广 「z是离子价数，$\phi$是相间电势差，F是法拉第常数。物质会朝这个势减少的方\par
向流的。与电位相同。不如说， 电化学势是同时，虑了电位差和由浓度差引起\par
的质量移的的势能。出色的东西。」\par
广 「把这个放在心上。 质量传递反应存在3 种。 一种是扩散 （diffusion），由浓度\par
梯度引起。第二种是迁移（migration），由电位梯度引起。这些都可用电化学\par
势差来，虑。第三种是对流（convection），由振的、搅拌等扰的引起的浓度梯\par
度所『。 」\par
广 「致出综合，虑这三种反应的一般公式后就结束。 物质j的通量（flux）会使化\par
学势的梯度减少。 通量J$\vec{}$j与化学势的梯度成正比。 以J$\vec{}$j $\propto$ $\nabla$uj表示。 这里的比例\par
系数，可以使用扩散系数Dj[m2/s]和表面浓度Cj，由经验规则致出。 」\par
J$\vec{}$j = - CjDj\par
RT $\nabla$uj (19)\par
广 「R 是理想气体常数，T 是温度。负号，是因为朝化学势增加方向的反方向流\par
的。 」\par
P 『这不是菲克定律 ?』\par
广 「这个式子是在扩散基上上还，虑了迁移。菲克第一定律包含在这里面。 」\par
广 「进而，溶液在点s处时以速度v$\vec{}$运的，所以需要与浓度梯度相关的项Cjv$\vec{}$。这\par
对应于对流。 」\par
J$\vec{}$j = - CjDj\par
RT $\nabla$uj + Cjv$\vec{}$ (20)\par
广 「最后，化学势与标准态的的。关系为uj = uj\par
$\ominus$ + RTlnaj，所以电化学势在点\par
s 处， 」\par
uj(s) = (uj\par
$\ominus$ + RT ln aj(s)) + ziF$\phi$(s) (21)\par
广 「成立。那么，由上面两式可得下式。 」\par

% --- PDF page 60 ---
59\par
Jj = -Cj$\nabla$Dj - ziF\par
RT DjCj$\nabla$$\phi$ + Cjv (22)\par
广 「这叫能斯特-普朗克方程。是模型化质量传递的最基本公式。当物质j带有电\par
荷时，通量Jj等于电流密度。由此也能致出与电流的关系式。 」\par
P 『来来如此， 因为电化学势中存在活度项， 展开后可看出第1 项对应菲克第一\par
定律呢。 』\par
3.3 氧化物、还原物的化学反应\par
广 「关于这个， 不同电池的行为差异很大。需要根据各个系统，构建实验色彩浓\par
厚的理论。这里不触及。 」\par
3.4 双电层与吸附\par
广 「电荷转移反应是氧化还来反应。 此反应要发生， 需要足够的自由能以超越活\par
化能，由式 4 知即需要电势。 」\par
P 『没有足够电势时，电池中会发生什么?』\par
广 「实际上， 会有微小的氧化还来反应发生， 或杂质被氧化还来。 什么都不发生，\par
是理想情况。这样的电极叫理想极化电极（IPE）。」\par
广 「即使反应不进行，不施加不充分的电势，电荷也会移的。但因为不反应，电\par
子即电流不会流的。 」\par
P 『像电容器一样呢。 』\par
广 「正是因此，认为在电极-电解质界面产生了类似电容器的机制。在任意电势\par
下， 金属电极上产生电荷qM， 电解质中产生电荷qS。 这个。的正负取定于外部\par
电压。 」\par
广 「qM由电极中电子的过剩或不足引起，是厚度小于 0.1 Å 的层存在。 这可以类\par
比为电荷堆积的、 带负电的金属板。qS由存在于电极表面周围的阳离子或阴离\par
子的过剩或不足引起。这可以类比为电荷不聚集的、带正电的金属板。 两者均\par
常除以电极表面积，表示为电荷密度$\sigma$M = qM/A等。单位是$\mu$C/cm2。」\par
广 「为了在界面整体呈电中性，qM = qS成立。这种类似电容器的机制叫双电层\par
（electric double layer） 。 」\par
广 「然后， 当离子等子子依电势从电解质侧移的到达界面时， 界面处的浓度会增\par
加。就像电荷集中于电容器的平板上一样。这叫吸附（adsorption）。」\par
广 「吸附有两种。 与带负电金属的相互作用是在好长距离产生的静电性的， 所以\par
仅依赖于电荷，不依赖于化学性质。 」\par
P 『库仑定律只在公式中包含电荷量，不，虑化学性质呢。 』\par

% --- PDF page 61 ---
60\par
广 「对。 这样的离子称为发生了非特性吸附 （non-specifically adsorbed） 。 与此相\par
对，依赖于子子化学性质的吸附叫特性吸附。 」\par
广 「经验规则提出， 双电层及离子的特性吸附会影响电极反应速度， 并构建了说\par
明此现象的理论。 」\par
4 最近锂离子电池的发展\par
广 「制作人，辛苦了。才才的就是电化学理论的介绍。 」\par
P 『非常感谢。不过，比起感觉上，更多是经由数学理解的部分啊。』\par
广 「…。 在所有的物理和化学领域， 都有从事理论工作的人和从事实验工作的人。\par
从这里开始，以实验说明。我想，大概在这里可以做到与实际生活相结合的理\par
解。 」\par
广 「电池所需的要素大『分为 6 个。 能量密度、 功率密度、 安全性、 寿命、 性能、\par
成本\par
[5]\par
。要迎合需求做出特化，或是要取得良好平衡，就必须尝试各种各样的\par
材料和制了方法。 」\par
P 『最近用 AI 似乎就能做到啊。』\par
广 「也看到了以 Google GNoME 为首的、利用机器学习技术来发现优良材料的\par
的向，但电池涉及的要素太多，还无法做到精度良好的模型化。现阶段，仍是\par
人类科学是在通过试错努力寻找。 」\par
广 「在，虑这 6 个要素的基上上，为什么使用锂离子呢?列举其优点：还来电压\par
低、 元素中第三轻、 离子半径小、一价\par
[6]\par
。 尤其是多价离子虽能带来更多容量，\par
但迁移性能无著下降， 所以不是主流。 这些优点在进行插层反应时都很重要。 」\par
P 『从资源方面看，也很难想象锂将来会匮乏。 』\par
广 「不过，有几个缺点。首，，在锂金属和正极所用到的钴等的开采中，童工和\par
腐败等社会问题，以及采矿劳的的健康问题等，被多份评估报告指出\par
[7]\par
。使用\par
钠的电池，除了不用稀有金属等带来的成本优势外，因无需采矿等，供应链风\par
…也相对好低，开发正在广进中。」\par
广 「第二个，容易劣化。因插层反应会产生各种劣化。例如电解质的劣化、活性\par
物质的减少、锂离子量的减少等\par
[8]\par
。要应对各个劣化因素，就需要掺杂或更改\par
材料本身等。 」\par
广 「另外，使用液体电解质的锂离子电池有爆炸的险…性。这个后面详细说明。\par
将电解质从液体更改是直接的对策，如Li-Po 电池和全固的电池等。还有，调\par
查副反应的路径，通过添加或的除什么，以防产生气体等。 」\par

% --- PDF page 62 ---
61\par
广 「触及一下近年的研究方向。首，，负极材料相对好少被人着手。因为已经找\par
到了石墨这样能量密度非常高、 成本极低的材料。虽然主要针对正极材料和电\par
解质进行了协同优化，但整体改进力度好小。」\par
广 「作为关于负极的研究事例之一， 有一种无需负极的无负极电池。这种电池在\par
循环过程中会自然形成锂金属负极， 从而能够降低制了成本、 减少工序和减轻\par
质量。尤其在持实储能能力的同时减轻重量，可无著提高能量密度。目前，多\par
是机构正在开展相关研究，但均处于试制阶段\par
[9]\par
。」\par
广 「正极正在尝试多种类型， 最主流的是过渡金属氧化物LiMO2和聚阴离子化合\par
物LiXO4两种。此外，硫化锂、 碘化锂，乃至硒之类的稀有金属都在尝试\par
[6]\par
。也\par
有方向是混合两种正极材料，以此取得各自的优点\par
[10]\par
。 电池的劣化比好大地依\par
赖于正极材料。能量密度与功率密度直接相关。 成本也很大程度上取定于是具\par
使用稀有金属。可以说正极材料是电池的关键。」\par
广 「电解质大多是不耐温度变化的， 所以被大量研究。 市售产品使用碳酸乙烯酯\par
或碳酸二甲酯， 但这些只能耐受-20 到 60℃左右\par
[11]\par
。 要驱的在地下或宙空等更\par
极端环下中使用的机器，就需要扩大作用温度的范围。2019 年，开发出了能\par
耐-125℃到 70℃的电解质\par
[12]\par
。」\par
5 电池爆炸的机理\par
广 「基于以上讨论，对电池爆炸进行，改。 」\par
广 「作为全世界召回金额最高的智能手机也闻名的某 S 公司电池爆炸事例中， 据\par
第三方调查结果，以下两个要素是主要来因\par
[13]\par
。」\par
广 「许多电池的右上部分存在细小凹陷。因为这个凹陷， 存在于正极和负极之间\par
的隔膜发生了歪曲。 」\par
广 「另外，谋求电池小型化时，削减了隔膜的厚度。 」\par
广 「这两个因素致『电池内部位于正极和负极下方的集流体相互接触， 从而发生\par
短路。主回路产生的焦耳热对爆炸也稍有贡献。」\par
广 「然而，在许多锂离子电池中， 对爆炸起主要贡献的，是来本不会发生的副化\par
学反应所放出的热，这会升高整个电池的温度\par
[14]\par
。这个反应是在正极或负极、\par
或两极同时发生的放热反应。超过临界温度后，热会进一步增加电流， 内部电\par
阻产生的热又进一步增加电流， 引发称为热失控 （thermal runaway） 的现象，\par
变得无法遏制，致『火灾或爆炸\par
[15]\par
。」\par
P 『最近也有看到掉落手实风扇结果爆炸的事例，大概也是同样的来因吧。 』\par
广 「对。可以认为，大多是由于力学上的凹陷损伤了隔膜，内部发生短路，最终\par
致『热失控。 」\par

% --- PDF page 63 ---
62\par
广 「最后是关于会成为我的周边的移的电源，这个叫Li-Po 电池，电解质使用了\par
凝胶态聚合物。因为中间有凝胶，与一般锂离子电池所用的液体电解质相比，\par
能在一定程度上起到缓冲作用， 所以不易短路。另外，耐热聚合物电解质有很\par
多\par
[16]\par
，即使不是这样，改善到能耐热也好容易\par
[17]\par
。所以，说它更安全。 」\par
P 『是这样啊。太好了。 』\par
广 「不过，无法持证不会短路。严重变形时也可能会贯穿隔膜了成短路， 或是施\par
加巨大外部热量、高电流时，还是会热失控。 」\par
P 『结之，正确使用电池的话，爆炸的可能性就接近 0 吧。』\par
广 「就是这么回事。 」\par
6 总结\par
广 「在这几个小时里， 解说了电化学的基上理论、锂离子电池的研究事例，以及\par
爆炸的机理。 」\par
P 『非常感谢。结算是，对这个领域有了大概的印象。最在意的爆炸话题，也加\par
深了理解。』\par
广 「电化学是理论和实验所要求的东西相当有区别的领域之一。 理论同时需要物\par
理和化学，实验也需要工程知识，因此有各种背景的人聚集，很有意思。 」\par
广 「以电池为首的电化学子领域， 无一不是对能源问题的直接对策， 能非常行烈\par
地感受到价。。只要活着，无论是谁都会使用能源，所因此能够开发出可能成\par
为全人类的幕后支实者的技术。」\par
广 「非常快乐，所以希望觉也来。 」\par
P 『这话是对我说的 ?』\par
广 「……。制作人是制作人。 」\par
P 『...???』\par
完\par

% --- PDF page 64 ---
63\par
参考文献\par
[1] Ye J, Yuan D, Ding M, et al. A cost-effective nafion/lignin composite membrane with\par
low vanadium ion permeation for high performance vanadium redox flow battery[J].\par
Journal of Power Sources, 2021, 482: 229023.\par
[2] Dutt S, Kumar A, Singh S. Synthesis of metal organic frameworks (MOFs) and their\par
derived materials for energy storage applications[J]. Clean Technologies, 2023, 5(1):\par
140-166.\par
[3] Wang M J, Kazyak E, Dasgupta N P, et al. Transitioning solid -state batteries from\par
lab to market: Linking electro -chemo-mechanics with practical considerations[J].\par
Joule, 2021, 5(6): 1371-1390.\par
[4] Bard A J, Faulkner L R, White H S. Electrochemical methods: fundamentals and\par
applications[M]. John Wiley \& Sons, 2022.\par
[5] Stan A I, Ś wierczy ń ski M, Stroe D I, et al. 2014 international conference on\par
optimization of electrical and electronic equipment (optim)[J]. 2014.\par
[6] Nitta N, Wu F, Lee J T, et al. Li -ion battery materials: present and future[J].\par
Materials today, 2015, 18(5): 252-264.\par
[7] Thies C, Kieckhä fer K, Spengler T S, et al. Assessment of social sustainability hotspots\par
in the supply chain of lithium-ion batteries[J]. Procedia CIRP, 2019, 80: 292-297.\par
[8] Dubarry M, Anseán D. Best practices for incremental capacity analysis[J]. Frontiers\par
in Energy Research, 2022, 10: 1023555.\par
[9] Wu B, Chen C, Raijmakers L H J, et al. Li -growth and SEI engineering for anode -\par
free Li-metal rechargeable batteries: A review of current advances[J]. Energy Storage\par
Materials, 2023, 57: 508-539.\par
[10] Whitacre J F, Zaghib K, West W C, et al. Dual active material composite cathode\par
structures for Li-ion batteries[J]. Journal of Power Sources, 2008, 177(2): 528-536.\par
[11] Ji Y, Zhang Y, Wang C Y. Li -ion cell operation at low temperatures[J]. Journal of\par
The Electrochemical Society, 2013, 160(4): A636-A649.\par
[12] Fan X, Ji X, Chen L, et al. All -temperature batteries enabled by fluorinated\par
electrolytes with non-polar solvents[J]. Nature Energy, 2019, 4(10): 882-890.\par
[13] Loveridge M J, Remy G, Kourra N, et al. Looking deeper into the Galaxy (Note 7)[J].\par
Batteries, 2018, 4(1): 3.\par
[14] Ren D, Feng X, Liu L, et al. Investigating the relationship between internal short\par
circuit and thermal runaway of lithium -ion batteries under thermal abuse\par
condition[J]. Energy Storage Materials, 2021, 34: 563-573.\par
[15] Conte F V, Gollob P, Lacher H. Safety in the battery design: the short circuit[J].\par
World Electric Vehicle Journal, 2009, 3(4): 719-726.\par
[16] Armand M. The history of polymer electrolytes[J]. Solid State Ionics, 1994, 69(3-4):\par
309-319.\par

% --- PDF page 65 ---
64\par
[17] Chen M, Yue Z, Wu Y, et al. Thermal stable polymer-based solid electrolytes: Design\par
strategies and corresponding stable mechanisms for solid -state Li metal batteries[J].\par
Sustainable Materials and Technologies, 2023, 36: e00587.\par

% --- PDF page 66 ---
65\par
关于不相溶的东西。 或者说是， 共晶系统的统计力学正当化。\par
休息日的午后。虽说是休息，但那是指练习，课还是照常上。所以到了下午就闲了\par
下来， 每周的这一天， 我都会走到制作人所在的这个房间。 制作人正对着电脑在做什么。\par
现在连打字声也停了，我直直地盯着他的脸， 他也依旧面无表情纹丝不的。 这种时候的\par
制作人并不是在思，，而是明知我在看他，却故意无视我。\par
对视游戏还是稍后再说吧。 我把视线落回手绪的平板电脑上，上面无示的是一个小\par
小的图表。是佑芽给我看的，硫酸钠的溶解度曲线。她说，昨天的课上讲了作为温度函\par
数的溶解度，复习的时候就发现了这个。\par
很多物质的溶解度会着温度单调变化。 盐之类的是单调增加， 二氧化碳则是单调减\par
少。而硫酸钠则以 32.4℃为界，从增加转向减少。她想知道这是为什么。\par
解的本身我当然会。 但是， 要用佑芽也能理解的前提知识……而且还得说得不太冗\par
长，我不知道该怎么做。所以我坐在椅子上，要么默默思，，要么把视线投向制作人的\par
脸，可是还是想不出来。\par
「制作人。我要讲个本科水平的统计力学课，觉起着。有问题我会回答。 」\par
「……怎么突然要讲课了?」\par
「制作人，能增加知识，难道不快乐 ?」\par
呵呵，一副无奈的表情。\par
言归正传，我想，既然想不出好方法，不如，把信息铺展开来，找个突破口。基于\par
统计力学的理解， 能够在某种程度上定性地明白， 那个平衡点往往就在那里。虽说如此，\par
我也不能为了这个就生硬地搬出什么『根据玻尔兹曼\par
1\par
来理S = kB log $\Omega$ 2 这一关系式成\par
立……』之类的说法来。要教，就应该从这个式子的广致开始，对吧。……kB是玻尔兹\par
曼常量，$\Omega$是组合数， 这还算比好直。， 但熵S的浅无说明就……结之，广致一下， 看看\par
需要多少信息吧。\par
「制作人。物理知识，觉懂多少?」\par
「我且且算是现役文科生。。高中物理那种简单的，我想我应该还是理解的。 」\par
「知道了。那就，教教千奈和佑芽时一样来。制作人坐我旁边。 」\par
虽然在白板上写写画画我还可以， 但要是不趁能休息的时候彻底让身体休息， 我再\par
次倒下也不奇怪。\par
1\par
译者注：作者来文中均持留了人名的英文，我再此处根据个人习惯译为了中文人名\par
2\par
译者注：此处来文为kB log W，国内教材一般使用$\Omega$代表组合数，故作此修改\par

% --- PDF page 67 ---
66\par
打开手写笔记应用，新建了空白页。要教别人东西，传达意象是最快的。这是佑芽\par
教给我的其中一件事。受过理科教育的人之间，只要有算式就能共享意象， 但对于不是\par
这样的人， 大概画图表示是最好的。 我在上面画了些线条构成方框， 里面画了好几个点。\par
其中一个点旁边还添了个向右的箭绪。\par
「统计力学致入时常用的就是这种图。 有一个子子， 子子………， 就当成小应好了。\par
小应在里面飞来飞的。子子里被墙壁围着，完全不受外界影响。这个时候，这样的……\par
说得素素一点，存在一个可以合理化为『 仅仅将其复制到别处也会表现出相同行为』的\par
特定区域，物理学里就叫做系统。那么首，，，虑子子里只有一个子子的情况。……要\par
是在分析力学的语下下，可能会想到谐振子………」\par
能懂 ?我用眼如问制作人。他普普通通地摇了摇绪。\par
「觉就把子子里的东西想成弹力应。 它弹跳到顶点附近时的作会非常慢， 但基本上，\par
通过弹跳它不会损失能量，也不受空气阻力。 」\par
「\par
泽同学，觉知道弹力应这种东西啊?」\par
「…。前子子的典典的时候，我玩了弹弹力应。……一个都没弹到。呵呵。 」\par
「我想也是。是花海同学大收收之后分给觉的，对吧?」\par
「答对了。很开心，。。然后，那个弹力应所具有的能量E，看作力学能量就好，所\par
以可以认为是的能和势能的和，设质量为m，位置为x，速度v = dx/dt，重力加速度为\par
g 的话，\par
$\epsilon$ = mv2\par
2 + mgx (1)\par
就可以结结为这个式子了。从能量守恒定律来说，这个结是持实恒定，对吧。 」\par
长时间的嘴是很累的。作为休息，我把式子写在了笔记上，但式子本身也不是什么\par
大不了的，一下就写完了。\par
我试着把绪靠在制作人的肩膀上，却被他无言地广了回来。不过这样一来， 我也能\par
把体重托付给他，也不算坏。，让制作人负担我一子子好了。\par
「在那种时候……关于某个态的的子子，就算能量一定，它也能取一定范围的。。\par
这就像轻轻放下弹力应和用力砸下的时， 结果不同一样。 ……制作人， 觉不看笔记这边，\par
我很困扰。。 」\par
就那样让他支撑着， 我猛地转过绪的看制作人的脸，结果他避开了。我的嘴角自然\par
地上扬。趁制作人的脸还没转回来，我把视线落回平板上。\par

% --- PDF page 68 ---
67\par
「然后……像这样，可以设想具有各种不同能量的弹力应。如果是弹力应，能量就\par
是从 0 到+$\infty$。这种时候，试想一下同时有很多个弹力应。设设一个非常大的子子里有\par
一定数量的弹力应， 它们几乎互不干涉， 所以每个弹力应的能量可以看作是各自大『独\par
立守恒的，给出这样一个系统。能想象 ?」\par
因为我还没法转过绪的， 不知道制作人点绪了没有。不过我想， 制作人的话肯定能\par
懂。呵呵。\par
「为了方一言，做一下准备， 我们也，，虑弹力应的结数N和能量的结和E = $\sum$ $\epsilon$i =\par
$\epsilon$1 + $\epsilon$2 + $\epsilon$3 + $\cdots$。 这个下标是指从能量低到高， 按适当宽度划分的一个个区域。 如果能\par
量是离散的，也可以把它想成那些离散的能级本身。这时候……。 」\par
「请等一下，那个操作是什么意图啊?」\par
「从这里开始， 再往下讲一点会更容易明白， 所以请再耐耐一下。 能量……比方说，\par
简单起见， 设设能量取-1、0、1 这三个。， 有16 个弹力应， 它们的能量结和是2 的话，\par
那么可以想出-1 有 7 个、1 有 9 个的组合；-1 有 6 个、0 有 2 个、1 有 8 个的组合；\par
5$\cdot$4$\cdot$7 的组合……0$\cdot$14$\cdot$2 的组合，一共 8 种组合。如果，虑在结数 N 和结能量 E\par
都不变的系统里，各个弹力应的能量是着机定定的，而且弹力应之间可以区分的话，才\par
才说的 8 种组合实际上还会变得更多。但是，这些组合并非均匀分配到那 8 种情况里，\par
而是会形成这样一种山峰一样的分布。和掷两个骰子时出现数。的分布，是一样的。。\par
然后， 当这个子子数大幅增加时，这座山就会在正中央变得非常尖锐。 ，虑到气体分子\par
之类的话，每 1 g 的子子数在1021到1022个的量级，所以除了频率最大的地方以外，其\par
他地方几乎都可以忽略不计了。\par
3\par
所以，这样分割能量，就能得出其中组合数最大的地\par
方，就是这么个道理。 」\par
我草草画了个正的分布，然后又画了个类似狄拉克$\delta$函数的东西。要找到这个频率\par
最。，用拉格朗日乘数法是最简单的……不过，对制作人还是算了。\par
4\par
「懂了 ?」\par
例子他好像懂了，但话题的核心似乎还没懂。等看到全貌之后，视野会更好，所以\par
我想，继续讲下的比好好，但是有没有更好的做法呢。\par
「然后，这种……由互相几乎独立的子子集合而成的系统，就叫做微正则系综。且\par
且不管结能量，，，虑对应于结数N的排布方式数$\Omega$。在前面那个同样的三能级系统的\par
例子里，设设进入各个能级的子子数分别是na、nb、nc，\par
3\par
另外，这里的前提是设设各能级被据据的概率是相等的。详情请搜索『等概率来理』或『遍历设\par
说』之类的\par
4\par
寻找待定常数 f、$\beta$的式子变形大概搜一下就能出来，有兴趣的话请查查看。或者着一找本面向统\par
计力学初学者的书翻开也行\par

% --- PDF page 69 ---
68\par
CN\par
na $\cdot$ CN-na\par
nb $\cdot$ CN-na-nb\par
nc = CN\par
na $\cdot$ CN-na\par
nb $\cdot$ Cnc\par
nc\par
= N!\par
(N - na)! na! $\cdot$ (N - na)!\par
(N - na - nb)! nb! $\cdot$ 1\par
= N!\par
na! nb! (N - na - nb)!\par
= N!\par
na! nb! nc!\par
(2)\par
同样地，对于多能级系统，也可以有\par
$\Omega$ = N!\par
n1! n2! n3! $\cdots$\par
log $\Omega$ $\cong$ N log N - $\sum$ ni log ni\par
i\par
(3)\par
这样的式子。\par
5\par
现在， 对各能级i，，虑该能级的子子数ni以及在该能级具有的能量$\epsilon$i，那么结子子\par
数N = $\sum$ nii 、结能量E = $\sum$ ni$\epsilon$ii 就可以成立了。用拉格朗日乘数法求解的话，利用玻尔\par
兹曼常量kB、温度T、待定乘数——虽然不知道内在是什么，但结之是常数的那个f，\par
ni = N\par
f exp ( -$\epsilon$i\par
kBT) (4)\par
可以得到，所以代入式(3)后，\par
log $\Omega$ = N log N - $\sum$ ni (log N - log f - $\epsilon$i\par
kBT)\par
i\par
= N log N - $\sum$ ni(log N - log f)\par
i\par
+ 1\par
kBT $\sum$ ni$\epsilon$i\par
i\par
= N log f + E\par
kBT\par
$\Leftrightarrow$ kB log $\Omega$ = kBN log f + E\par
T\par
(5)\par
就成了这样。清爽了。 」\par
式子清爽了， 但意思还不明白。 我等着制作人读了式子之后提出那个问题。 在等待\par
的这段时间里就是休息。\par
5\par
log W使用了斯特林公式。取log应该是因为……拉格朗日乘数法的关系。不过，因为最后式子出来\par
时是带log的形式，所以也可能是我看漏了什么\par

% --- PDF page 70 ---
69\par
能起到时钟指针走的的声音。 此外还能起到什么的话……现在也没有耳鸣， 大概就\par
只有从远处传来的喧嚣声了。啊，传来了欢呼声。是有谁在开演唱会 ?\par
「……式(5)是怎么回事啊?」\par
「制作人，觉有热力学的知识 ?」\par
「……觉怎么会觉得我有呢?」\par
「这样啊，制作人是1 年级生，所以不知道啊。 」\par
「不，都说了我是文科生，不修热力学的啊。 」\par
「…诶?不学热力学的话，怎么……理解世间的热现象呢?」\par
「\par
泽同学，世上的人们……是不理解，热现象的?对\par
泽同学来说，这个世界或许\par
很困难吧……」\par
是这样啊。这么一说，或许真是如此。 我想起前子子作为参，资料看的一个综艺节\par
目里，有个东大的人在空调问题上搞错了。我当时还以为大概是节目效果，但说不定，\par
他其实真的只是弄错了而已。\par
「…，那，我简单说一下热力学的基本概念。 『熵』这个词，觉起说过的，对吧。符\par
号是S， 就是这样一个物理量。 ……所谓物理量呢， 就是像温度T啦， 体积V啦， 那种……\par
被认为物理上存在也不会盾。的、一定的量。 要点是，系统的熵在什么都不做时通常是\par
会增加的， 达到一定。后就会取极大。而不再变的。 这和封闭系统中温度不再变化是一\par
回事，只是说平衡的时的特性如此罢了。然后，实际上式(5)的右边可以归结为S。\par
用系统的内能U （系统内部的能量结和，所以这里和E一『） ，将亥姆霍兹自由能\par
6\par
A\par
定义为A = U - TS的话\par
7\par
……利用热机的微小功dW，在热力学里有一系列广致\par
8\par
，\par
dA = dU - TdS - SdT\par
= dW - SdT (6)\par
就会出来这个式子。话说回来，根据省略掉的拉格朗日乘数法的结果，\par
d log f = - E\par
N d ( 1\par
kBT) - 1\par
NkBT dW (7)\par
6\par
就是这样一个物理量。是热机在恒温变化时能够取出的能量。反过来，TS 就是无法取出的部分\par
7\par
译者注：来文使用 F 代表亥姆霍兹自由能，国内普通高校教材一般使用 A，故译者做出如上修改\par
8\par
着一翻开哪本热力学书籍都会写，所以请参，任意一本书。对于只用书本的初学者来说，马塞马之\par
类的比好浅无易懂，广荐\par
↑译者注：马塞马(マセマ)为日本一是出版社，主要出版各类教材。对于国内读者，译者个人比好广\par
荐南京大学傅献彩教授主号的物理化学，对于物理化学各部分都讲解的比好透彻，对于，研备，人群也\par
是相当合适。\par

% --- PDF page 71 ---
70\par
可以得到 这个式子 ，于是为了能出现 dT 和d$\prime$W ，而且为了方一地处理系数，把\par
-A/kBT全微分一下的话，\par
d (- A\par
kBT) = - dA\par
kBT + A\par
kBT2 dT\par
= - 1\par
kBT (d$\prime$W - SdT) + 1\par
kBT2 (U - TS)dT\par
= U\par
kBT2 dT - dW\par
kBT\par
(8)\par
就得到这个了。比好式(7)和式(8)，\par
A = -NkBT log f (9)\par
就明白了，所以代入式(5)的话，\par
kB log $\Omega$ = kBN log f + E\par
T\par
$\Leftrightarrow$ kB log $\Omega$ = - A\par
T + E\par
T\par
therefore kBT log $\Omega$ = -(U - TS) + U\par
= TS\par
therefore S = kB log $\Omega$\par
(10)\par
会得出这么个结果。 这个式子被称为玻尔兹曼公式。在统计力学的入门阶段，大多\par
都是以广致出这个式子为目标， 所以我就绕道经过了一下。 ……但是， 式子变形太多了，\par
教千奈和佑芽的时候不合适啊。怎么法呢。 」\par
不看也知道制作人已经绪晕眼花了。不过，我还是想且且补充一下。\par
「这个广致的前提是设设各子子几乎互不干涉，但即使是互相干涉的子子，只要对\par
那些小系统同样地思，，同样的道理也能成立。详细内容会用所谓的正则系综啦、 巨正\par
则系综之类来表示……不过现在，不管。 」\par
「……觉能这么做，那可真是帮大忙了。 」\par
能明白熵与组合数的对数成正比会为此感到高兴， 大概……是在热力学中对熵有了\par
直。感受之后吧。那么，或许通过构建具体例子，把意义弄得更清楚一点比好好。\par
「有一个例子，形的类似水溶液，但可以更清地地公式化，我们来处理那个。固溶\par
体。制作人知道琥珀金 ?」\par
「不，没起说过。那是什么样的东西?」\par

% --- PDF page 72 ---
71\par
「金和的的天然合金。Au-Ag 系能以任意比例制成合金，也没有发生相变的浓度之\par
类。 这种就叫做无，固溶体。 在Ag 浓度低的区域， 可以看作是Ag 溶在了 Au 里， 反过\par
来也一样，因为是溶在固体里的东西，所以叫固溶体。当然，既然是溶进的的，通常想\par
来，当它以 Au 99\%对 Ag 1\%的浓度比进行溶解时，很难想象 Ag 会集中在中心，或者\par
浮到表面上来。也就是说，我们可以知道，Ag 很可能是分散存在于晶体中可以安置来\par
子的那些位置中的合适的地方， 对吧。 因此， 可以像才才做的那样， ，虑排布的组合数。\par
那么做的话，\par
$\Omega$ = (NAu + NAg)!\par
NAu! NAg!\par
therefore log $\Omega$ $\cong$ (NAu + NAg) log(NAu + NAg) - (NAu log NAu - NAu) - (NAg log NAg - NAg)\par
= (NAu + NAg) log(NAu + NAg) - NAu log NAu - NAg log NAg\par
(11)\par
……之类的就可以，虑了。 之所以取log， 就像才才说的， 是因为熵与log $\Omega$成正比。\par
然后，这个排布数，因为我们取的是不区分 Au 与 Au 彼此、Ag 与 Ag 彼此时的。，所\par
以在计 Au 或计 Ag 的情况下，它会变成 1，取 log 就是 0。因为有这么个方一的特性，\par
所以特别是在想，虑与 计品态的的差时，光准备这个就足够了。对于比率 xAu = NAu/\par
(NAu + NAg)，xAg = 1 - xAu，\par
log $\Omega$ = -(NAu + NAg)(xAu log xAu + xAg log xAg)\par
therefore kB log $\Omega$ = -kB(NAu + NAg)(xAu log xAu + xAg log xAg)\par
= $\Delta$S\par
(12)\par
像这样，就能知道熵S 的差分。光看这个式子，意思还不是很明白，所以在这里再\par
引入吉布斯自由能能G = U - TS + pV = F + pV = H - TS。\par
9\par
H = U + pV是焓。这里，\par
如果，虑到$\Delta$H = $\Delta$U + $\Delta$(pV)，第二项几乎为零\par
10\par
，所以$\Delta$H $\cong$ $\Delta$U = $\Delta$E。因为我们要比\par
好的是相同温度下的计金属和合金，所以$\Delta$T可以舍的，于是就有$\Delta$G = $\Delta$H - T$\Delta$S了。\par
那么，觉觉得$\Delta$H $\cong$ $\Delta$E该怎么求呢?」\par
为了在平板上写几个公式而稍微抬起的肩膀， 再次靠向了制作人。 因为平衡不太好，\par
身体稍微有点扭曲。另一侧的肩膀碰到了椅背，停了下来。正好形成由座面、靠背和制\par
作人这三点支撑的姿势。制作人不仅具备椅子的功能，还能让我开心、培育偶像，以及\par
整理各种信息，不过……\par
9\par
在亥姆霍兹自由能之上，加上了亥姆霍兹设设中因膨胀而损失的那部分能量。结之，就是设想因非\par
膨胀作用而能取出的能量，合金的过剩能量等就相当于这个。\par
10\par
固体和液体与气体相比，这一项大『视为零也没什么问题。\par

% --- PDF page 73 ---
72\par
「诶?………这里的 E 是各子子能量的结和对吧 ?那么，设法算出能级的平均之\par
类……。 」\par
「怎么做呢?」\par
「……我不知道啊。因为我连哪些方法是可行的都不清楚。 」\par
「…——60 分。所谓虽不中亦不远矣，大概就是这样。因为是合金系统，所以尽量\par
想用比率来做。 这样想来，素素而容易用的方法， 就是的数来子间键的能量。设Au-Au\par
间、Ag-Ag 间、Au-Ag 间的键能分别为eAuAu、eAgAg、eAuAg，设某个来子的配位数（键\par
数）为 z 的话，那么某个来子的能量可以写成\par
$\epsilon$ = zeAuAu\par
这样。键的起点和终点都是 Au 的模式的比例简单地用xAu $\cdot$ xAu就可以，所以对于\par
多算了的部分用 1/2 处理，\par
E = Nz (1\par
2 xAu\par
2 eAuAu + 1\par
2 xAg\par
2 eAgAg + xAuxAgeAuAg)\par
= 1\par
2 NzxAueAuAu + 1\par
2 NzxAgeAgAg + NzxAuxAg (eAuAg - eAuAu + eAgAg\par
2 )\par
(14)\par
$\Delta$E = Nz (eAuAg - eAuAu + eAgAg\par
2 ) xAuxAg = $\Omega$xAuxAg (15)\par
就得到了。这里的$\Omega$有时会被称为相互作用参数。\par
11\par
读一下式子就能明白，通常，\par
它是由用于合金的元素所定定的参数。$\Omega$ = 0的话， 就表示混合完全没影响，$\Omega$ > 0的话\par
就倾向于分离，$\Omega$ < 0的话就倾向于混合。最终，\par
$\Delta$G = $\Delta$H - T$\Delta$S\par
= $\Omega$xAuxAg + kBT(NAu + NAg)(xAu log xAu + xAg log xAg) (16)\par
这么一个式子就广致出来了。把这个式子以浓度为横轴， 适当定定常数，然后一边\par
变的$\Omega$一边画图， 其影响就一目了然了。\par
12\par
……制作人， 一直说话我有点口渴了。 呵呵。 」\par
「正好有发放的灾害储备用用水，是瓶装的，请喝那个吧。……我要站起来，请不\par
要再把体重压上来了。 」\par
我不要。\par
「……\par
泽同学?很险…的。 」\par
我不要。\par
11\par
与组合数$\Omega$不同，请注意区分\par
12\par
在 Desmos、WolframAlpha 等网站上在线画图很简一，广荐。\par

% --- PDF page 74 ---
73\par
……啊，制作人抓住我的双肩固定住了。着着『…』地一下，我大『回到了垂直态\par
的。稍微有点往反方向倾斜，不过这样一来……怎么说呢。\par
「我要放手了，拜托了。，\par
泽同学……。 」\par
「呵呵。怎么法呢。 」\par
直到肩绪被按住的感觉消失， 这段时间里连一次时钟的滴答声都没起到。 不小心反\par
应慢了，却维实住了平衡，于是我就这么眼睁睁地让制作人站起来了。算了，也好。喝\par
下制作人拿回来的水，小休息也结束了。\par
「………因为琥珀金是无，固溶体 （能以任意比例混合的二元系合金） ， 所以不适合\par
用来解的硫酸钠水溶液的溶解度。 会产生那种曲线 （指硫酸钠溶解度曲线） 的固溶体系，\par
叫做共晶系统，而这世上反而是这种更普遍。身边的例子的话，就是焊锡\par
13\par
之类的。无\par
铅焊锡\par
14\par
也是共晶系统。在这样的系统里，会以特定的浓度为界，某一方的成分会优，\par
单独析出。就如同冷却盐水时氯化钠会析出一样，例如，冷却锌过量的焊锡时，直到达\par
到共晶温度 183℃之前，锌会不断析出。等到缓慢冷却到 183℃时，那时的锌浓度就是\par
61.9\%了。 铅过量也是同理， 在那个浓度附近， 固溶体会稳定下来。 表现出这类行为的，\par
就是共晶系统。\par
这一点，从才才那个式子有时会形成 W 形图形这一点上可以一定程度地理解。，\par
虑 Pb-Sn 体系在 200°C 的情况，以 Sn 浓度 60\%附近为中心，只有那一带液化时吉布\par
斯自由能会更低。除那以外的地方，则是固体的能量更低。着着温度下降，T$\Delta$S项的贡\par
献逐渐减少，液相稳定的区域也着之缩小。最终到达共晶点时，稳定区域变为零。大『\par
类似的现象也发生在硫酸钠的溶解度曲线中， 在转变点附近，硫酸钠十水合物与无水物\par
之间出现稳定的区分，因此可以画出同样的相图。\par
15\par
……实际画一下吉布斯自由能曲线的话，大概就是这样。在图上，取公共切线。这\par
是因为会朝着极小。分解。那个平衡点并非出现在能量曲线上，而是产生于多个相\par
16\par
的\par
共同切线上。于是，正好在共晶点的地方，以纵轴为温度、横轴为浓度时，像这样。不\par
过，温度在 300K 下，关于kB(NAu + NAg)项设为 1 mol 时的kBN = R。R 是理想气体常\par
数。这个时候，在能量上，就处于一个任意的浓度下，固相也好液相也好都稳定的三相\par
点上了。\par
13\par
Pb-Sn 系\par
14\par
以 Sn-Ag 系为基上。应该，就算是在此之上的多元系也基本上是以此为基本。\par
15\par
因为画图很烦了，关于这一点请参， https://chemistrystory.com/sodium-sulphate-solubility-curve/\par
16\par
这次设想的是液相和固相\par

% --- PDF page 75 ---
74\par
图 1. 常温下的适当能量曲线\par
稍微留有点液相余地的情况下，温度在 500 K 时，是这样。\par
图 2. 将上图仅提高温度后的结果\par

% --- PDF page 76 ---
75\par
图正中间附近，被液相接触点所包围的区域是只有液相稳定的浓度区域。 才才作为\par
例子的琥珀金，因为相互作用参数$\Omega$小，所以不会变成像这张图黑线那样的 W 字形。\par
结果就是，虽然能产生固液混合相\par
17\par
，但基本不会变得像共晶系统那样，依浓度不同而\par
发生固液变化。\par
……所以， 说到最后的结论就是……温度在一定程度以上时， 无水物比十水合物更\par
稳定，就是这样。我觉得，其根据已经通过吉布斯自由能曲线，用更容易理解的方式解\par
的出来了，。。 」\par
「那么，这个……觉打算全都这样，解的给花海同学起 ?」\par
「…，就连我也知道这太勉行了。怎么法呢，呵呵。制作人觉得呢?」\par
「请不要问我啊。说起来，才才的讲解我自己都还没完全理解呢。 」\par
他眨了眨眼。没错，我带着实感明白了，来来这个表情就是那么回事。但是，\par
「才才的，很难懂 ?」\par
「用公平一点的说法来说……是的，要说难懂，确实是有点难懂。不过，那大概是\par
因为我本人缺乏这种物理素养， 而说明里也穿插了例子， 所以我想能懂的人会觉得很容\par
易懂吧。 」\par
来来如此。也就是说，才才那样对佑芽也行不通，是这么回事。就算设设她能起我\par
讲到最后，恐的中途就已经理不清逻辑关系了。\par
「知道了。那，我从绪再说明一遍，如果有理论不太通便的地方，就告诉我。 」\par
「……能不能让我，，稍微，休息一下……?」\par
我扭过绪，看了制作人的脸，就那样从背上朝他倒了下的。他没把我广开，所以，\par
我觉得挺好。\par
参考文献\par
[1]\par
,\par
,\par
:\par
(\par
5),\par
, (2005)\par
[2]\par
: 2\par
,\par
https://ist.ksc.kwansei.ac.jp/\textasciitilde{}nishitani/Lectures/2005/NewMaterialDesign/PhaseDi\par
agram.pdf, (2005)\par
17\par
实际的固溶体并非完全规则的固溶体，所以吉布斯能曲线在浓度 0.5 时并不对称，而是会倾向某一\par
方。因此，液相和固相取极小。的浓度会改变，从而能够画出共同切线。\par

% --- PDF page 77 ---
\noindent\textit{[第 77 页没有可抽取文字层。]}\par

% --- PDF page 78 ---
77\par
请\par
泽广教我李群与李代数\par
献给不顾专业领域不同而着我讨论的友人 H.O.。\par
登场人物介绍\par
广：\par
泽广。初星学园高中部偶像科 1 年 2 组，15 岁。来理论物理学者。离开学\par
界成为偶像。\par
P：广的制作人。就读于初星学园制作人科，专科大学一年级生。稍懂线性代数。\par
晶叶：池袋晶叶。14 岁。机器人工程学者。\par
数学符号\par
对本书使用的数学符号进行汇结。在正文中出现而不明白时，请看此处。\par
下表是本书使用的字母。来为大写的表示矩子，为粗体的表示为向量。\par
名称 集合 元素\par
体 K  k, l, s, t\par
群 G, H  a, b, c, d,单位元e\par
李代数 g, h  x, y, z, p, q, r\par
向量 V  u, v, w\par
整数 Z  m, n, i, j\par
$\cdot$点ḟ(t)：时间函数的时间微分。即：\par
ḟ(t) = d\par
dt f(t)\par
1 游乐园\par
今天和\par
泽同学来到游乐园玩。\par
广 「写着什么。 『回答谜题获得奖品』 ……机会难得，做做看吧。 」\par
P 「明白了。 」\par

% --- PDF page 79 ---
78\par
看来\par
泽同学似乎对谜题产生了兴趣。\par
P 「不过真意外。我还以为\par
泽同学会觉得这种东西太简单而无聊呢。 」\par
广 「没有的事。历史啦艺术啦，不知道的事情更多。 」\par
2 群\par
广 「问题 1。」\par
P 「哪个哪个………?」\par
广 「答案是……」\par
P 「请等一下。说到底，群是什么啊?」\par
泽同学看来立刻明白了答案，但我连问题的意思都不懂。\par
广 「不是教科书的定义， 」\par
P 「来来如此……?」\par
广 「那么，我们来，虑整数与差。封闭性很明确吧。接下来试代入恰当的。，，\par
虑是具成立。比如，\par
(2 - 3) - 1 = 2 - (3 - 1)\par
下列集合与运算中，构成群的是哪个。\par
1. 整数与差\par
2. 自然数与和\par
3. 有理数与积\par
4. 绝对。为 1 的复数与积\par
问题 1\par
群是满足以下性质的集合G与二元运算$\cdot$的代数结构(G,$\cdot$)。如果a, b, c是集合\par
G 的元素，则存在：\par
封闭性：a $\cdot$ b也是集合G的元素\par
结合律：(a $\cdot$ b) $\cdot$ c = a $\cdot$ (b $\cdot$ c)\par
单位元：存在集合G的元素e，满足a $\cdot$ e = e $\cdot$ a = a\par
逆元：存在G的元素a-1，满足a $\cdot$ a-1 = a-1 $\cdot$ a = e\par
群的定义\par

% --- PDF page 80 ---
79\par
正确 ?」\par
P 「………左边是-2，右边是0，所以不成立啊。 」\par
广 「那下一个。自然数与和。封闭性和结合律明确，单位元是 0。那么，2 的逆\par
元是?」\par
P 「-2。 这不是自然数， 所以不对。 那么有理数与积。 封闭性和结合律明确， 单\par
位元是 1 吧。2 的逆元是 1/2，看来逆元也存在啊。 」\par
广 「但是，0 的逆元不存在。0 也是有理数。 」\par
P 「来来如此。那最后是绝对。为 1 的复数与积 ……」\par
广 「从欧拉公式， 复数可以用绝对。和辐角表示。 不绝对。为1， 辐角分别为$\theta$,\par
$\omega$的话，\par
ei$\theta$ei$\omega$ = ei($\theta$+$\omega$)\par
因此封闭。 」\par
P 「结合律很明确，单位元和逆元也与有理数一样啊。0 的绝对。不是 1，所以\par
这就是群吧。 」\par
广 「群的运算也称为（群的）乘法。不过请注意，依具体情况加法或其他运算也\par
会成为群的乘法。 」\par
3 群的作用\par
P 「群与具的判别方法我懂了，但终究想要做什么，太抽象了，不太明白啊。 」\par
广 「那试着思，一个新的视角。 」\par
P 「什么意思?」\par
广 「映射是函数的一般化，不过替换成函数来理解也无妨。函数有，入和，出。\par
合成函数是将函数的，出连接到另一个函数。 」\par
群即映射，群的运算表示映射的合成。 此时，由结合律能以简洁形式求得合成\par
映射。\par

% --- PDF page 81 ---
80\par
P 「来来如此。可是，那一点和群有关系 ?」\par
广 「合成函数并不那么简单。例如，海伦公式一般写作\par
\{ s = a + b + c\par
2\par
S = $\sqrt{}$s(s - a)(s - b)(s - c)\par
但将其合成而用\par
S = 1\par
4 $\sqrt{}$(a + b + c)(-a + b + c)(a - b + c)(a + b - c)\par
表示的情况基本没有。 」\par
4 李群\par
广 「问题 2。」\par
P 「那个啊。 」\par
广 「答案是……」\par
P 「所以说请等一下。我连选项都不明白。 」\par
这次\par
泽同学也似乎立刻明白了答案， 但我对问题的意思完全不懂。 选项里排列着\par
见都没见过的字母。\par
广 「SO(n)是 n 维特殊正交群。 所谓正交群， 是指由转转向量的矩子构成的群。 」\par
P 「矩子不都是转转向量的 ?」\par
下列群中，不是李群的是哪个。\par
1. SO(n)\par
2. SE(n)\par
3. SU(n)\par
4. GL(n, Z)\par
问题 2\par

% --- PDF page 82 ---
81\par
广 「确实如此。 但是正交群在转转多个向量时， 也会持实它们之间的形态不变地\par
转转。如果只是把人偶在空中转一下就压扁了的话就不好了。」\par
P 「来来如此。那么，具体是什么样的形式呢?」\par
广 「那是……。制作人还记得内积表示什么 ?」\par
P 「内积 ……。记得是各自绝对。和夹角余弦的乘积，对吧?」\par
广 「正是。所以，当存在两个n 维实向量v, w和n $\times$ n的实矩子A时，至少希望内\par
积在作用前后一『。即希望\par
v $\cdot$ w = (Av) $\cdot$ (Aw)\par
成立。将此改写为矩子乘法的话，\par
vTw = vTATAw\par
所以，这仅当\par
ATA = E\par
时成立。如果v = w，才才的式子就是持绝对。的变换。内积与绝对。相同则夹角\par
也相等， 可以持实形态不变地转转。 也就是说， 只是逆矩子等于转置矩子这样的矩子。 」\par
P 「来来如此。那么，这就是SO(n) ?」\par
广 「不对。仅凭这个条件，无法排除像镜像那样翻转的转转。因此，这只被称为\par
n 维正交群O(n)。便一，正交群的行列式必为$\pm$1。知道为什么 ?」\par
P 「………因为转置也不会改变行列式，对吧?」\par

% --- PDF page 83 ---
82\par
广 「对。 而且， 是具翻转可由行列式的正负得知， 行列式恰为+1的矩子就是 n 维\par
特殊正交群SO(n)。不向量是复向量，内积则为\par
v $\cdot$ w = vTw\par
如此将共轭转置汇结以†表示的话，\par
A†A = E\par
这样的复矩子就是转转复向量的矩子。 这称为酉矩子， 其构成的群为n 维酉群U(n)。\par
同样，行列式恰为+1的矩子是 n 维特殊酉群SU(n)。」\par
P 「那么，SE(n)是什么?」\par
广 「SE(n)是特殊欧几里得群。首，，矩子能表示转转，却无法表示平行移的。\par
所以，是使其也能表示平行移的的矩子。具体而言，当存在特殊正交矩子A和\par
两个实向量v, b时，\par
[\par
A b\par
0 1\par
] [\par
v\par
1\par
] = [\par
Av + b\par
1\par
]\par
通过如此，能对v施加转转和平行移的。 」\par
P 「矩子里面还进了矩子和向量呢。 」\par
广 「如果v是 2 维向量，具体表示成分的话，则有\par
[\par
A b\par
0 1\par
] = [\par
a11 a12 b1\par
a21 a22 b2\par
0 0 1\par
]\par
这样的结果。」\par
P 「为什么矩子下方附有 0 或 1?像这样\par
[A b] [\par
v\par
1\par
] = [Av + b]\par
不也可以 ?」\par
广 「如果不是方子就无法计算逆矩子。没有逆元就不成群。 实际上，已知这个矩\par
子的逆矩子可以简单计算，就像这样\par
[\par
A b\par
0 1\par
]\par
-1\par
= [AT -ATb\par
0 1\par
]\par
乘上来矩子能切实得到单位矩子。 」\par

% --- PDF page 84 ---
83\par
P 「那，GL(n, Z)是什么?」\par
广 「这是 n 维整数一般线性群，由矩子元素全为整数且行列式为$\pm$1的矩子构成\par
的群。这个群的逆矩子也还是整数矩子，且关于矩子之间的积也封闭。 」\par
P 「来来如此。选项我懂了，但问题里的李群是什么?」\par
广 「李群是什么虽然很难解的， 但它是既是群又是流形的数学对象。 流形的说明\par
也很难，但例如，虑地应。地应虽然是三维的，但表面是二维的。流形是虽然\par
整体上看是弯曲的，但在局部上看是直的数学对象。」\par
P 「为什么要，虑这种东西?」\par
广 「那是为了忽略约束。例如地应虽然是三维的，但从某个地点前进时， 通常不\par
会向正下方或正上方走。 像这样在流形上移的的东西会有，制。 将这些移的到\par
可以忽略这些约束的空间是目的之一。」\par
P 「来来如此。那么，问题的答案是什么?」\par
广 「流形要求是连续的。这是因为要在表面上，虑微分。整数是离散的，所以\par
GL(n, Z)不是李群。便一一提，问题 1 里出现的『绝对。为 1 的复数与积』也\par
是李群。 」\par
5 李代数\par
广 「看。问题 3。」\par
P 「差不多注意到了，结感觉有些奇怪……」\par

% --- PDF page 85 ---
84\par
P 「又是李群……不对，是李代数 ?而且，字体好特别啊。 」\par
广 「这是称为 Fraktur 的字体。李代数是李群的单位元处的切向量空间。 」\par
P 「切向量空间是?」\par
广 「切向量空间是收集了切向量的空间。 首，， ，虑在李群上活的的时间函数。 」\par
P 「???」\par
广 「那么，试想一下道口。 」\par
3 维特殊欧几里得李代数se(3)是什么结构?\par
1.\par
[\par
0 x1 x2 y1\par
x1 0 x3 y2\par
x2 x3 0 y3\par
0 0 0 0 ]\par
2.\par
[\par
1 x1 x2 y1\par
x1 1 x3 y2\par
x2 x3 1 y3\par
0 0 0 1 ]\par
3.\par
[\par
0 x1 x2 y1\par
-x1 0 x3 y2\par
-x2 -x3 0 y3\par
0 0 0 0 ]\par
4.\par
[\par
1 x1 x2 y1\par
-x1 1 x3 y2\par
-x2 -x3 1 y3\par
0 0 0 1 ]\par
问题 3\par

% --- PDF page 86 ---
85\par
广 「道口处， 列车与汽车或行人以不同速度通过铁路和公路这不同的轨道。 但是，\par
两种轨道都经过地应上， 且在道口这一个地点相交。像这样，，虑经过李群上\par
的时间函数。 此时， 将某地点a处时间函数的微分， 也就是速度作为速度向量。\par
这就是切向量。因为时间函数只要经过地点 a 就行，所以存在无，个，切向\par
量也存在无，个，所有切向量构成的空间称为切向量空间。 」\par
P 「……。 」\par
广 「在，虑欧几里得群之前，，，虑正交群。正交群的李代数，由该李群的单位\par
元 （单位矩子E） 处的切向量空间表示。 具体地， ，虑正交群上的时间函数A(t)。\par
此处， 时间函数A(t)带有在时刻t = t0时经过单位矩子E的制约， 即A(t0) = E。\par
此时，时间函数A(t) 在时刻t = t0 的时间微分就是李代数的元素。将这李代数\par
的元素称为X吧。也就是说，Ȧ(t) = X。首，由正交群的性质，\par
ATA = E\par
这里对两边时间微分，\par
ȦTA + ATȦ = 0\par
此处令t = t0，由前提A(t0) = E，Ȧ(t) = X，\par
Ẋ TE + ETẊ = 0\par
X = -XT\par
就得到结果了。 」\par
P 「来来如此，但这具体是什么形式呢?」\par
广 「这称为反对称矩子或交错矩子，\par
X =\par
[\par
0 x12 x13 $\cdots$ x1n\par
-x12 0 x23 $\cdots$ x2n\par
-x13 -x23 0 $\cdots$ x3n\par
$\vdots$ $\vdots$ $\vdots$ $\ddots$ $\vdots$\par
-x1n -x2n -x3n $\cdots$ 0 ]\par
是这样的形式。 」\par
P 「像是对称矩子， 但对角成分为0， 且反侧符号相反啊。 可是， 明明是切向量，\par
却是矩子 ?」\par
广 「向量只要是在加法和标量乘法下封闭的集合就行。例如，即一是矩子，像\par

% --- PDF page 87 ---
86\par
[\par
x11 x12\par
x21 x22\par
] → [\par
x11\par
x21\par
x12\par
x22\par
]\par
那样变形后也会变成向量。特别地，像右侧那样把数字排成一列的向量，称为数。\par
向量。 」\par
广 「着眼于问题来思，的话，，虑欧几里得群的其他元素，因为 0 和 1 都是常\par
数，所以微分后变成 0。平移向量没有约束，所以即使微分也不会减少维度。\par
这并不意味着在李群和李代数中平移部分是一『的，而只是变量的约束不变。\par
所以答案是……」\par
P 「是 3 啊。李代数是向量我懂了，但和李群是什么关系呢?」\par
广 「将李代数的元素沿李群卷绕，就成为李群的元素。回想地应与地图的例子，\par
地应是李群， 地图是李代数。这种情况下虽然不知道单位元在哪里，但将那里\par
对齐，把比例尺 100\%的地图卷在地应上，那么地应上的点与地图上的点应该\par
一『。」\par
6 变换映射\par
P 「…——但还是觉得不太明白啊。 」\par
广 「那么，一边看问题 4 一边说明。 」\par
广 「从李代数向对应李群的变换，是将李代数的元素卷绕到李群上的操作。\par
lim\par
m→$\infty$\par
(E + X\par
m)\par
m\par
就是这样的式子。 」\par
关于作为李代数元素的矩子X，向李群的变换映射是以下哪个?\par
1. $\sum$ 1\par
n! Xn$\infty$\par
n=0\par
2. $\sum$ (-1)n+1\par
n! (X - E)n$\infty$\par
n=1\par
3. $\sum$ (X-1)n$\infty$\par
n=0\par
4. $\sum$ 1\par
n! (X - nE)n$\infty$\par
n=0\par
问题 4\par

% --- PDF page 88 ---
87\par
P 「什么意思?」\par
广 「…，比方说，设设有像下面这样的李群和李代数组。 」\par
广 「看最左边，设李群的元素A 与李代数的元素X 对应。现在，虑的是从李代数\par
向李群的变换，所以X已知而A未知。接着看正中间，这里，比方说不知道$\sqrt{}$A\par
将其平方就能知道A。」\par
P 「但是，连A都不知道，不可能知道$\sqrt{}$A吧?」\par
广 「对。但能知道近似的。。看图右。那是 X 的一半的点。以算式而言就是E +\par
X/2。」\par
P 「为什么要加E?」\par
广 「因为李代数是李群单位元处的切向量空间，所以偏移了那部分。 」\par
P 「但能进行得那么便利 ?」\par
广 「的确， 即使分割成2 个也不会那么接近。 但如果无，分割下的， 。就会接近。\par
，虑到不$\theta$足够小则存在tan $\theta$ $\doteq$ $\theta$，不觉得很妥当 ?」\par
P 「也许是吧。但是，问题选项里并没有啊。 」\par
广 「这个式子收敛比好慢。 所以一般使用别的式子。 首，， 将它用二项式定理展\par
开的话，\par
lim\par
n→$\infty$ Cm 0Em + Cm 1Em-1 X\par
m + Cm 2Em-2 (X\par
m)\par
2\par
+ $\cdots$\par
= lim\par
m→$\infty$\par
$\sum$ Cm n\par
mn\par
$\infty$\par
n=0\par
Xn\par
关于所有n，Xn的系数存在\par
lim\par
m→$\infty$\par
Cm n\par
mn\par

% --- PDF page 89 ---
88\par
= lim\par
m→$\infty$\par
m(m - 1)(m - 2) $\cdots$ (m - n - 1)\par
n! mn\par
= 1\par
n! lim\par
m→$\infty$\par
mn + k1mn-1 + k2mn-2 + $\cdots$ + kn\par
mn\par
= 1\par
n! lim\par
m→$\infty$\par
\{1 + k1\par
m + k2\par
m2 + $\cdots$ + kn\par
mn\}\par
= 1\par
n!\par
成立。其中ki是mn-1的有，系数。 」\par
P 「极，与无，和可以交换 ?」\par
广 「不太好。但这次就忽略吧。最终答案是 1，\par
$\sum$ 1\par
n! Xn\par
$\infty$\par
n=0\par
。。」\par
P 「嘿——」\par
广 「便一，不觉得很像什么 ?」\par
P 「像指数函数的麦克劳林展开呢。 」\par
广 「对。 所以这些映射被称为指数映射， 写法如指数函数。 逆映射是对数映射。 」\par
7 李群的伴随表示\par
广 「是李群的伴着表示呢。 」\par
P 「伴着表示是什么?」\par
当存在李群某点a处的切向量xa时，将其变换为单位元处切向量xe的式子是哪\par
个。\par
1. xa $\cdot$ a $\cdot$ xa-1\par
2. xa-1 $\cdot$ a $\cdot$ xa\par
3. a $\cdot$ xa $\cdot$ a-1\par
4. a-1 $\cdot$ xa $\cdot$ a\par
问题 5\par

% --- PDF page 90 ---
89\par
广 「来本， 李代数是李群单位元e处的切向量空间。 所以， 当存在李群另一点a处\par
的切向量xa时， 要作为李代数讨论， 就必须将其变换为单位元处的切向量xe。」\par
广 「在李群里定义了群的运算， 所以首，将a与exp xe相乘。 最后乘上a-1返回来。 」\par
P 「但那样的话，\par
exp xe = a $\cdot$ exp xa $\cdot$ a-1\par
会变成这样。。 」\par
广 「…。将它具体展开看看。运算符就省略了，\par
a (e + xa + xa2\par
2 + $\cdots$ ) a-1\par
aea-1 + axaa-1 + 1\par
2 axa2a-1 + $\cdots$\par
这里因为ae = a，第 1 项是aea-1 = aa-1 = e。相反，因为a-1a = e，所以也能将\par
其夹到第 3 项以后的xa的幂乘之间。\par
e + axaa-1 + 1\par
2 axaa-1axaa-1 + $\cdots$\par
= e + axaa-1 + (axaa-1)2\par
2 + $\cdots$\par
= exp(axaa-1)\par
既然这被变换成了exp xe，那答案就是 3 的，\par
a $\cdot$ xa $\cdot$ a-1\par

% --- PDF page 91 ---
90\par
。。」\par
8 终\par
稍走片刻，有一位双马尾戴眼镜的女孩。\par
晶叶「呀。谜题玩得开心 ?为了寻找精通机器人工程的人，我出了李群和李代数\par
的题目。。 」\par
漂亮地全部正解的\par
泽同学，得到了奖品磨泥器， 高兴地回的了。似乎是要用来做\par
搭配野草天妇罗的萝卜泥。真是弄不太明白的一天。\par
参考文献\par
[1] Selig J M. Lie groups and lie algebras in robotics[M]//Computational\par
Noncommutative Algebra and Applications. Dordrecht: Springer Netherlands, 2004:\par
101-125.\par
[2] Sola J, Deray J, Atchuthan D. A micro lie theory for state estimation in robotics[J].\par
arXiv preprint arXiv:1812.01537, 2018.\par
作者\par
https://sites.google.com/view/noisy-daemon\par

% --- PDF page 92 ---
91\par
Random mode!!\par
让\par
泽广讲解随机性的故事\par

% --- PDF page 93 ---
\noindent\textit{[第 93 页没有可抽取文字层。]}\par

% --- PDF page 94 ---
93\par
1 computably random\par
可 计 算 随 机\par
广：呐，千奈。\par
千奈：怎么了，\par
泽同学?\par
广：像我们这样积攒体力与元气、在最后一回合打出高伤害的风格，虽然常被人说\par
容易受运气左右。\par
千奈：觉在说什么啊?\par
广：千奈，觉知道着机是怎么定义的 ?\par
千奈：诶?着机…… ?\par
广：对。其实，着机在数学上是可以定义的。。\par
千奈：………突然这么问，我很为难啊……\par
广：那给觉提示。试着想想，着机序列，会具有什么样的性质，或许不错。\par
千奈：着机序列……是呢，比如无法言测接下来会发生什么，怎么样?\par
广：…，不错。着眼于这个性质而定义的，就是被称为可计算着机的东西。\par
千奈：可、可计算着机……\par
广：那么，我来讲解一下可计算着机的定义。。\par
千奈：请、请讲得通俗易懂些好 ?\par
广：首，，试想一个如下的游戏。\par
广： 事，准备好一个由0 和 1 组成的无，序列， 玩是要在第n 回合猜中序列的第 n\par
项。\par
广：游戏开始时，玩是拥有初始资金 1。然后，每一回合都将手中资金分成两份，\par
一份押 0，另一份押 1。\par
千奈：……。\par
广：下注后，公布第 n 项，押中那一项的钱会翻倍后返还。。\par
千奈：也就是说，设设在第 n 回合手实 100 元，0 押 30 元、1 押 70 元的话……\par
广：序列的第 n 项，如果出现 0 就返还 60 元，出现 1 就返还 140 元。。\par
千奈：我明白了。也就是说，大量押注在言测会出的那一方就好了呢!\par
广：…，就是这么回事。\par
广：那么，试想把这个游戏无，次进行下的。\par
千奈：无，次!?\par
广：因为一开始准备了无，序列嘛。那么，当无，次进行这个游戏时，千奈能把资\par
金增加到无穷大 ?\par
千奈：…，事，准备的序列，是着机排列 0 和 1 的 ?\par
广：对。\par

% --- PDF page 95 ---
94\par
千奈：那不可能不是 ?着机排列的东西是无法言测的啊!\par
广：是呢。也就是说，可以这样定义：在这个游戏中，使玩是无法无，赚钱的那种\par
序列，正是着机序列。\par
千奈：来来……如此……?\par
广：那么，把才才的内容变成数学上严格的定义吧。\par
广：玩是事，准备好策略，并按照该策略进行游戏时的资金变化，用从 01 有，列\par
到非负整数的函数m(x)来表示。\par
广：这时，m满足以下条件。\par
m(x) = (m(x0) + m(x1))/2\par
广：具有这个性质的函数称为鞅。各个鞅函数就对应于玩是的各个策略。\par
广：也就是说，这个游戏中玩是无法无，赚钱的序列，就是玩是策略所对应的鞅函\par
数不发散的序列。\par
广：特别地，对于所有可计算鞅函数都不发散的序列，就称为可计算着机序列。。\par
千奈：可计算，是什么意思?\par
广：要讲严谨的定义会累， 所以粗略来说， 能用某种号程语言， 比如C 语言或 python\par
等实现的函数，就称为可计算函数。\par
千奈：………\par
广：怎么样，懂了?\par
千奈：大『……?\par
广：是 ，太好了。着机性的定义还能再讲两种……\par
千奈：诶\par
广：今天累了，剩下的明天讲吧。\par
2 Kolmogorov complexity\par
柯 尔 莫 哥 洛 夫 复 杂 度\par
广：那么，接着讲昨天的后续。。\par
千奈：是着机性的第二种定义，对吧?\par
广：对。为此，要着眼于着机性的不同性质。\par
千奈：记得昨天用的是「无法言测」这个性质，对吧?\par
广：是呢。今天要用的性质是「无法压缩」。。\par
千奈：压缩……这和着机有关系 ?\par
广：这个看具体例子更快吧。\par
广：例如，虑如下序列。\par
00000000001111111111\par

% --- PDF page 96 ---
95\par
广：这是一个不太着机的例子，但它能用更短的方式表达 ?自然语言也行。。\par
千奈：………「10 个 0 之后是 10 个 1」，这样可以 ?\par
广：不错，那下一个例子。\par
10111011111010011100\par
广：这个怎么样?\par
千奈：这个……似乎很难用简短的话来表达啊……\par
广：…。也就是说，着机性越高的序列，就越难压缩，就是这么回事。\par
广：那么，这个也从数学上来严格定义。。\par
广：将从 01 有，序列到 01 有，序列的可计算函数D，称为解压器。\par
广：有，序列x关于解压器D的复杂度，定义为CD(x) = min\{l(y): D(y) = x\}。\par
广： 说得容易懂一点， 复杂度就是指长度为l(x)的有，序列x， 能通过解压器D被压\par
缩到长度CD(x)。\par
千奈：也就是说，复杂度越大……\par
广：就越难压缩，也就是着机性越高，就是这么回事。\par
广：更严谨地说，还会取压缩率最高的最优解压器等等，但讲到那里会累，所以今\par
天且且割爱。\par
广：再说，复杂度虽然是针对有，序列定义的，但其实讨论有，序列是具着机并不\par
太有益。\par
千奈：这是什么意思?\par
广：比如，长度为 3 的 01 序列有 000、001、010、011、100、101、110、111 这\par
8 个，但哪个是着机序列，哪个不是，这不太清楚对吧。\par
千奈：的确……\par
广：就那样， 要把复杂度用于着机性的严格定义还需要再的一点脑筋， 不过那方面\par
改日有机会再说吧。\par
广：那么，明天讲着机性的第三种定义，马丁-洛夫着机。。\par
3 Martin-Löf random\par
马丁 - 洛夫随机\par
广：那么是着机性的第三种定义。使用的性质是「只具有平凡的特征」。。\par
千奈：这我还是完全不懂啊……\par
广：例如，具有「1 的后面必定教着 0」或者「1 的比例比 0 的比例多」之类特征\par
的 01 序列，觉觉得是着机序列 ?\par
千奈：那个……似乎不太着机啊……\par
广：对。反过来讲，所谓着机序列，就可以说是完全不具有这类特殊特征的序列。\par

% --- PDF page 97 ---
96\par
千奈：来来如此……!\par
广：那么，为了利用这个性质来定义着机性，需要定义特征的一般性，或者说特殊\par
性。\par
广：这可以定义为，具有相同特征的元素的集合的大小。\par
千奈：……?\par
广：就是说，具有相同特征的元素越少，该特征就越特殊，越多，就越是普遍的特\par
征。\par
广：集合的大小就是测度。所以，「序列具有更特殊的特征」这句话，可以换言为\par
「序列包含在测度更小的集合中」。\par
千奈：速度……?\par
1\par
广：是 measure，不是 speed。千奈在初中没学过 ?\par
千奈：没学过。!我觉得\par
泽同学也不是在初中学的吧!\par
广：也许是那样吧。\par
广：重新来，定义马丁-洛夫着机。。\par
广：首，，「序列不是着机的」，就是「序列具有特殊特征」，而这即是说「序列属\par
于测度小的集合」。\par
广：将此铭记在心，来定义如下的马丁-洛夫检验。\par
广：将马丁-洛夫检验，定义为关于可计算测度$\mu$，满足$\mu$(Ui) $\leq$ 2-i的集合的一『递\par
归可枚举序列Ui。\par
广：也就是说，是测度越来越小的集合序列。\par
广：然后，对于任何马丁-洛夫检验Ui，都不属于$\cap$i Ui的 01 序列，才是 「不包含在\par
测度小的集合中的 01 序列」，也就是「着机的 01 序列」。这就是马丁-洛夫着机的定义\par
。。\par
广：怎么样，懂了?\par
千奈：完全不懂啊～!\par
广：呵呵，不如意呢。\par
千奈：不过，看到\par
泽同学讲得那么开心，我很高兴。!\par
广：…………呵呵。\par
4 Equivalency\par
等价性\par
千奈：话说\par
泽同学，我有个问题。\par
1\par
此处是同音。\par
与\par
发音都是\par
。\par

% --- PDF page 98 ---
97\par
广：好啊，我在可能的范围内回答。\par
千奈： 到此为止，我学到了三种着机性的定义……结果， 哪个定义才是那个……最\par
好的定义啊?\par
广： 非常好的问题。 其实， 只要凑齐条件， 这些分别定义的三种定义就会彼此等价。\par
广：首，，可计算着机与马丁-洛夫着机的关系。\par
广：可将所有可计算的鞅函数中不发散的序列定义为可计算着机，但是，不定义为\par
所有从下方可计算的鞅函数中不发散的序列，则这就成为马丁-洛夫着机。\par
广：其次是复杂度与马丁-洛夫着机的关系。\par
广： 这里不用柯尔莫哥洛夫复杂度C(x)， 而使用由无前缀解压器定义的前缀复杂度\par
K(x)。\par
广：此时，以下定理成立。\par
定理（莱文-施诺尔） 。\par
$\alpha$ $\in$ 2$\omega$为 ML-着机，\par
当$\exists$c，$\forall$x $\preceq$ $\alpha$，l(x) - K(x) $\leq$ c成立，则两者等价。\par
千奈：……这是什么啊?\par
广： 粗略来说， 这是个定理： 马丁-洛夫着机序列的复杂度会变大， 反过来， 复杂度\par
足够大的序列会成为马丁-洛夫着机。\par
广：以上就说明了，三种着机性的定义，在适当的条件下彼此等价。\par
广：这样，关于计算论着机性理论，最基上的基上部分就讲完了。照这个样子学下\par
的，成为出色的数学是吧。\par
千奈：才不会啦!?\par
广：从这里开始，是给在大学里学过数学，特别是分析学的人的小赠品。。\par
广： 敏锐的读者， 看到马丁-洛夫着机的定义， 或许已经想到， 着机序列就是除的部\par
分包含于零测集中的序列，即「几乎处处序列」。\par
广：其实，以这种形式定义的着机性概念，与分析学中常用的 「几乎处处」概念亲\par
和度很高。\par
广：例如，我想觉们起过「几乎处处可微的函数」，关于此，以下定理成立。\par

% --- PDF page 99 ---
98\par
定理。以下等价。\par
$\bullet$ z $\in$ [0,1]是可计算着机实数。\par
$\bullet$ 任意可计算非递减函数f: [0,1] → R在z处可微。\par
广：啊啊，所谓着机实数，就是在 01 着机序列前加上 0.，看作二进制小数展开的\par
东西。。\par
广：像这样，可以思，把着机性理论和其他领域连接起来。\par
广：那么，这次真的结束了。Thank you for reading.\par

% --- PDF page 100 ---
99\par
两位广与巴拿赫-塔斯基定理\par
导入\par
广： 佑芽，早上好。\par
佑芽： 广酱!早上好!!!!\par
广： 佑芽今天也从早上起就精如饱满。我已疲劳困惫。\par
（门的声音）\par
佑芽： 啊，千奈酱也来了!\par
广： 千奈也早上好。\par
千奈： 两位都已经到齐了啊。 比起这个， 请起我说!今天下了车到这儿来的路\par
上，我碰巧和\par
泽同学一道了……可是……\par
千奈： …………????\par
千奈：\par
泽同学…………\par
千奈： 为什么有两个人啊?\par
广 1： 不知怎么，一醒来就变多了。\par
广 2： 真是不可思议呢。\par
千奈： 诶，诶诶!?\par
广 1： 对了。佑芽，闭上眼睛。\par
佑芽： ?\par
（两位广来来来的）\par
广 1： 小测验。我们左右交换了 ?\par
佑芽： ?没交换吧?\par
广 2： 佑芽好厉害。为什么能知道?\par
佑芽： 穿的衣服不是不一样嘛。\par
广 1： 啊，对。。\par
广 1： 那脱掉吧。\par
千奈： 不、不能脱啊!?\par
佑芽： 这是规定啊\par
千奈： （规定……?）\par
这种事，别管， 必须快点找恢复来态的法法啊!\par
泽同学， 为什么变成\par
了两个，心里有什么绪… ?\par

% --- PDF page 101 ---
100\par
广 2： 变多的理由……说到东西变多……\par
广 2： 现在开始讲解巴拿赫-塔斯基定理。\par
千奈： 太冷静了啊～!!!!\par
关于体积的思，\par
佑芽： 巴拿赫-塔斯基定理是啥?\par
广：\par
1\par
设设有 1 个应。 粗略来说， 这个定理是说， 那个应可以被分解成有，\par
个碎片，通过转转和平移（以及镜像） ，能重新拼成和来来应同样大\par
小的 2 个应。\par
佑芽： 不可能有那种事!\par
广： 被具定了。\par
千奈： 但是， 正如花海同学所说， 那种事不能法到就是革命了!能量守恒定律\par
跑哪里的了啊?\par
广： 我虽说是分解成有，个碎片，但并非物理上可行的分解。 还有， 现在\par
，虑的是内部塞得严严实实的（数学上的）应，而现实中的「应体」\par
（比如应）是「基本子子集合构成的稀疏的核」与电子构成的稀疏的\par
来子的集合，和数学上的应情况完全不同。\par
千奈： 数学有时和日常生活情况不同啊……!\par
千奈： （那不就和\par
泽同学变成两个人的事没关系 ……）\par
佑芽： 可是， 教科书上登的毕达哥拉斯定理\par
2\par
的证明之类， 会用到即使转转面\par
积也不变这件事吧?\par
3\par
现在的情况也一样，应该成立（来来应的体积）\par
=（碎片体积的合计）=（将碎片组合后的结果的体积） ，不可能让应\par
增加才对!\par
广： 不对。因为分解后的各碎片并不存在「体积」 。\par
千奈： 是说体积为 0  ?\par
广： 不是 0，而是无法赋。。\par
佑芽： 有那种事 ?\par
1\par
以下将两位篠泽广视为同一人。\par
2\par
本稿采用人名以英字书写的惯例，可能看起来不熟悉，但这就是毕达哥拉斯定理（勾股定理）\par
↑译者注：这是作者的说法，我给各位都改回中文人名了。\par
3\par
中等教育不太重视定理的证明，完全可能有人从未见过证明。为这样的读者，附录中载有证明。附\par
录不在本刊，而在本稿末尾。\par

% --- PDF page 102 ---
101\par
广： 作为长度、面积、体积的一般化，我们来，虑测度。才才佑芽，虑了\par
体积的合计，但为了让测度能用作对「大小」的公式化，需要把「部\par
分合计的测度等于部分测度的合计」 作为条件\par
4\par
。 平时我们还会进一步\par
使用满足「平移和转转下不变」的（各维度的）勒贝格测度，但不对\par
所有集合都赋予勒贝格测度，就会产生奇怪的事\par
5\par
*5。\par
佑芽： 为什么?\par
广： 我来试着了一个无法规定勒贝格测度的R的子集吧。现在，在R上致\par
入等价关系\textasciitilde{}s. t.  a, b $\in$ R，a\textasciitilde{}b\par
def\par
$\Leftrightarrow$ a - b $\in$ Q。并设V为[0,1]/\textasciitilde{}的完\par
全代表系\par
6\par
。\par
千奈： 咒……咒文啊!\par
广： 不用集合论的语言来说明的话，就是，0 以上 1 以下的实数中，将像\par
$\sqrt{}$2 - 1和 $\sqrt{}$2 -\par
1\par
2 这样只错开了有理数的数视为同一物的结果， 就是\par
V。这是勒贝格不可测的，即无法赋予长度\par
7\par
。\par
佑芽： 是这样 ?\par
图 1：V与V$\infty$的示意图。\par
广： 试将V用从-1 到 1 的所有有理数平移并合并。 （下面将此记为V$\infty$。 ） 这\par
个结果会收敛在-1 以上 2 以下的范围内，对 ?\par
佑芽： 把V所在的范围 （[0,1]） 和平移幅度的范围 （[-1,1]） 的最小。与最大\par
。加起来，就行了呢!\par
广： 那么，它包含[0,1] ?\par
4\par
此处未提及，但对空集赋予 0 的测度也作为条件。\par
5\par
此事（即勒贝格不可测集合的存在性）不使用选择公理是无法证明的。\par
6\par
此V被称为维塔利集合。\par
7\par
选择公理的准确表述，以及（在设定选择公理下）从商集能取完全代表系的证明，载于附录。\par

% --- PDF page 103 ---
102\par
佑芽： ?\par
千奈： 是说，0 以上 1 以下的无论哪个实数，都包含在V$\infty$里 ?\par
广： …。 严密的证明省略\par
8\par
， 但可以认为， 是把仅因有理数偏差而被视为同\par
一的实数们，再一次平移恢复了来态。那么，V$\infty$的长度\par
9\par
就在 1 以上\par
3 以下对吧?\par
佑芽： 毕竟它包含[0,1]（长度 1） ，又被包含在[-1,2]（长度 3）里。\par
广： 这里， 不给V赋予长度 0，则V$\infty$的长度就成了 0； 不给V赋予正的长度，\par
则V$\infty$的长度就成了$\infty$\par
10\par
，所以，应赋予V的长度不存在。\par
佑芽： …～……虽然很别扭，但不得不承认……!\par
千奈： 可是\par
泽同学，应增加的故事里，，虑的是分解成有，个碎片， 像这\par
样收集起无，个时的问题，难道不是可以不，虑 ?\par
广： 千奈说的也有一定道理。 不区域和的测度等于测度和的性质只在有，\par
和时成立就好， 那么在1 维的情况下， 有种叫巴拿赫测度的东西存在\par
（是非平凡的测度\par
11\par
），它可以在平移下不变，并且能对所有数轴的子\par
集赋予测度。\par
佑芽： 那个叫巴拿赫测度的东西，1 维以外就没有 ?\par
广： 2 维的话有。所以，关于圆，就正如佑芽所说， （来来圆的巴拿赫测\par
度）= （碎片巴拿赫测度的合计）= （经平移、转转后的碎片巴拿赫测\par
度的合计）=（操作后结果的巴拿赫测度）成立，即使分解后转转、\par
平移再组合， 也不会从来来的圆增加测度，特别是，不会增加成和来\par
来的圆同样大小的 2 个圆。\par
佑芽： 既然 2 维有，那么 3 维呢?\par
广： 没有。\par
佑芽： 为什么啊～～????\par
广： 2 维和 3 维的区别，可能之后会聊到。\par
8\par
载以略证。任意[0,1]的元素，存在与其被视为等价的V的元素，其差。是包含于[-1,1]的有理数。\par
9\par
不是包含V$\infty$的区间的长度，而是V$\infty$本身的长度（赋予V$\infty$这个集合的测度） 。\par
10\par
这广理由用了区域和的测度等于测度的和的性质，以及勒贝格测度在平移下不变。\par
11\par
给所有区域都赋予 0 的测度的叫平凡测度。如果经操作而「增加」的东西其测度为 0，那么即使它\par
变成 2 个也不会致出盾。，所以这点很重要。\par

% --- PDF page 104 ---
103\par
将东西「变多」\par
千奈： 因为并不结能给分割出的碎片赋予体积， 所以『「转转和平移不会增加\par
体积」这一直觉不可用，这一点我懂了。但尽管如此，对做了点什么\par
东西就变多了这件事，还是完全没实感啊。\par
广： 作为这类话题的典型例子，来说说把自然数\par
12\par
变多吧。\par
佑芽： 把自然数变多?\par
广： 把自然数分成偶数和奇数。偶数除以 2，这就成了自然数全体。因为\par
只是把第 n 个偶数送到 n 而已。\par
图 2：把自然数变多的样子\par
千奈： 嘛，是……这样呢。\par
广： 同样，奇数减 1 再除以 2，这也成了自然数全体。因为这也只是把第\par
n 个奇数送到 n 而已。这样一来，把自然数分成2 部分再操作，来来\par
的自然数全体就成了 2 个对吧?\par
13\par
佑芽： 无法具定……\par
广： 这有时被表现为 「在定员无，的客满酒店，即使新来无，位客人，也\par
能接待所有客人」\par
14\par
，被称为希尔伯特无，旅馆悖论。\par
千奈： 没有不能住宿的人，真是美妙呢。\par
广： 接着，来看圆周（下面设为S）的掉 1 个点的东西（下面设为S∗） ，将\par
那 1 点补上吧。\par
佑芽： 已经能把东西变成 2 倍了，增加 1 个点这种事，不费吹灰之力!\par
广： 佑芽干劲十足呢。那么，觉知道怎么做 ?\par
佑芽： 不知道!\par
广： 那我来做。 首，， 在S∗上打点pn(cos n , sin n)。（n 是正整数） 。 这不会\par
与(0,0)或其他的pn点重合，没问题吧?\par
12\par
包含 0。\par
13\par
此处往后，话题将转向对某集合 X，使用从 X 到 X 的可逆双射将东西变多，但请留意，此例并不\par
属于该情况。 （各个操作作为映射，其定义域并未覆盖自然数全体。 ）\par
14\par
这里的「无，」是可数无，的意思。\par

% --- PDF page 105 ---
104\par
千奈： ?n=360 之类的时候，pn不会变成(0,0) ?\par
广： 不会，这个cos n里的 n 不是n°，而是n rad。\par
佑芽： 弧度?\par
广： $\pi$ rad = 180°。 比方p1就是(cos\par
180\par
$\pi$ ° , sin\par
180\par
$\pi$ °)。 在图上大概是这一带\par
（参照下图） 。\par
千奈：\par
180\par
$\pi$ °的话是个相当半吊子的角度， 所以确实不会和(0,0)重合， 或2 次\par
通过同一点的样子呢\par
15\par
。\par
广： 把这pn全部向便时针方向转转1 rad\par
16\par
。 也就是说， ，虑将p1移到(0,0)，\par
pn+1移到pn。这样一来，这就成了完整的圆周。\par
图 3：左边缺了(0,0)点的是S∗。在将那个缺的点补上。\par
佑芽： 也就是，像这样（参照图 3） ，用p1填最初的。，用p2填那里空出的\par
。……就这样无，重复下的对吧?但那样的话，任何时候都空着 1 个\par
。，不会变成完全的圆周不是 ?\par
广： 不不。pn有无，多个，用来填。的点可以无，取出，所以没问题。\par
佑芽： 诶～～????好别扭啊～……\par
千奈： 无，，好深奥啊……!\par
15\par
此证明省略。 （略证：用反证法。利用此类pn可广致出n为有理数。 ）\par
16\par
作为从S到S的映射看时，这是双射。\par

% --- PDF page 106 ---
105\par
操作的集合\par
广： 从这里开始， 我们想看看更接近正题的例子， 但在那之前， 我想把 「做\par
些操作让应增加」这个目标，用更合适的说法重新建立一下。\par
千奈： 更合适的说法?\par
广： 现在，设某集合X到X的映射全体的子集G，满足以下条件\par
17\par
。\par
$\bullet$ G包含恒等映射。\par
$\bullet$ G的各元素都是双射。\par
$\bullet$ G的元素g1, g1的合成映射g1 $\circ$ g2\par
18\par
是G的元素。\par
$\bullet$ G的各元素的逆映射是G的元素。\par
就称这为「G具有群的结构」吧\par
19\par
。\par
广： 而且，对于X 的子集A ，将其分割为A1, A2, $\cdots$ $\cdots$ , An 和对应的G 的元素\par
g1, g2, $\cdots$ $\cdots$ , gn，满足：\par
$\bullet$ g1(A1), g2(A2), $\cdots$ $\cdots$ , gn(An)互不相交\par
$\bullet$ g1(A1), g2(A2), $\cdots$ $\cdots$ , gn(An)的并集是B，\par
此时，我们就称A与B是 G-裁剪同余的。\par
佑芽： 没有具体性，完全不知道在说啥!\par
广： 那举例子。设从R2到R2的映射中，可用(x\par
y) $\mapsto$ (ax + b\par
cy + d)（a, b, c, d是\par
实数，|a| = |c| = 1） 表示的全体集合为H。 这个满足才才给出的4 个\par
条件，没问题吧?\par
千奈： 当(a, b, c, d) = (1,0,1,0)时，这是恒等映射呢。双射这点……\par
广： 展示各元素的逆映射包含于H 时也连带展示存在性， 所以待会再说。\par
佑芽： 关于映射的合成， 将2 个映射(x\par
y) $\mapsto$ (ax + b\par
cy + d) , (x\par
y) $\mapsto$ (a$\prime$x + b$\prime$\par
c$\prime$y + d$\prime$)合成，\par
就成了(x\par
y) $\mapsto$ (aa$\prime$x + ab$\prime$ + b\par
cc$\prime$y + cd$\prime$ + d)……\par
20\par
千奈： 因为|aa$\prime$| = |a||a$\prime$| = 1 ∗ 1 = 1 （同样的|cc$\prime$| = 1） ， 所以确实似乎包含\par
于H呢。\par
17\par
这是作用于X的群G这一想法的雏形。把它抽象化后就是群作用，但，虑到讨厌过度抽象议论的读\par
者很多，这次不涉及「真正的」群作用。\par
18\par
映射的合成一般用$\circ$作符号，但本稿有时将其省略。\par
19\par
与*17 同样，这也不是「真正的」群的定义。本稿下文将不加言告地使用「群」这一用语，但阅读\par
本稿时，无需深入思，此「群」与真正群的区别。\par
20\par
请注意，写在左侧的操作被后施行。\par

% --- PDF page 107 ---
106\par
佑芽： 对于 逆映射，设 ax + b = X, cy + d = Y 逆向解出的话， x =\par
1\par
a X -\par
b\par
a , y =\par
1\par
c Y -\par
d\par
c，所以是(x\par
y) $\mapsto$ (\par
1\par
a x -\par
b\par
a\par
1\par
c y -\par
d\par
c\par
)……\par
21\par
千奈： |1/a| = 1/|a| = 1（同样的|1/c| = 1） ，所以这也是H的元素呢。\par
图 4：A与B。这些是 H-裁剪同余。\par
广： …。接下来，想请觉们证明，上图所示的R2的子集A, B是 H-裁剪同\par
余。\par
佑芽： 也就是说，把A 分解成有，个碎片，然后将各个碎片用H 的元素移的\par
做成B就行了吧?加油喽～～～～!\par
21\par
需要注意a, c $\neq$ 0。\par

% --- PDF page 108 ---
107\par
图 5：A与B是 H-裁剪同余的样子\par
千奈： 2 个的公共部分好像可以持实来样呢\par
22\par
。\par
佑芽： 剩余部分确实是同余的，但……\par
广： 意识到 H 是x轴翻转、y轴翻转、平移的组合，应该就能明白了。\par
佑芽： 这个是沿x 轴翻转，然后移的(1\par
1) 后的东西……所以用( x + 1\par
-y + 1) 移的\par
的话就恰好嵌上了!!!!\par
广： G-裁剪同余这个词的意思，明白了 ?\par
千奈： …!\par
广： 想再定义一个用语。 称X为 G-复制可能， 是指存在X的裁剪X1, X2，使\par
得X1, X2两者都与X是 G-裁剪同余的。设R3到R3的映射中，不改变 2\par
点间距离的映射全体\par
23\par
为E(3)，这具有群的结构。那么，巴拿赫-塔斯\par
基定理的断言就可以重新表述为「单位应体B是 E(3)-复制可能的」。\par
佑芽： 回到复制可能与裁剪同余的定义，就是「单位应体B能被分成 2 个部\par
分，各个部分分别被分割成有，个碎片，用E(3)的元素移的后，2 个\par
都变成B自身」 。这就和最初广酱说的意思一样了吧!\par
22\par
请注意 H 包含恒等映射（『「持实来样」的映射） 。\par
23\par
即是平移、转转、镜像及这些的组合。\par

% --- PDF page 109 ---
108\par
广： 最后再举一个把东西变多的例子。设有某种操作$\sigma$, $\tau$ 和抵消它们的操\par
作$\sigma$-1, $\tau$-1的（0 个以上）\par
24\par
的组合构成的集合\par
25\par
为F。 设设这些操作除\par
了群的结构以外，不具有其他结构。这里，设集合存在F∗，其元素是\par
与F的各元素g对应的F到F的映射g $\circ$， 那么F是F∗-复制可能的。 我想\par
展示这个。\par
佑芽： ??????\par
广： 操作是可以合成的，但不将合成的左侧（后施行的）固定，合成就能\par
看作 「将操作映射到操作的操作」 。 比如， 以才才，虑的H来说，( x\par
-y) $\circ$\par
就是 「进行该操作之后， 再沿x轴镜像」 这一合成的操作。 比如将H的\par
操作(x\par
y) $\mapsto$ (x + 1\par
y - 2)移到(x\par
y) $\mapsto$ ( x + 1\par
-y + 2)。\par
千奈： 从操作到操作的操作，我明白了。\par
广： 以H为例，比方设$\sigma$ = ( x\par
-y) , $\tau$ = (-x\par
y )的话，就是$\sigma$2 = $\tau$2，但 「不具\par
有 群 的 结 构 以 外 的 结 构」 就 是 指 没 有 这 类 情 况 。F 是 禁 止 了\par
$\sigma$$\sigma$-1, $\sigma$-1$\sigma$, $\tau$$\tau$-1, $\tau$-1$\tau$（即忽略取消前一个操作的东西，而区分除群\par
结构所含关系之外的所有操作）的$\sigma$, $\tau$, $\sigma$-1, $\tau$-1 的序列。如果不好理\par
解， 只要把F的元素想成单计的$\sigma$, $\tau$, $\sigma$-1, $\tau$-1这些字母的序列， 而F∗的\par
元素则是在字符串左侧添加特定字符串的东西就行了。\par
佑芽： 那么，用这个怎么变多呢?\par
广： 首，，将F分为以$\tau$或$\tau$-1开绪的(A)和除此以外的(B)。要展示这 2 个\par
部分都和F全体是F∗-裁剪同余见下图。\par
图 6：将A1, A2, B1, B2可视化的图。\par
24\par
0 个操作的组合因为不能什么都不写，所以方一起见记为 e。\par
25\par
这作为映射，其定义域和。域为何无需在意。这是怠于抽象化、未，虑「真正的」群所致『的弊\par
端。\par

% --- PDF page 110 ---
109\par
千奈： 这就是F∗-复制可能的定义呢。\par
广： 将A 分割成以$\tau$ 开绪的(A1) 和以$\tau$-1 开绪的(A2) 。对A1 施以F∗ 的操作\par
$\tau$-1 $\circ$，就成了除的A2的F全体。\par
佑芽： $\tau$-1 $\circ$就是往左边加$\tau$-1吧。A1是以$\tau$开绪的， 所以这会被的掉呢!?因为\par
$\tau$-1$\tau$是被禁止的， 所以$\tau$的紧后面不会接$\tau$-1。 所以就成了除A2 （以$\tau$-1\par
开绪的）之外的东西?\par
广： 正是!\par
千奈： B这方该怎么法呢?\par
广： 这方也几乎一样。设B1为e、以及所有以$\sigma$开绪的东西、以及$\sigma$-1单独\par
连续着的东西的集合，B2为其剩余\par
26\par
。这里，对B2施以F∗的操作$\sigma$ $\circ$，\par
就成了除B1之外的F全体。\par
千奈： 和方才一样，B2 的要素全都以$\sigma$-1 开绪，所以除了$\sigma$-1 单独连续的部\par
分缺了以外，它似乎会和所有以$\sigma$开绪的东西的集合一『呢。\par
佑芽： 就是恰好嵌合，对吧!\par
广： 这样，就展示了A 和B 也都与F 是F∗ -裁剪同余，所以展示了F 是F∗ -复\par
制可能。\par
巴拿赫-塔斯基定理\par
广： 才才作为例子展示了F 是F∗ -复制可能，但其实，这就展示了巴拿赫-\par
塔斯基定理的关键部分。\par
千奈： 是这样的 ?\par
佑芽： F不是那莫名其妙、毫无具体性的操作的群嘛!\par
千奈： 完全搞不懂怎么就能和平移、转转连上!\par
广： …。所以，从现在开始让它具有具体性。\par
佑芽： 什么意思?\par
广： 我们对群感兴趣的， 与其说是该群包含什么操作， 不如说是那些操作\par
间的结构。所以，从现在起，我们来，虑在（轴过来点的）3 维转转\par
操作中，具有与F相同结构的转转操作的集合。\par
千奈： 和F相同， 是指方才的 「具有群的结构， 除此以外不具有非平凡结构」\par
这回事 ?\par
26\par
当然是以 B 为全体，而非 F。\par

% --- PDF page 111 ---
110\par
广： … 。 取 来 两 个 完 全 独 立 的 转 转 。 具 体 来 说 ， 设$\sigma$ 和$\tau$ 分 别 为\par
1\par
7 (\par
6 2 3\par
2 3 -6\par
-3 6 2\par
) ,\par
1\par
7 (\par
2 -6 3\par
6 3 2\par
-3 2 6\par
)那么，，虑由这个$\sigma$和$\tau$（和它们的\par
逆元）组合而成的转转构成的群，就能行得通。 （这个转转操作的集\par
合下面记为 Sato\par
27\par
。）\par
佑芽： 这排着 9 个数字的是啥?\par
广： 矩子。\par
佑芽： …～，感觉起过……线性代数那东西里出来的对吧?\par
千奈：\par
泽同学，我们对线性代数完全不懂啊?\par
广： 认识成， 是一种能用9 个数字表示 3 维转转 （虽然也能表示转转以外\par
的东西）的方法，就足够了。\par
佑芽： 那么， 这两个转转和它们的逆转转的组合， 只要组合次数和便序不同\par
就全都不一样，是这样说的吧?\par
28\par
为什么?\par
广： 证明请让我省略。\par
佑芽： 诶～～～～!!!!\par
广： 我想大概话会变得太长。\par
佑芽： 懂了。那且且，承认这个，可这教应变多有啥关系?\par
广： 在把应（体）变多之前，首，来把应壳\par
29\par
（下面设为S）变多吧。应壳\par
中，除的与组合$\sigma$和$\tau$构成的转转的转转轴相交的点。。 （这个下面设\par
为S∗。）\par
佑芽： ……。\par
广： 现在在这里， 在S∗上致入等价关系\textasciitilde{}， 定义为对p, q $\in$ S∗, p\textasciitilde{}q\par
def\par
$\Leftrightarrow$ $\exists$g $\in$\par
Sato, g(p) = q。设M为S∗/\textasciitilde{}的完全代表系。\par
千奈： 又是咒文啊～!\par
广： 把这个不用集合论用语翻译一下的话就是， 从S∗的要素中选一个要素，\par
将作为 Sato 要素的转转作用上的，把由此经过的点从候选中的掉，\par
然后从剩下的要素中， 选一个没被从候选中的掉的要素， 重复同样操\par
作……如此反复， 实续到S∗的所有要素要么被选中要么被从候选中的\par
掉为止，这样被选中的要素就是M。\par
27\par
取佐藤健治之名，故如此标记。\par
↑译者注：参加参，文献，佐藤建治是\par
28\par
这是「具有群的结构，除此以外不具有非平凡结构」的换种说法。\par
29\par
意思是中空的应。\par

% --- PDF page 112 ---
111\par
佑芽： S∗ 几乎是应壳所有的点对吧?我感觉要把它的要素全都取尽，不可能\par
的吧。\par
广： 不不，可以。没错，只要选择公理在。\par
佑芽： 选择公理啊……几乎只知道名字!\par
30\par
广： 大『能将这类超越操作正当化的免罪符\par
31\par
。\par
佑芽： 了解!\par
千奈： 太、太着一了吧!?\par
广： 由定义， 将Sato 的各元素作用于M， ，虑它们的并集的话， 这就成了\par
S∗32\par
。\par
佑芽： 这是和作没法测定勒贝格测度的集合时同样的感觉， 因为是把从候选\par
中的掉的点全都「收了回来了」对吧。\par
广： …。 把这个形式地记为Sato(M) = S∗吧\par
33\par
。 和才才对F所做的同样， 将\par
Sato 分割为A1, A2, B1, B2的话，A1(M), A2(M), B1(M), B2(M)的并集就\par
是S∗了。\par
千奈： 那倒也是呢。\par
广： 这里，是利用操作的集合其本身曾能通过操作的操作来变多这一点。\par
$\tau$-1 $\circ$ (A1(M)) $\cup$ A2(M) = ($\tau$-1 $\circ$ A1 $\cup$ A2)(M) = Sato(M) = S∗， 同样\par
的B1(M) $\cup$ $\sigma$ $\circ$ (B\_2(M)) = S∗，所以S∗是 Sato-复制可能的。\par
34\par
到这里\par
没问题吧?\par
佑芽： 就是说， 把Sato 的全部要素作用于M就成了S∗全体， 但我们在关注的\par
是，只把 Sato 的部分要素作用得到的碎片，对吧。\par
千奈： 通过M，我们成功把关于操作（Sato 的要素）的讨论，转移到了关于\par
S∗的讨论!\par
佑芽： 但是，这还是关于S∗的话吧?S的话怎么做?\par
30\par
当然，数I是从「集合与拓扑」开始的，首，处理了公理集合论的简史。 （请注意，一般来讲这个\par
描述是有误的。 ）\par
↑译者注：应该是指日本的数学教材数学I，本人对日本教材不甚了解，请有兴趣者自行查阅。\par
31\par
此处如何使用选择公理，与*7 同样。\par
32\par
这里单计说「并集」的必须是非交并集，为此从S的除了与转转轴相交的点。 （详细的证明只能割\par
爱了）\par
33\par
请留意 Sato 其本身是操作的集合，并非操作。\par
34\par
与 Sato(M)这一标记同样，使用了大量略记。\par

% --- PDF page 113 ---
112\par
广： 详细的证明略过，但使用不包含于 Sato 的转转（设轴过来点的 3 维\par
转转的全体集合为SO(3)）中合适的转转，就能和补全圆周缺的 1 点\par
时的讨论同样，填上S∗的。，能够展示S∗与S是SO(3)-裁剪同余的，\par
所以S\textasciitilde{}S∗\textasciitilde{}2S∗\textasciitilde{}2S\par
35\par
成立，S是 SO(3)-复制可能。\par
千奈： 第二个\textasciitilde{}就是才才的S∗是 Sato-复制可能那件事吧。\par
佑芽： SO(3)-复制可能和 Sato-复制可能，不一样没关系 ?\par
广： 写在这个『「-复制可能」前面的操作的集合， 是能用于复制的操作的候\par
选，所以着一增多些是没有问题的。\par
佑芽： 确实! SO(3)比 Sato 广嘛。\par
图 7：M的示意图（这不反映$\circ$的正确作用，也不准确反映S∗所缺的点。 ）\par
千奈： 话虽如此，\par
泽同学，到此为止应壳是变多了，但最初觉说的是把应\par
体变多吧。\par
广： 现在，半径 1 的应壳能变多了，那么出于同样理由，0 以外任意半径\par
的应壳都能变多，所以把（以来点为中心的）半径r （0 < r $\leq$ 1）的\par
应壳收集起来的东西，即除的来点的应体，能用SO(3)复制。\par
佑芽： 这个来点也能补上 ?\par
35\par
此处\textasciitilde{}表示 SO(3)-裁剪同余。\par

% --- PDF page 114 ---
113\par
广： …。，虑让过来点的圆上的点转转那样的（不过来点的）3 维转转，\par
就能用和补圆的 1 点时同样的讨论补上。\par
千奈： 结感觉最后太过干脆，很没有实感，但莫非这就算……\par
佑芽： 证明出来了!?\par
广： …。\par
佑芽： 诶～～～～!?道理是懂了，但好别扭啊～～～～\par
广： 我倒是希望觉能的然。\par
图 8：将补来点的方法可视化的图。在圆上进行的操作与图 3 同样。\par
少许关于2 维的话题\par
佑芽： 结果，F能变多， 以及允许使用的操作中包含F的结构， 这两点很重要\par
对吧?\par
广： 对。\par
佑芽： 结感觉这样的话，圆板也像是能变多!用 2 维的转转，，虑 2 个独立\par
的，比如 1 rad 和$\sqrt{}$2 rad 的转转不就行了?\par
广： 那 2 个的话，存在即使交换操作便序结果也不变这一非平凡的关系\par
性，所以不行。这种不介意操作便序的叫可交换性。\par
佑芽： 啊～!!!确实……\par
广： 而且可以说，那就成了圆板无法复制的根本理由。\par
千奈： 这怎么说?\par

% --- PDF page 115 ---
114\par
广： ，虑把圆板变多时所用的 2 维等距变换群E(2)的元素，并非可交换，\par
对吧?\par
千奈： 比如，向x方向的 1 然后转转90°（逆时针）的结果，和转转90°后再\par
向x方向的 1 的结果，是不同的东西对吧。\par
广： 但是，E(2)是平移、2 维转转以及镜像的组合，而各个分类中的要素\par
各自是可交换的。\par
佑芽： 平移可交换， 是因为向量的加法可交换，所以是那么回事，转转才才\par
说过了。镜像呢?\par
广： 镜像只有翻转或不翻转 2 种，因此不用讨论，就是可交换的。\par
图 9：左为 2 维转转可交换的样子。右为镜像的一例。\par
千奈： ……?没有选择以哪个轴来做镜像的余地 ?\par
广： 严密的内容请看附录。\par
千奈： （附录……?）\par
广： 而众所周知， 像这样由各要素可交换的群彼此组合构成的群， 是无法\par
进行增殖的。所以，2 维中圆板不能变多，与其说因为是 2 维，不如\par
说是因为正试图作用于 2 维空间的E(2)的结构。\par
佑芽： 都到这一步了，那个证明我也想知道!\par
广： 饶了我吧。\par
佑芽： 广酱小气!\par
广： 会赶不上上课的\par
36\par
*。\par
佑芽： 诶～～!!!!\par
36\par
指截稿期，。\par

% --- PDF page 116 ---
115\par
千奈： 那么说， 就是说2 维中也有『「乍看面积好像不会增加， 却能让圆板增\par
殖」的法法喽?\par
广： …。如果剪切也允许的话，圆板可以变多。\par
佑芽： 剪切是啥?\par
图 10：剪切的一例。\par
广： 比如，三角形的面积是底∗高/2，只要不改变这两个，像这样往横侧\par
挪，面积也不变对吧。这就是剪切\par
37\par
。\par
千奈： 和应体的时候一样，我想知道这个分割成怎样的碎片就能变多!\par
广： 千奈，上课\par
38\par
。\par
千奈： 明、明白了……\par
结束\par
广： 时间上也好，话题的段落也好，都才好，就到这里吧。\par
佑芽： 广酱……\par
佑芽： 没有选择公理就证明不了\par
39\par
这点呢??????\par
广： …。\par
千奈： 是啊，\par
泽同学!我们，对巴拿赫-塔斯基定理，除了没有选择公理就\par
证明不了以外，什么都不知道啊!\par
40\par
广： ………我们平时通常会思，「从设定能致出怎样的结论」 ，但要思，\par
「能具从这设定致出这结论」这种元问题，就需要用到平时不会使用\par
的特殊技巧……\par
37\par
请注意，剪断后面积不变的情况不仅，于三角形。 （参见卡瓦列里来理）\par
38\par
同*36\par
39\par
在 ZF 上的意思。ZF 是公理系的名称，不含选择公理。\par
40\par
同*39。\par

% --- PDF page 117 ---
116\par
佑芽： 能力不足……就这么回事!?\par
41\par
广： 诶～…………呵呵。\par
广： 真不如意呢。\par
附录\par
1.勾股定理的证明\par
，虑直角三角形ABC($\angle$C = $\angle$R)。 设各边长为AB = c, BC = a, CA = b。 在半直线CB\par
上取点B$\prime$， 满足CB$\prime$ = a + b。之后，过B$\prime$且垂直线段CB的线段上， 在夹CB$\prime$与A同侧处，\par
取点D满足B$\prime$D = a。\par
这里，CB = B$\prime$D = a, AC = BB$\prime$ = b, $\angle$ACB = $\angle$BB$\prime$D = $\angle$R， 所以由两边夹角相等，\par
$\triangle$ ACB与$\triangle$ BB$\prime$D全等。therefore $\angle$CAB = $\angle$B$\prime$BD。\par
因而，$\angle$ABD = 2$\angle$R - ($\angle$CBA + $\angle$B$\prime$BD) = $\angle$R。\par
这里，$\triangle$ ABD的面积(=\par
1\par
2 c2)是\par
梯形ACB$\prime$D的面积- $\triangle$ ACB的面积- $\triangle$ BB$\prime$D的面积\par
= 1\par
2 (a + b)(a + b) - 1\par
2 ab - 1\par
2 ab\par
41\par
从某种意义上说确实如此，但大多数数学学习中并不需要这类知识。\par

% --- PDF page 118 ---
117\par
= 1\par
2 (a2 + b2)\par
如此表示。\par
以上，c2 = a2 + b2。\par
这里，暗中承认了碎片合计的面积等于碎片面积的合计。 （另外，梯形面积和直角\par
三角形面积如上方所示被表示，是利用面积在转转、平移下不变而致出。 ）\par
2.选择公理的陈述与完全代表系\par
选择公理是如下之物。\par
对任意集合（的族\par
42\par
）X，存在满足以下条件的选择函数f : X\textbackslash{}\{$\emptyset$\} →$\cup$\{x$\in$X\} x。\par
对任意非空集合x $\in$ X，f(x) $\in$ x。\par
使用此公理，对[0,1]/\textasciitilde{}（\textasciitilde{}的定义请参见正文中的相应部分） ，存在选择函数\par
f : [0,1]/\textasciitilde{} →$\cup$\{x$\in$[0,1]/\textasciitilde{}\} X(= [0,1]) 。 这 里 设V = \{f(X) $\in$ X|X $\in$ [0,1]/\textasciitilde{}\} ， 则 这 成 了\par
[0,1]/\textasciitilde{}的完全代表系。\par
现在展示这点。对[0,1]/\textasciitilde{}的任意元素X，当存在两个元素x1, x2都包含于V时，展示\par
它们相等即可。\par
由两元素包含于V ，意味着存在某[0,1]/\textasciitilde{} 的元素X1, X2，f(X1) = x1 且f(X2) = x2 。\par
由此，X $\cap$ X1 $\neq$ $\emptyset$且X $\cap$ X2 $\neq$ $\emptyset$被致出。 （请留意x1, x2都是X的元素。）\par
由商集的定义，可以得出，商集的两个元素有交时它们相等\par
43\par
，所以X = X1 = X2 ，\par
且x1 = f(X1) = f(X2) = x2。\par
在此， 也就更近自然语言的解的略述一二。[0,1]/\textasciitilde{}的元素里各自包含无，个元素。\par
（例如，包含于[0,1]的实数中，将$\sqrt{}$2\par
2 的差为有理数的全部收集来的作为[0,1]/\textasciitilde{}的元素，\par
其中包含$\sqrt{}$2\par
2 -\par
1\par
2 , $\sqrt{}$2\par
2 -\par
1\par
3 , $\cdots$等无，个要素。 ）现在想从这各自的元素中选出 1 个元素来作\par
完全代表系，但从各自的元中选什么出来，比如从与$\sqrt{}$2\par
2 差为有理数、0 以上 1 以下的实\par
数的集合中，选出的是$\sqrt{}$2\par
2 -\par
1\par
2还是$\sqrt{}$2\par
2 -\par
1\par
3，还是别的东西，这必须定定。\par
不是有，个，我们能简单提供其基准。例如，取最小的东西就好。然而，一旦变成\par
无，个，要想出那样的好基准似乎很难，且实际以建构的方法未必结能法到。于是，选\par
择公理就贡献于「基准的提供」 。 「由选择函数所选出的东西」这一基准， （某种意义上\par
仿佛从天而降地）可以将之用于这类选定的基准，选择公理持证了这一点。\par
42\par
在公理集合论中，万物皆被想作集合，所以某对象是集合还是集合的族，不过是语下问题。\par
43\par
此证明好为冗长，且且割爱。\par

% --- PDF page 119 ---
118\par
3.群的可解性\par
广：像这样由各要素可交换的群彼此组合构成的群，无法进行增殖，这是已知的。\par
这意味着，可解\par
44\par
的群是便从的。 （群是便从的与不产生悖论有深刻关系。 ）\par
E(2)是可解的。即E(2)具有如下的正规部分群列：\par
E(2) $\geq$ SE(2) $\geq$ T(2) $\geq$ 1\par
且商群\par
\{1, -1\}, SO(2), T(2)\par
各自是可交换的。其中，SE(2), SO(2), T(2) 分别是特殊欧几里得群、特殊正交群、\par
平移群。\par
广：「镜像只有翻转或不翻转 2 种，因此不用讨论，就是可交换的」这句话的意思\par
是\{1, -1\}对应于（排除了平移和转转要素的）镜像的要素。\par
参考文献\par
Grzegorz Tomkowicz, Stan Wagon\par
https://alg-d.com/math/ac/ac.html (2024/10/11\par
)\par
https://youtu.be/I3\_9oAmIv04 (2024/10/11\par
)\par
44\par
不仅如此，本节所用的群论术语将不作解的。敬请理解，此为面向已掌握相关知识的读者而设的\par
附录。\par

% --- PDF page 120 ---
119\par
请\par
泽广教教我!狭义相对论\par
Alkyne\par
2024 年 11 月 30 日\par
Day0 研讨会的开始\par
某天练习的休息时间。我为了吃午饭而走进休息室，发现\par
泽同学正坐在椅子上，\par
不知往笔记本上写些什么，而仓本同学则在她身旁，脸色发青。仓本同学一看到我，一\par
立刻跑了过来。\par
千奈「老、老老、老师!!!\par
泽同学她从才才起样子就很奇怪啊!!!」\par
千 P「仓本同学。\par
泽同学怎么了 ?」\par
千奈「那、那个，她从才才起就一直在笑着做笔记，我教她说话她也没反应……」\par
起她这么一说， 我看向\par
泽同学， 确实， 她微微张着嘴， 脸上挂着异常开心的笑容，\par
铅笔一直在的。的确，旁人看来，这明无是个怪人。\par
千奈「或、或许，是因为接受了过于严酷的训练，大脑产生了还在一直训练的错觉\par
吧!!!得赶快把她送到持健室——」\par
千 P「请、请冷静一点。 」\par
我一边安抚慌张的仓本同学， 一边偷看\par
泽同学的笔记本。 上面， 在摊开的两页上，\par
密密烦烦地写满了一长串似乎是数学符号的文字。\par
广「做出来了，呵呵……」\par
千 P「训练辛苦了。居然摊开笔记本，真少见啊。 」\par
我趁\par
泽同学才好写完一长串式子的时机，试着教她搭话。说来，我可能还是第一\par
次看到她这副模样。回想起来，\par
泽同学 14 岁就从海外大学毕业，作为物理学是前途\par
备受期待， 这一天才事迹在学园内曾轰的一时， 但我却从未见过她在学园里拿着物理学\par
的教科书或笔记本。\par
广「啊，千奈的制作人。 」\par
千 P「仓本同学很担心觉，说教觉说话也没反应。觉在做什么呢?」\par
广「做什么，如觉所见，只是在做表达式变形而已。 」\par
千 P「为什么觉要那样笑嘻嘻地做呢?」\par
起我这么问，她再次向我投来之前对着笔记本时的那种闪闪发亮的眼如。\par

% --- PDF page 121 ---
120\par
广 「呵呵……我现在是在做这本教科书上写的式子的中间表达式变形， 但这非常困\par
难，很辛苦，又很快乐。 」\par
说着，\par
泽同学从包里拿出一本书。那是一本白底带弯曲蓝条的、看起来很难的物\par
理学书。封面上大大地写着《场的经典理论》 。\par
千 P「咦，觉不用看书也能做 。 」\par
广「…。就算不看书，这种程度的式子我也能再现。 」\par
这种程度的式子……?怎么看都不像是「这种程度」的东西，但真不愧是她那天才\par
的样子。\par
千奈「真是的!教觉说话至少给点反应啊!我好担心觉啊!」\par
仓本同学也略带怒气地加入了对话。\par
千 P 「说起来，我之前起志摩，生说过，觉开始当偶像是因为不擅长，而且平时做\par
的物理学已经做腻了。 」\par
志摩，生，指的是\par
泽同学的制作人。\par
广「…，是啊。 」\par
那她为什么现在又这么开心地做式变形呢。\par
广「确实，比起物理，当偶像更快乐。但是，并不是说我讨厌物理学，只是不想过\par
无聊的日子罢了。 」\par
千 P「是这样啊。那觉现在也在学习物理学 ?」\par
广「…。 」\par
千 P「嘿。那觉书架上的书全都是物理学相关的 ?」\par
广「那倒不至于。我在大学大体都学过了。虽然有时会读读论文什么的。而且，就\par
算是我也会读小说啦着笔啦，普通的文库本也会读的。 」\par
千 P「我还以为觉肯定只会读物理学和数学的书呢……」\par
广「那是偏见。 」\par
确实， 这或许是我的看法有点偏颇了。志摩，生很喜欢物理，时常会把和她对话的\par
内容讲给我起，但每次内容都太过浓密且高深，是些物理和数学的话题，害我不禁以为\par
他、她、他们一直都在思，这类东西。\par
广「便一一提，最近基本只在看论文，没怎么读物理学的书。 」\par
千奈「之前觉和制作人，生聊天的时候，说过大学的物理学书大多很无聊对吧。 」\par
广「…。至今为止在高中和大学读过的大学物理书，每一本都浅无易懂，是很不错\par
的书，但行间太窄，没有做的价。，老实说很无聊。但这本书，很厉害。这本书行间很\par
宽，做表达式变形很有价。，我一直都很兴奋。 」\par
泽同学稍微探出身子，语速也变快了。\par

% --- PDF page 122 ---
121\par
千 P「是、是 。嘿……。……话说回来，式子还真是长啊。 」\par
广「是 ?」\par
千 P 「是的。因为我只学到高中数学为止。特别是这里，一个多项式据了8 行，我\par
见都没见过。 」\par
不过，这式子可真长啊。就算是喜欢数学的我，也被震撼到了。\par
千奈「这是关于什么的式子的变形啊?」\par
广「粗略来说，是广义相对论。 」\par
千奈「广……广义相对论?」\par
广「对。相对论，觉起过 ?」\par
千奈「啊，起过!就是那个，叫爱因斯坦的人创立的很厉害的物理学理论对吧?」\par
广「对。 」\par
千 P「为什么觉现在在学相对论呢?」\par
广「说是学习，不如说，是我的制作人问我要不要用这本书开个研讨会，我这是在\par
做准备。 」\par
千 P「研讨会?」\par
广「…。据说是训练的一环。 」\par
真是相当奇特的训练。\par
广「我，没什么体力。一开始按照制作人的指示，少进行些练习，剩下的时间做想\par
象训练，但作为范本的偶像视频我基本都看完了，素材用光了。 」\par
千 P「是这样啊。那练习的态况如何?」\par
广「老样子。 」\par
千 P「也就是说?」\par
广「体力耗尽，然后倒下了。 」\par
千 P「哎……」\par
广「制作人每次练习时都用干巴巴的眼如看着我，十分火大。」\par
千奈「结觉得有非常险…的气息啊!」\par
这点也还是老样子。\par
千奈「那么，在研讨会?……上，觉们训练些什么呢?」\par
广「当然是体力。 」\par
千奈「体力?」\par
广「…。虽说是体力，或许该叫毅力吧。体力很少的我，努力过绪就会马上倒下。\par
所以，为了不倒下来，要锻炼心灵的体力。呵呵，非常合理。 」\par

% --- PDF page 123 ---
122\par
说起来，我虽觉得她来本就是个颇有毅力的人，但对她本人而言，是因为喜欢挑战\par
不擅长的事， 所以既不觉得辛苦， 也并非靠毅力在耐耐， 这点我现在才注意到。 ………?\par
但仔细一想，她都笑嘻嘻地在做式变形了，这根本不能算是毅力的训练吧……?\par
千奈「呵呵呵，\par
泽同学的制作人，生，非常为\par
泽同学着想呢。真棒啊!」\par
千 P「觉这话是说我不太为仓本同学着想 ?」\par
千奈「诶，诶诶!?不、不是，我不是那个意思，那个，该怎么说呢……请恕我失礼\par
了～!!」\par
千奈「……话说回来，相对论，这个词起起来好帅啊!我，很感兴趣!」\par
仓本同学双眼闪闪发光地说道。\par
千奈「那个，研讨会?……我也想试着做做看!请务必让我加入!」\par
千 P「诶诶!?话虽如此，但觉高中数学之类的还没学吧……?」\par
千奈「是的，………」\par
千 P「那不会相当困难 ……?」\par
广 「是啊。 如果对数学一窍不通， 那这个研讨会恐的还是稍微有点难。 千奈， 抱歉。 」\par
千奈「说、说得也是呢……」\par
被勾起了好奇心的仓本同学的吸收能力令人刮目相看， 但不是对高中数学一窍不通，\par
恐的还是相当困难吧。\par
广「…。不过，虽然用这本书开研讨会确实很难，但我觉得相对论本身，觉是能理\par
解的。 」\par
千奈「真的 !?」\par
仓本同学的眼中再次点亮了光芒。\par
广「说起来，千奈的制作人，数学有多拿手?才才我起到觉说学过高中数学了。 」\par
话题突然抛向了我。\par
千 P 「诶?啊，是的。数学我既拿手又喜欢，且且选了理科，学到了数III。不过在制\par
作人科，那些知识并没怎么用到。 」\par
广「来来如此。便一，觉放学后有空 ?」\par
千 P「…，不过是到晚上了。 」\par
广「千奈呢?」\par
千奈「我也，和\par
泽同学的训练日程基本一样，没训练的日子多半是闲着的……」\par
广「是 。那样的话……」\par
说着，\par
泽同学取出了智能手机。\par
千 P「怎么了?」\par
广「教制作人确认一下。……啊，已经回信了。 」\par

% --- PDF page 124 ---
123\par
不愧是制作人科的人，回信真快。\par
广 「………， 好像没问题。 而且我的制作人放学后似乎也相应有空。 那样的话……」\par
稍过片刻，\par
泽同学通过信息发来了链接。\par
广「那样的话，晚上就连线通话，再开一个研讨会吧。点这个链接，视频通话就开\par
始了。日程就酌情调整。 」\par
千 P「请、请等一下。那样的话，\par
泽同学的休息时间就……」\par
广「物理就是休息，所以没关系。教制作人也确认过了，他反倒说请务必作为训练\par
的一环让我来做。 」\par
我一瞬间怀疑，那绝对是打着训练的幌子，其实只是想传播物理学吧。但，他大概\par
也有他自己的想法吧。而且，我学生时代也未曾想过自己会进入这样的专科大学， 也是\par
曾对大学里那些艰深的数学和物理抱有兴趣的人，所以趁此机会学习一下也不坏。\par
千 P「…我明白了。那么，我就和仓本同学一起参加吧。 」\par
千奈「太好了～!好期待啊～!」\par
就这样， 以回是路上绕个远路便一逛逛般轻松的心情， 相对论研讨会的举法就定了\par
下来。\par
Day1 数学准备\par
一进入视频通话，\par
泽同学已经和志摩，生在说话了。\par
广 P(志摩)「……那么，一周前布置的舞蹈课题，觉完成到什么程度了?」\par
广「还不到两成。 」\par
广 P「……来来如此。 」\par
广「那个停顿是什么意思。 」\par
广 P「……老实说，我来本希望觉能至少完成五成左右……」\par
广「呵呵……失望了?」\par
广 P「…，算是吧。 」\par
广「制作人，结是很严格呢。」\par
广 P「因为要是表扬觉进步了，觉不就要情…低落了 。 」\par
广「并没有那回事。不过，起到严厉的意见，我很高兴。 」\par
广 P「真是个难懂的人啊……」\par
泽同学，果然是个不可思议的人。不论是在这所初星学园内外，我根本就没遇见\par
过这种类型的人。\par
千 P「辛苦了。 」\par
广 P「啊，辛苦了。从今天起的研讨会，请多关照!」\par

% --- PDF page 125 ---
124\par
千 P「话说回来，\par
泽同学为什么叫了志摩，生来这个研讨会呢。 」\par
广「我没叫。 」\par
广 P「是的。是我自己申请，希望让我参加的。 」\par
果然只是想传播物理学吧。\par
广 P「主要言定是由\par
泽同学来说明，我则打算做做补充说明之类的。 」\par
广「制作人，明明平时结对我那么严格，一牵扯到物理，就特别积极。 」\par
确实。\par
广 P「我有在好好为\par
泽同学着想。?这也是训练的一环。 」\par
广「嘿。……」\par
广 P「……有什么不满 ?」\par
广「前子子，休息日，我说想的水族馆，觉却拒绝我了。 」\par
广 P「休息日就该好好让身体休息。 」\par
广「……制作人。 」\par
广 P「什么?」\par
广「下周一定要的，水族馆。 」\par
广 P「觉这么想的啊。明白了，就按这个计划来安排吧。 」\par
这瞬间窥见了志摩，生培育\par
泽同学时的辛苦。 之后又过了一会， 仓本同学也进来\par
了。看来，她现在是按照\par
泽同学的指示，在请是庭教师特别教授她相对论所需的高中\par
数学知识。话说回来，她那想到就做的行的力，以及将之付诸实践的财力，每每令人，\par
叹。\par
千奈「那么，现在，请给我们讲相对论的故事吧!」\par
广「…。不过，在那之前，我想，从数学讲起。因为千奈想知道，所以一开始会配\par
合千奈的水平来讲解。 」\par
说着，\par
泽同学立刻通过画面共享，将在线服务的白板映在了视频画面上。\par
千 P「那么，研讨会期间，我就在这边安静地守护仓本同学吧。 」\par
我一边这么说，一边也翻开了笔记本。 既然机会难得，我也悄悄和千奈一起学习相\par
对论好了。\par
广 P 「我也，只会在需要补充说明时插嘴，基本打算安静地旁起，所以研讨会就请\par
以\par
泽同学和仓本同学为主体进行吧。 」\par
广「知道了。那么千奈，我们开始吧。 」\par
千奈「拜托觉了!」\par
～（研讨会开始）～\par

% --- PDF page 126 ---
125\par
广「千奈，已经学过坐标了 ?」\par
千奈「是的，且且学过了!」\par
广「来来如此。那样的话，今天的话题看来没问题。今天是坐标的话题。 」\par
广「首，，从1 维开始，来想想距离的测量方法。有点突然，千奈现在的身高是多\par
少?」\par
千奈「是 148 cm!」\par
广「来来如此。那和台应的应杆长度大『一样呢。 」\par
千奈「……诶?」\par
广「也就是说，用台应应杆来量长度的话，千奈大概长度就是 1 了。 」\par
千奈「请、请等一下!!!就没人吐槽 ～!?」\par
广 「嘛， 才才的， 是关于长度基准的一个小例子。 这里我想说的是， 以什么为基准，\par
测量长度的方法就会改变。 也就是说， 千奈的身高， 说148cm 也好， 说1 根应杆也好，\par
都可以。但用来表示那个长度所用的数字，这个例子里，就是 148 cm 或 1 根应杆之类\par
的数字，是会着基准而变的。 」\par
千奈「……嘛、嘛，确实是那么回事呢。 」\par
广「坐标也是一样的。也就是说，在某一方，1 的长度，在另一方那里并不一定非\par
得是 1。到这里没问题吧。 」\par
千奈「没问题。。 」\par
广「那么，接下来是2 维。变成 2 维之后，不仅距离的基准，其他的基准也可能变\par
化，觉觉得什么会变?」\par
千奈「………我不知道……」\par
广「那，试着画一下 2 维的坐标吧。 」\par
千奈「………是这样的 ?」\par

% --- PDF page 127 ---
126\par
广「…，这样倒是可以，那我提问。坐标轴的 x 和 y，为什么要直角相交呢。 」\par
千奈「高中数学里是这么教的吧……?」\par
广「那倒也是，可直角相交有必要 ?」\par
千奈「啊，确实!斜着相交也没问题啊!」\par
广「对。也就是说，除了距离的基准，两根轴的相交方式的基准也可能改变。 」\par
千奈「确实……啊，我又发现了一点!坐标轴也未必非得是直线啊!就算坐标轴是弯\par
的，感觉也能定义坐标!」\par

% --- PDF page 128 ---
127\par
广 「。。。 发现得不错。 其实这个今后会提到， 但我觉得现在还太难， 所以往后放。\par
现在且时把轴当成直的东西来，虑。 」\par
广「那么，从至今为止的讨论来看，可以明白，在一般的n 维里也同样可以有各种\par
坐标的取法。试着写一下 3 维坐标的取法，应该就能整理出至今的内容了。到了 n$\geq$4\par
的时候，就很难具体画图来想了，但放心，用同样的思，方式不会有问题。 」\par
千奈「我明白了，坐标可以有各种取法。那么，坐标一变，位置的表示法会怎么样\par
呢?」\par
广「现在就来讲解这个。千奈，觉懂向量 。 」\par
千奈「啊，稍微懂一些!就是箭绪对吧?加法、减法，还有乘以正数倍啦负数倍啦之\par
类的也会!内积?还有外积?什么的我还不太懂……」\par
广「对。就是箭绪。箭绪有『「方向」和『「长度」两个要素。只要明白这点，就基本\par
没问题了。 」\par
广 「现在， ，来，虑轴正交的坐标吧。 设设坐标的来点O 处有觉自己， 要走到点P\par
的。 」\par
广「那么，首，是第 1 题。设设千奈现在只能沿着这个 x 轴、y 轴的方向走。这时\par
候，千奈能到达点 P  ?」\par

% --- PDF page 129 ---
128\par
千奈「能，这样走就能到了!」\par
广「答对了。那么第2 题。不论点P 在哪里，千奈遵守才才的，制规则，都能到达\par
点 P  ?」\par
千奈「…——，是的，我是这么想的。 」\par
广「来来如此。那觉为什么这么想呢?」\par
千奈「觉问我为什么，…——，很难……但是，结感觉是那么回事。 」\par
广 「在这种情况下， 想的的地方， 也就是点P 的坐标可以表示为(x, y) = (a, b)这样\par
的感觉。比如，，沿 x 走 a，再沿 y 走 b，按这种感觉走，就能满足规则移的过的。所\par
以，不论点 P 在什么地方，都能移的过的。 」\par
千奈「来来如此，的确是呢!」\par
广「那么，第 3 题。接下来有点难，但像才才说的那样，把坐标轴弄斜吧。 」\par

% --- PDF page 130 ---
129\par
广「像这种情况时，不管点 P 在哪里，都能到达 ?」\par
千奈「…——，这个也很难，但还是用教才才一样的想法试试看呢。画通过点P 且\par
平行于 x 轴、y 轴的直线的话………，好像能行!」\par
广「就是这样。斜的似乎也没问题。那么，最后的问题。其实，在斜的情况下，有\par
时会变得只能的有，的地方。那是什么时候呢?」\par
千奈「诶!?那、那有点不知道啊……」\par
广「那给个提示。轴倾斜过度的话会怎么样呢。 」\par
千奈「…——……啊!轴就会重叠在一起!」\par
广「就是那样。一旦轴重叠，就只能往那重叠的轴的方向走了。正确答案就是，轴\par
重叠的情况。 」\par
千奈「来来如此，完全明白了!」\par

% --- PDF page 131 ---
130\par
广 「那么， 在这里稍微讲解一下术语。 至今的问题里， 出现了『「沿轴前进」的规则，\par
但所谓前进， 必须知道是 「往哪个方向」、『「走了多少长度」。 于是， 把『「往轴的正方向」、\par
「走了 1 的长度」的向量称为基向量。 注意， 基向量本身的数量和轴的数量一『。 还有，\par
「1 的长度」这东西，也像才才那样，会根据坐标的取法而变，这点也请注意。 」\par
广 P 「便一一提， 把才才的内容换个说法， 就成了 「能具通过组合基向量的常数倍，\par
到达坐标上的各个点」这个话题。也就是说，平常我们说的坐标的 x 分量、y 分量，就\par
表示要到达那个点，需要按各轴的基向量，走多少长度。另外，把基向量用分量表示，\par
就成了(1\par
0)和(0\par
1)。这从定义看是不言而喻的。 」\par
广「…。还有，这个希望觉记住，像这样能的到任何地方的基向量的组合，称为线\par
性无关。基向量的组合，结之就是轴的取法。另一方面，只能的特定地方的情况，比如\par
才才那个轴重叠了的情况，这样的基向量的组合，就称为线性相关。 」\par
千奈「……。也就是说，是这样一回事呢!」\par
广「…，2 维的情况下是这样。便一，3 维的时候，想象一下大概就是这种感觉。 」\par
千奈「啊，确实，3 维的情况下，会有只能在平面上移的的情形!真是盲点啊!」\par
广「相对论里只处理线性无关的情况， 所以线性相关忘了也没事。在这里希望觉记\par
住的是， 「只要别取奇怪的坐标取法，哪个点都能的到」。能的到点，就意味着那个点的\par
坐标可以用分量表示，记住这点。 」\par
千奈「是!!」\par
广「那么接下来，讲解一个叫矩子\par
行列\par
的东西。 」\par
千奈「行列……?是拉面店前面排的那种 ?」\par

% --- PDF page 132 ---
131\par
广「不不，不是那个，是把数字纵横排列起来的东西。便一，觉学过把向量的分量\par
竖着排列来写 ?」\par
千奈 「啊， 是的， 学过了……但是， 我不太明白为什么要纠结是竖着还是横着……」\par
广「来来如此。那么，那个疑问，学了矩子应该就能解定了吧。在之后的讨论里，\par
向量的分量就定成竖着排列来写吧。 」\par
广「那么，这次试着来想想这样一件事吧。至今为止，千奈都是以自己为中心的直\par
角坐标来测量东西的位置什么的，但有一天， 制作人对觉说 「要从更多样的角度的。改\par
竞争对手」，于是觉从直角坐标，换成了以斜交的斜交坐标来测量东西的位置。这里，\par
坐标的中心当然是固定在自己身上，所以设设变更前后是不变的。 」\par
千奈「也就是说，是把自己的坐标的基向量给替换掉了，对吧!」\par
广「对对，就是那么回事。此时，设设在替换前的世界，也就是直角坐标里，测量\par
了替换后的世界，也就是斜交坐标的基向量是什么样的向量。 」\par
广 「来来直角坐标的基向量是(1\par
0)和(0\par
1)。 这从定义看是理所当然的。 设设这被替换\par
成了在直角坐标中表示为(a1\par
b1\par
)和(a2\par
b2\par
)的向量。此时，将这个坐标的变换写作下式。 」\par
(a1 a2\par
b1 b2\par
)\par
广「这就是基的变换矩子。 」\par
千奈「???」\par
广「有点不太明白吧。其实这个矩子经常用在向量的计算上，那么就来讲解一下它\par
的规则吧。 」\par
广「首，，从向量的内积开始。内积的做法是，同分量彼此相乘之后，把它们全部\par
加起来。例如，如果有两个含2 个分量的向量，结果如下。符号$\cdot$的意思教乘号$\times$或加号\par
+差不多，表示「计算这个向量和这个向量的内积」。」\par
(x1\par
y1\par
) $\cdot$ (x2\par
y2\par
) = (x1\par
y1\par
)\par
T\par
(x2\par
y2\par
) = (x1 y1) (x2\par
y2\par
) = x1x2 + y1y2\par
广「式子里出现的，第一个向量右上角标的T，是个表示把列向量变成行向量的符\par
号。也就是说，列向量右上角标了这个符号的话，就变成行向量就好。我们定义，当行\par
向量在左侧，列向量在右侧时，就像那样进行乘法。有 3 个分量的话，就是这样。 」\par
(\par
x1\par
y1\par
z1\par
) $\cdot$ (\par
x2\par
y2\par
z2\par
) = (\par
x1\par
y1\par
z1\par
)\par
T\par
(\par
x2\par
y2\par
z2\par
) = (x1 y1 z1) (\par
x2\par
y2\par
z2\par
) = x1x2 + y1y2 + z1z2\par

% --- PDF page 133 ---
132\par
广「一般地，有n个分量的话，可以这样写。这里，$\sum$  n\par
j=0 是个表示「j 从 1 到 n 变\par
化，请把它们全加起来」的符号。 」\par
(\par
a1\par
$\vdots$\par
an\par
) $\cdot$ (\par
b1\par
$\vdots$\par
bn\par
) = (\par
a1\par
$\vdots$\par
an\par
)\par
T\par
(\par
b1\par
$\vdots$\par
bn\par
) = (a1 $\cdots$ an) (\par
b1\par
$\vdots$\par
bn\par
) = $\sum$ ajbj\par
n\par
j=0\par
广「希望觉且且记住的是，有2 个同样维数的向量，当行向量在左侧、列向量在右\par
侧时，就计算内积。 」\par
广「那么，来道题。请计算(\par
1\par
2\par
3\par
4\par
)和(\par
5\par
4\par
3\par
2\par
)的内积。 」\par
千奈「是这样 ……?」\par
(\par
1\par
2\par
3\par
4\par
) $\cdot$ (\par
5\par
4\par
3\par
2\par
) = (\par
1\par
2\par
3\par
4\par
)\par
T\par
(\par
5\par
4\par
3\par
2\par
) = (1 2 3 4) (\par
5\par
4\par
3\par
2\par
) = 5 + 8 + 9 + 8 = 30\par
广「…，答对了。 」\par
广「那么，接下来是矩子的计算。首，说明一下矩子的术语，矩子的切取方式，有\par
用行向量来切和用列向量来切两种。这些切出的行向量叫行，列向量叫列。然后，行和\par
列如果不清楚是第几行、第几列，就无法指定向量。因此，以左上为基准，行的话就指\par
定是从上数第几行，列的话就指定是从左数第几列，称为 「第i行的向量」啦， 「第j列的\par
向量」啦。 如果想指定第i行、 第j列的分量， 就叫第i行j列的分量什么的。 矩子本身， 不\par
结共有m行、n列，则称为m行n列的矩子或m $\times$ n矩子。尤其m = n时，称为m(n)次方\par
子。我在白板上，写了个 5 行 5 列的矩子，也就是 5 次方子的例子。 」\par
1\par
1\par
译者注：此处作者使用   和heart来指代行列数，译者出于美。和讨论方一更改为了 i 和 j，此后同类\par
更改不再赘述。\par

% --- PDF page 134 ---
133\par
广 「接下来是计算， 首，关于矩子和矩子的乘法规则和结果， 有些事想让觉，知道。\par
就是下面这些。 」\par
1. 乘积结果矩子的第i行第j列分量，是左侧矩子的第i行向量与右侧矩子的第j列\par
向量的内积的。。\par
2. 左侧矩子的列数，和右侧矩子的行数，如果不一『，乘法来本就无法进行。\par
3. l $\times$ m矩子与m $\times$ n矩子的乘积结果，会变成l $\times$ n矩子。\par
4. 方子彼此相乘时，左右互换，结果一般会变。\par
千奈「渐、渐渐变得复杂起来了啊……」\par
广「乍看是复杂的规则，但其实，只要记住第1 条和第 3 条规则，其他规则都能广\par
致出来。一开始很难，我们来，虑 2 次方子彼此的积吧。 」\par
(3 5\par
2 6) (-1 4\par
1 0) = (2 12\par
4 8 )\par
广「这个这样想会容易计算。 」\par
(3 5\par
2 6) (-1\par
1 |4\par
0) = (2 12\par
4 8 )\par
广 「实际来看看计算吧。 首，， 这个左侧是2 $\times$ 2矩子， 右侧是2 $\times$ 2 矩子， 所以积\par
的结果是 2$\times$2 矩子。第 1 行 1 列分量的计算，就是这么回事。 」\par
(3 5) (-1\par
1 ) = -3 + 5 = 2\par
广「同样地，第 1 行 2 列分量的计算，就是这么回事。 」\par
(3 5) (4\par
0) = 12 + 0 = 12\par
广「第 2 行也，同样计算就能求出来，试试看。 」\par
千奈「是!…，第 2 行 1 列是……」\par
(2 6) (-1\par
1 ) = -2 + 6 = 4\par
千奈「第 2 行 2 列是……」\par
(2 6) (4\par
0) = 8 + 0 = 8\par

% --- PDF page 135 ---
134\par
千奈「做好了!」\par
广「…。就是那样。那么，接下来试试算这个吧。 」\par
(\par
1 5\par
2 -1\par
3 4\par
) (-1 6\par
-2 1) =\par
千奈「…，首，结果是，左侧是 3 $\times$ 2矩子，右侧是2 $\times$ 2 矩子，所以应该是3 $\times$ 2\par
矩子!各自计算之后……做好了～!」\par
(\par
1 5\par
2 -1\par
3 4\par
) (-1 6\par
-2 1) = (\par
9 11\par
-4 11\par
5 22\par
)\par
广「就是那样。那么，接下来是这个。 」\par
(\par
1 5\par
2 -1\par
3 4\par
) (\par
-1 6\par
-2 1\par
-3 0\par
) =\par
千奈「同样地思，，首，结果是，左侧是3$\times$2 矩子，右侧是3$\times$2 矩子，所以应该\par
是 3$\times$2 矩子，计算一下……咦?结觉得好奇怪啊……」\par
广「什么奇怪?」\par
千奈「内积没法取……」\par
广「就是那样。内积，不非向量的分量数相同，是无法计算的。这就等于，规则的\par
第 2 条，是由第 1 条规则致出的。 」\par
(\par
1 5\par
2 -1\par
3 4\par
) (\par
-1 6\par
-2 1\par
-3 0\par
)→无法定义!!!!\par
广「那么，接着是这个。把才才例子里的左右互换后的。这会怎么样呢。 」\par
(-1 4\par
1 0) (3 5\par
2 6) =\par
千奈「我来计算看看……咦咦咦?」\par
(-1 4\par
1 0) (3 5\par
2 6) = (5 19\par
3 5 ) $\neq$ (2 12\par
4 8 )\par
广「…，那个就是正确答案。 」\par
千奈「为什么不一『呢?」\par

% --- PDF page 136 ---
135\par
广「那是因为，与才才矩子的切取方式不同。 」\par
(-1 4\par
1 0) (3 5\par
2 6) = (5 19\par
3 5 )\par
千奈「啊，真的!是因为纵横不同才变成这样的啊!」\par
广 「对。 这个性质虽说是第4 条规则， 但也能看出， 它也是从第1 条广致出来的。 」\par
千奈「嘿……结感觉好像稍微懂点矩子了。 」\par
广「太好了。关于矩子计算的话题到此结束。那么，，休息一下吧。 」\par
～（休息）～\par
广「……那么，重新开始研讨会。。千奈，还记得定义矩子时的内容 ?」\par
千奈「完全不记得了……」\par
广「那么，复习一下。我们来，虑，将。测位置的坐标，从直角坐标替换成斜交坐\par
标。这样吧，直角坐标的轴叫x轴和y轴，斜交坐标的轴叫x$\prime$轴和y$\prime$轴。设设斜交坐标系\par
里， 也就是交换后的基向量， 用直角坐标系来量， 基向量由直角坐标系下的(1\par
0)和(0\par
1)分\par
别替换成了(a1\par
b1\par
)和(a2\par
b2\par
)。」\par
广「首，，来想一下在斜交坐标系的世界里，坐标为(x$\prime$, y$\prime$) = ($\alpha$, $\beta$)的位置向量。\par
此时，这个位置向量不在来来的直角坐标系里测量，会变成什么样子呢?」\par
千奈「……x'轴的基向量是(a1\par
b1\par
)，y'轴的基向量是(a2\par
b2\par
)所以……应该变成这样!」\par
$\alpha$ (a1\par
b1\par
) + $\beta$ (a2\par
b2\par
)\par

% --- PDF page 137 ---
136\par
广「…，就是那样。那么，这个能用矩子和向量表示 。 」\par
千奈「诶!?那、那有点……」\par
广「沉下心就没问题。不把(a1\par
b1\par
) 定为e1\par
$\prime$ ，把(a2\par
b2\par
) 定为e2\par
$\prime$ ，这个式子就能这样表示。\par
便一，这种写成粗体字的，意思就是这文字是向量。高中为止是用箭绪写，但大学里经\par
常用这种写法。 」\par
$\alpha$e1\par
$\prime$ + $\beta$e2\par
$\prime$\par
千奈「…——……虽然没什么自信，但这个感觉似乎能用内积来表示……」\par
$\alpha$e1\par
$\prime$ + $\beta$e2\par
$\prime$ = (e1\par
$\prime$ e2\par
$\prime$ ) ($\alpha$\par
$\beta$)\par
广「。。。锐利。那么，把这当中的e1\par
$\prime$, e2\par
$\prime$用向量表示写出来，会怎样呢。 」\par
千奈「………啊!变成像矩子一样了!」\par
(e1\par
$\prime$ e2\par
$\prime$ ) ($\alpha$\par
$\beta$) = (a1 a2\par
b1 b2\par
) ($\alpha$\par
$\beta$)\par
广 「答对了。 也就是说， 斜交坐标系的坐标(x$\prime$\par
y$\prime$)， 不在直角坐标系(x\par
y)里测量， 会表\par
示如下。 」\par
(x\par
y) = (a1 a2\par
b1 b2\par
) (x$\prime$\par
y$\prime$)\par
广 「便一， 把直角坐标系里的基分别定为e1: = (1\par
0) , e2: = (0\par
1)， 斜交坐标系里的基分\par
别定为e1\par
$\prime$: = (a1\par
b1\par
) , e2\par
$\prime$ := (a2\par
b2\par
)的话，下面的式子就成立。这里，: =这个符号，是「定义\par
为」的意思，不写1 := 2a，意思就是「将 1 定义为表示2a」。」\par
(e1 e2) (a1 a2\par
b1 b2\par
) = (1 0\par
0 1) (a1 a2\par
b1 b2\par
) = (a1 a2\par
b1 b2\par
) = (e1\par
$\prime$ e2\par
$\prime$ )\par
广 「这可以看作， 是基底从之前的直角坐标系， 替换到斜交坐标系时的矩子。 所以，\par
这个矩子就叫基的变换矩子。还有，这是闲话，在途中式变形里出现的矩子(1 0\par
0 1)，在\par
矩子乘积的世界里， 扮演着 「$\times$ 1」般的角色。 便一， 关于这个，下面的式子也成立。 这\par
个矩子叫单位矩子。 」\par
(1 0\par
0 1) (a b\par
c d) = (a b\par
c d) (1 0\par
0 1) = (a b\par
c d)\par
广 「单位矩子， 必定是方子， 且是对角线上放1， 其他位置放0 的矩子， 要记住。 」\par

% --- PDF page 138 ---
137\par
(1 0\par
0 1) , (\par
1 0 0\par
0 1 0\par
0 0 1\par
) , (\par
1 0 0 0\par
0 1 0 0\par
0 0 1 0\par
0 0 0 1\par
) , $\cdots$\par
千奈「是!我明白了!」\par
广「那么，接下来来想， 直角坐标系下坐标为(x, y) = ($\alpha$, $\beta$)的位置向量， 在基底替\par
换之后的斜交坐标系里测量，会变成怎样呢。 」\par
广「这次稍微有点难。为什么呢，因为斜交坐标系的基底是(a1\par
b1\par
) , (a2\par
b2\par
)，这基底是\par
在直角坐标系里测量时的东西， 至今为止的讨论都是以直角坐标系为基准来，虑的。 但\par
这次的基准是斜交坐标系，而此基底在斜交坐标下来量，自然应该是(1\par
0)和(0\par
1)。所以，\par
不不，，虑一下， 用这组向量， 那些在直角坐标系下测量时是(1\par
0)和(0\par
1)的向量， 在斜交\par
坐标系下会怎样表示，就没法把直角坐标系的坐标改写成斜交坐标系的。 」\par

% --- PDF page 139 ---
138\par
广「这个内容，不为了好懂而在直角坐标系下重新表述一下，首，该做的事就是，\par
使用(a1\par
b1\par
)和(a2\par
b2\par
)来表示(1\par
0)和(0\par
1)。千奈，试试看。 」\par
千奈「……首，，且且只持留x分量，做出 y=0 的向量。这个向量，这样做似乎就\par
能做出来了。 」\par
b2 (a1\par
b1\par
) - b1 (a2\par
b2\par
) = (a1b2\par
b1b2\par
) - (a2b1\par
b2b1\par
) = (a1b2 - a2b1\par
b1b2 - b2b1\par
) = (a1b2 - a2b1\par
0 )\par
千奈「那么，把这个向量的长度变成 1……做出来了!」\par
(1\par
0) = 1\par
a1b2 - a2b1\par
(b2 (a1\par
b1\par
) - b1 (a2\par
b2\par
))\par
广「就是那样。那么，这用斜交坐标系来表示会怎样呢。 」\par
千奈「…，把(a1\par
b1\par
)和(a2\par
b2\par
)分别改写成(1\par
0)和(0\par
1)就行了吧……」\par
1\par
a1b2 - a2b1\par
b2 (1\par
0) - b1 (0\par
1) = 1\par
a1b2 - a2b1\par
( b2\par
-b1\par
)\par
千奈「做出来了!……咦，这个是什么来着?」\par
广「这就是，把直角坐标系的基(1\par
0)用斜交坐标系的坐标测量时的坐标。。 」\par
千奈「来来如此!那接下来对(0\par
1)也做同样，虑的话……」\par
a2 (a1\par
b1\par
) - a1 (a2\par
b2\par
) = (a1a2\par
b1a2\par
) - (a2a1\par
b2a1\par
) = (a1a2 - a1a2\par
b1a2 - a1b2\par
) = ( 0\par
b1a2 - a1b2\par
)\par
(0\par
1) = 1\par
b1a2 - a1b2\par
(a2 (a1\par
b1\par
) - a1 (a2\par
b2\par
))\par

% --- PDF page 140 ---
139\par
1\par
b1a2 - a1b2\par
(a2 (1\par
0) - a1 (0\par
1)) = 1\par
b1a2 - a1b2\par
( a2\par
-a1\par
) = 1\par
a1b2 - a2b1\par
(-a2\par
a1\par
)\par
广「…，对的。由此，从斜交坐标替换到直角坐标所用的基的变换矩子，把才才求\par
出的向量排列起来就行了，具体如下。 」\par
1\par
a1b2 - a2b1\par
( b2 -a2\par
-b1 a1\par
)\par
广 P 「因为至今都没说明过，所以在此补充，矩子的常数倍，就是把该矩子的各分\par
量做常数倍。教向量的常数倍一样操作即可。 」\par
广「确实，是那样呢。……千奈，虽然唐突，但请觉求一下这个矩子，与才才那个\par
直角坐标到斜交坐标的变换矩子的乘积。 」\par
千奈「我来算算看……啊!」\par
1\par
a1b2 - a2b1\par
( b2 -a2\par
-b1 a1\par
) (a1 a2\par
b1 b2\par
) = 1\par
a1b2 - a2b1\par
( a1b2 - a2b1 a2b2 - a2b2\par
-a1b1 + a1b1 -a2b1 + a1b2\par
)\par
= (1 0\par
0 1)\par
千奈「这么说来，难道……果然!」\par
(a1 a2\par
b1 b2\par
) $\cdot$ 1\par
a1b2 - a2b1\par
( b2 -a2\par
-b1 a1\par
) = 1\par
a1b2 - a2b1\par
(a1b2 - a2b1 -a1b1 + a1b1\par
a2b2 - a2b2 -a2b1 + a1b2\par
)\par
= (1 0\par
0 1)\par
广「…，就是这么回事。就像这样，，虑反向的变换矩子，的乘才才的变换矩子，\par
就会变成单位矩子。 这个， 相对于变换矩子的反向变换的矩子， 也就是， 对某个矩子A，\par
使A-1A, AA-1都成为单位矩子的矩子A-1，就称为逆矩子。 」\par
千奈「才才的例子里，就是A = (a1 a2\par
b1 b2\par
) , AA-1 =\par
1\par
a1b2-a2b1\par
( b2 -a2\par
-b1 a1\par
)对吧!」\par
广「就是这么回事。那么，这里有提问，对某个矩子，逆矩子结是存在的 。 」\par
千奈「…——，虽然很难……啊，但是，我觉得不能那么说!」\par
广「嚯。为什么那么想?」\par
千奈「仔细看逆矩子，发现附带着系数!分母变成 0 的时候，会不会就没有逆矩子\par
了?我记得曾像念咒一样被教过，说不可以除以 0!」\par
广 「答对了。 其实， 要使逆矩子存在， 必须满足a1b2 - a2b1 $\neq$ 0。 那么， 来想想a1b2 -\par
a2b1 = 0表示什么意思吧。 」\par
广「千奈，比。，觉懂 。 」\par

% --- PDF page 141 ---
140\par
千奈「就是那个，几比几，的东西对吧?」\par
广「对。然后，觉还记得内项之积和外项之积相等 。 」\par
千奈「内项、外项……?」\par
广「就是这样子。 」\par
a: b = c: d的时候，ad = bc\par
千奈「啊!那个，把比的外侧和内侧乘起来，结果一样的那东西对吧!」\par
广「对，就是那个。这次的情况下，不a1b2 - a2b1 = 0，则a1b2 = a2b1，所以下面\par
的比就成立。 」\par
a1: a2 = b1: b2\par
广「这里，不设k := a2/a1，由于比。一『，就有a2 = ka1, b2 = kb1，因此变成下\par
面这样。 」\par
(a2\par
b2\par
) = (ka1\par
kb1\par
) = k (a1\par
b1\par
)\par
广「这个式子表明，基的其中一个是另一个的常数倍向量。也就是说，这对应于轴\par
重叠了的情况。这就说明了，它表示的是线性相关。 」\par
千奈「来来如此!不是线性相关的情况，也就是线性无关时，逆矩子就存在对吧 !」\par
广「就是那么回事。 」\par
广 P「在此补充，严格来说，a1b2 = a2b1 $\Rightarrow$ a1: a2 = b1: b2并不成立。为什么呢，\par
因为比是在各分量都不为 0 时才被定义的，而a1, a2, b1, b2之中有谁为 0 的可能性是存\par
在的。不过， 不单侧有0 分量，那么反侧也会出现0 分量，所以必定有2 个以上分量为\par
0。例如，不a1 = 0，则a2或b1为 0。a2 = 0时，基向量为( 0\par
b1\par
)和( 0\par
b2\par
)，这明无是线性相\par
关。另外，b1 = 0时，一方的基向量就变成了(0\par
0)，而(0\par
0) = 0 (a2\par
b2\par
)，所以这也是线性相\par
关。，虑b2 = 0的情况，也完全是同样讨论。因此，即一是这类例外情况，也能说是线\par
性相关，所以结之，线性相关的话逆矩子就不存在，这么说就对了。 」\par
广「那么，这里我想稍微延申一下单位矩子的意义。本来，在坐标变换里，把基底\par
交换掉就是变换矩子的意义，但这个单位矩子如果是变换矩子的话，会变成怎样呢。 」\par
千奈「这意思是从(1\par
0)和(0\par
1)替换成(1\par
0)和(0\par
1) ……」\par
广「是啊。也就是说，什么也没做。 」\par
千奈「啊，我明白才才乘以逆矩子会变成单位矩子的意义了!」\par

% --- PDF page 142 ---
141\par
千奈「从某个坐标系 1，移的到另一个坐标系 2 时的变换矩子的逆矩子，表示的是\par
从坐标系 2 到坐标系 1 的变换。 把这些矩子相乘， 也就是， ，从坐标系1 用变换矩子变\par
换到坐标系 2 的世界之后，再用其逆矩子返回到坐标系1，所以结果上，等于是什么都\par
没做!」\par
广「就是那样。这下，关于矩子的性质就大概明白了吧。 」\par
广「那么，接下来，该转到反变\par
はんへん\par
分量与共变分量的话题了。 」\par
千奈「半\par
は\par
、半片\par
はんぺん\par
2\par
……?半片的成分是鱼啊……」\par
广「不是那个反变，是相反的变化，写作反变。那么，我来解的。 」\par
广「至今为止，我们，虑的是，斜交坐标系下的向量在直角坐标系下如何表示，也\par
就是在坐标变换前后，用同样坐标。表示的两个向量，具有怎样的关系。但这次，我们\par
要，虑同一个向量，在直角坐标系与斜交坐标系下，分别如何表示。 」\par
千奈「………也就是说，觉想，虑什么?」\par
广「比方说，千奈在来点，想测量她眼中所见物体的位置。即使千奈改变了测量用\par
的坐标， 那个物体的位置也不会的。 也就是说， 物体的位置向量， 在千奈变换坐标前后，\par
应该是大小和方向相同的向量。 但是， 由于坐标的取法不同，那个向量如何表示却是会\par
变的。现在就是想，虑这个。 」\par
2\par
译者注：半片（\par
）是日本一种用鱼肉、山药等制成的白色软糕点，此处为谐音梗\par

% --- PDF page 143 ---
142\par
千奈「来来如此!」\par
广「那么，至今我们当作斜交坐标系，虑的，是那种画平行于各坐标轴的直线，来\par
查点在哪里的方法。 」\par
千奈「也就是说，就是那种把网格变形为平行四边形般的坐标对吧?」\par
广「对。首，确认一下，在这个世界里，当然能做向量的加法减法、以及乘以常数\par
倍。 」\par
广「那么，在直角坐标系下为(x, y) = ($\alpha$, $\beta$)的向量，不在斜交坐标系下测量，会变\par
成怎样呢。 」\par

% --- PDF page 144 ---
143\par
千奈「这个教才才用逆矩子思，时是一回事!如果用逆矩子，应该会变成这样!」\par
1\par
a1b2 - a2b1\par
( b2 -a2\par
-b1 a1\par
) ($\alpha$\par
$\beta$) = 1\par
a1b2 - a2b1\par
( b2$\alpha$ -a2$\beta$\par
-b1$\alpha$ a1$\beta$ ) = (x$\prime$\par
y$\prime$)\par
广「…。这里，请注意一下才才为了致出坐标。所用的矩子。这是基的变换矩子的\par
逆矩子。也就是说，坐标的。接受的是与基底相反的变换。 」\par
千奈「来来如此!这或许就是反变……?」\par
广「对。像这样，在坐标变换中，接受与基的变换相反的变换的分量，称为反变分\par
量。反过来，接受与基的变换同样变换的分量，称为共变分量。 」\par
千奈「我理解了!那么，斜交坐标系的坐标。就是反变分量对吧!」\par
广「…，就是这么回事。 」\par
广 P 「在物理学，特别是相对论中，往往不会一一表示基向量，而是只用该基向量\par
的系数来讨论问题， 因此只要意识到在坐标变换中隐藏着这样的数学步骤， 专业书就会\par
变得易读多了。 」\par
广「那么，在斜交坐标系里，其实还可以，虑另一种坐标的表示方法。那就是，把\par
想要表示的向量，通过与基作内积来表现的方法。 」\par
千奈「内积又出来了啊。说起来，内积到底有什么意义……?」\par
广 「其实， 不有两个向量u, v， 设它们所成角为$\theta$， 则关系式u $\cdot$ v = |u||v| cos $\theta$成立。 」\par
千奈「……来来如此?」\par
广 「我想， 光看算式大概不明白， 所以且且做一下几何的说明。 对其中一个向量u，\par
从垂直于另一向量v的方向打光，把u的影子映射到v上，这个操作叫投影。这个影的长\par
度是|u| cos $\theta$。然后，再乘以被投影那侧v的长度，就是内积的大小。 」\par

% --- PDF page 145 ---
144\par
千奈「来来如此，稍微有点复杂呢……」\par
广 「…。 当v的长度为 1 的时候， 还会更简单些， 它就单计变成了投影在v所示轴上\par
影子的长度。这就是教基作内积的几何意义。 」\par
千奈 「也就是说， 和大小为1 的向量的内积， 表示的是在轴上的影子长度， 对 ?」\par
广 「大『是这样。 再稍微细一点说， 影子可能形成在v的同方向， 也可能在反方向，\par
有两种情况，cos $\theta$的正负号就表示了方向。正的话就同方向，负的话就反方向。 」\par
千奈「来来如此，大『明白了!」\par
广 P 「内积的性质(u $\cdot$ v = |u||v| cos $\theta$)，可以由余弦定理简单证明。另外，余弦定\par
理本身，又可由勾股定理简单证明，有兴趣的话，请各自试试看。 」\par
广「那么，把想表示的向量，向各轴做投影。这样产生的坐标。的组，也能表示向\par
量。 」\par

% --- PDF page 146 ---
145\par
广「那么提问。首，，这组坐标。和想表示的向量，是 1 对 1 对应的 。 」\par
千奈「凭感觉似乎是那样……但是，我不知道该怎么展示。 」\par
广「如果 「教这个坐标的x$\prime$分量相同的线」与 「教这个坐标的y$\prime$分量相同的线」，它\par
们相交于一点， 那么对于坐标， 向量就会唯一确定。 这次的情况下， 垂直于x$\prime$轴的直线，\par
与垂直于y$\prime$ 轴的直线，将各被画出一条。由于两者都是直线，当然只会产生一个交点。\par
所以，坐标一定下来，向量就唯一确定了。 」\par
千奈「真是这样!」\par
广「那么，下个问题。用这种画垂直于各坐标轴的线来测量的方法，来来在直角坐\par
标系下表示为(x, y) = ($\alpha$, $\beta$)的向量，会如何变换呢。 」\par
千奈 「首，， 必须得在直角坐标系下， 用才才那种画平行于各坐标轴的直线的方法，\par
把向量表示出来……咦?」\par
千奈「之前画平行线的方法，和画垂直线的方法，向量的表示方式不是没变 !」\par
广「对。在直角坐标系下，因为坐标轴是直角相交的，所以两者是相同的表示法，\par
也就是用哪种都变成(x, y) = ($\alpha$, $\beta$)。所以，在直角坐标系里， 两种方法的坐标。，是相\par
同的。。 」\par
千奈 「来来如此， 的确， 起觉一说还真是!那么， 教斜交坐标系的基向量取内积， 投\par
影一下的话……」\par
($\alpha$\par
$\beta$) $\cdot$ (a1\par
b1\par
) = $\alpha$a1 + $\beta$b1, ($\alpha$\par
$\beta$) $\cdot$ (a2\par
b2\par
) = $\alpha$a2 + $\beta$b2\par
广「是啊。这个似乎能用这样的式子来表示。 」\par
(a1 a2\par
b1 b2\par
) ($\alpha$\par
$\beta$) = ($\alpha$a1 + $\beta$b1\par
$\alpha$a2 + $\beta$b2\par
)\par
千奈「啊，结感觉这个很有共变的感觉!虽然是转置了……」\par

% --- PDF page 147 ---
146\par
广「确实。但是在这里，有件事我们忘了。取内积的时候，不被投影那侧向量的长\par
度不是 1，影子就会伸缩。用r1 := $\sqrt{}$(a1\par
2 + b1\par
2), r2 := $\sqrt{}$(a2\par
2 + b2\par
2)的话，影子的长度会怎\par
样呢。 」\par
千奈「啊，对啊!那样的话……是这样 ?」\par
($\alpha$\par
$\beta$) $\cdot$ 1\par
r1\par
(a1\par
b1\par
) = $\alpha$a1 + $\beta$b1\par
r1\par
, ($\alpha$\par
$\beta$) $\cdot$ 1\par
r2\par
(a2\par
b2\par
) = $\alpha$a2 + $\beta$b2\par
r2\par
广「…，是那样。那么，用我们之前一直用的斜交坐标的基，x$\prime$, y$\prime$的坐标大概会变\par
成多少呢。 」\par
千奈 「…， 沿x$\prime$轴方向的基向量长度是r1，沿y$\prime$轴方向的基向量长度是r2， 所以……」\par
x$\prime$ = $\alpha$a1 + $\beta$b1\par
r1\par
2 , y$\prime$ = $\alpha$a2 + $\beta$b2\par
r2\par
2\par
千奈「啊啊，明明难得看起来共变得挺漂亮的……」\par
广「那个心情，我懂。那么，下一个问题。在用这样的坐标来表示向量时，能用才\par
才那种基向量的和，来表示向量 。 」\par
千奈「觉、觉的意思是……?」\par
广 「之前， 用斜交坐标的基e1\par
$\prime$, e2\par
$\prime$时， 我们求过x$\prime$, y$\prime$会变成什么。 想要表示的向量，\par
是具确实由下式表示了?」\par
$\alpha$a1 + $\beta$b1\par
r1\par
2 e1\par
$\prime$ + $\alpha$a2 + $\beta$b2\par
r2\par
2 e2\par
$\prime$\par
千奈「………………???结觉得有点奇怪啊?」\par
广「觉觉得什么奇怪?」\par

% --- PDF page 148 ---
147\par
千奈「用这个向量的加法，应该会的到沿着平行于各轴的刻度的坐标才对。可是，\par
这次我们是用垂直于各轴的刻度来，虑的，所以刻度的交点位置不同。 因此，这个式子\par
的向量，就和本来想表示的向量变成不同向量了!」\par
广「对。这就等于说，在眼下这种，虑坐标的方式下，是不能用这组基的。所以，\par
必须，虑另一组基。于是，就引入对偶基向量。 」\par
千奈「对、偶、基、向量?」\par
广「…。对偶基向量。对于才才的基向量e1\par
$\prime$, e2\par
$\prime$，我们把对偶基向量写为e$\prime$1, e$\prime$2。此\par
时，这样来取e$\prime$1, e$\prime$2，以满足下面条件。 」\par
e$\prime$1 $\cdot$ e1\par
$\prime$ = 1, e$\prime$1 $\cdot$ e2\par
$\prime$ = 0, e$\prime$2 $\cdot$ e1\par
$\prime$ = 0, e$\prime$2 $\cdot$ e2\par
$\prime$ = 1\par
广「首，，内积为 0，而两个向量大小都不为 0，所以cos $\theta$ = 0，也就是说，两个\par
向量所成角是直角。 由此， 不在直角坐标系测量e$\prime$1, e$\prime$2， 利用常数k1, k2， 应能表示成这\par
个样子。 」\par
e$\prime$1 = k1 ( b2\par
-a2\par
) , e$\prime$2 = k2 ( b1\par
-a1\par
)\par
广「由此，用剩下的两个条件，把 k1, k2 求出来，应该就是这样。 」\par
e$\prime$1 $\cdot$ e1\par
$\prime$ = k1 ( b2\par
-a2\par
) $\cdot$ (a1\par
b1\par
) = k1(a1b2 - a2b1) = 1 $\Leftrightarrow$ k1 = 1\par
a1b2 - a2b1\par
e$\prime$2 $\cdot$ e2\par
$\prime$ = k2 ( b1\par
-a1\par
) $\cdot$ (a2\par
b2\par
) = k2(a2b1 - a1b2) = 1 $\Leftrightarrow$ k2 = 1\par
a2b1 - a1b2\par
广「因此，对偶基向量求得如下。 」\par

% --- PDF page 149 ---
148\par
e$\prime$1 = 1\par
a1b2 - a2b1\par
( b2\par
-a2\par
) , e$\prime$2 = 1\par
a1b2 - a2b1\par
(-b1\par
a1\par
)\par
广「从直角坐标系映射到使用这对偶基向量的斜交坐标系时的、 基的变换矩子，就\par
是这样。 」\par
(e$\prime$1 e$\prime$2) = 1\par
a1b2 - a2b1\par
( b2 -b1\par
-a2 a1\par
)\par
广「这个矩子，不觉得在哪见过 。 」\par
千奈「啊!教逆矩子非常像!……咦，但是，它被转置了呢。 」\par
广「其理由之后会渐渐明白。结之，这里冒出了个逆矩子的转置。也就是说，这组\par
基自身，接受的是教之前基向量e1\par
$\prime$, e2\par
$\prime$的变换相反的变换。这是反变的。 」\par
广「那么，接下来，试着用这对偶基向量，来表示(x, y) = ($\alpha$, $\beta$)吧。 」\par
千奈「是!虽然有点累了，但我感觉到似乎只差一点了，拼毅力做式变形啦!」\par
～式子变形后\par
3\par
～\par
千奈「做出来啦～!」\par
($\alpha$\par
$\beta$) = ($\alpha$a1 + $\beta$b1)e$\prime$1 + ($\alpha$a2 + $\beta$b2)e$\prime$2\par
广「太棒了。这个分量，不觉得在哪见过 。 」\par
千奈「啊～!是那个我感觉很有共变味道的内积!!!」\par
广 「是啊。 才才， 好不容易的共变形崩掉了， 真遗憾， 不过用了这对偶基向量之后，\par
就漂亮地成了共变形。 」\par
广「其实，把至今所做的变形全部用图来表示，就是下面这样。一边回顾至今为止\par
的内容一边看这张图，我想理解会更进一步的。 」\par
3\par
作者注：不是一直读到这里的人，我想这附近的式变形是能自己独立完成的。具体而言，，虑一下\par
直角坐标系的基如何用这些对偶基向量表示，就能做出来。请务必自己试试看。\par

% --- PDF page 150 ---
149\par
千奈「清爽了～!啊——但是，还有个没法的然的谜团留着……」\par
广 「…。 从才才开始， 不知为何转置的形式会冒出来。 最后解开这个谜团， 结束吧。 」\par
广「那么，结结一下至今为止的内容，通过坐标变换，定向量v0 = ($\alpha$, $\beta$)可表示如\par
下。 这里，e1, e2是直角坐标系的基向量，e1\par
$\prime$, e2\par
$\prime$是斜交坐标系的基向量， 进而e$\prime$1, e$\prime$2是其\par
对偶基向量。请注意，用e1\par
$\prime$, e2\par
$\prime$的坐标，和用e$\prime$1, e$\prime$2的坐标，是不同的斜交坐标系。 」\par
v = $\alpha$e1 + $\beta$e2 = $\alpha$b2 - $\beta$a2\par
a1b2 - a2b1\par
e1\par
$\prime$ + -$\alpha$b1 + $\beta$a1\par
a1b2 - a2b1\par
e2\par
$\prime$ = ($\alpha$a1 + $\beta$b1)e$\prime$1 + ($\alpha$a2 + $\beta$b2)e$\prime$2\par
广「把这个，像内积那样以一于。看的方式来写，就成了这样。 」\par
v = (e1 e2) ($\alpha$\par
$\beta$) = 1\par
a1b2 - a2b1\par
(e$\prime$1 e$\prime$2) ( $\alpha$b2 - $\beta$a2\par
-$\alpha$b1 + $\beta$a1\par
) = (e$\prime$1 e$\prime$2) ($\alpha$a1 + $\beta$b1\par
$\alpha$a2 + $\beta$b2\par
)\par
广「进而，用矩子以一于。看的方式来写，就成了这样。 」\par
v = (e1 e2) ($\alpha$\par
$\beta$) = 1\par
a1b2 - a2b1\par
(e$\prime$1 e$\prime$2) ( b2 -a2\par
-b1 a1\par
) ($\alpha$\par
$\beta$) = (e$\prime$1 e$\prime$2) (a1 b1\par
a2 b2\par
) ($\alpha$\par
$\beta$)\par
广「那么，在坐标变换中，把(e1 e2) → (e1\par
$\prime$ e2\par
$\prime$ )的基的变换矩子设为T。此时，T\par
与T-1是长这样的。 」\par
T = (e1\par
$\prime$ e2\par
$\prime$ ) = (a1 a2\par
b1 b2\par
) , T-1 = 1\par
a1b2 - a2b1\par
( b2 -a2\par
-b1 a1\par
)\par
广「用这个T，就能更好看了。 」\par
v = (e1 e2) ($\alpha$\par
$\beta$) = (e1\par
$\prime$ e2\par
$\prime$ )T-1 ($\alpha$\par
$\beta$) = (e$\prime$1 e$\prime$2)TT ($\alpha$\par
$\beta$)\par

% --- PDF page 151 ---
150\par
广 「这里，T是转置矩子的符号。 那么， 到此为止， 所有项都变形得左侧是基向量，\par
右侧是坐标。了。可以看出，e1\par
$\prime$, e2\par
$\prime$坐标系下的坐标是反变分量，e$\prime$1, e$\prime$2坐标系下的坐标\par
是共变分量。 」\par
广「进而，关于基的变换矩子，也致出了下式成立。 」\par
(e1\par
$\prime$ e2\par
$\prime$ ) = (e1 e2)T, (e$\prime$1 e$\prime$2) = (e1 e2)T-1T\par
广「也就是说，对于基底，e1\par
$\prime$, e2\par
$\prime$坐标系是共变向量，e$\prime$1, e$\prime$2坐标系是反变向量。教\par
坐标的时候相反。 」\par
广「到这里的话，就差一点了。其实，虽说由计算规则立刻就能知道，但对于矩子\par
A, B，有 AB=BTAT成立，所以用上这个的话……」\par
(e1\par
$\prime$ e2\par
$\prime$ )T-1 ($\alpha$\par
$\beta$) = (e1 e2)TT-1 ($\alpha$\par
$\beta$) = (e1 e2) ($\alpha$\par
$\beta$)\par
(e$\prime$1 e$\prime$2)TT ($\alpha$\par
$\beta$) = (e1 e2)T-1T\par
TT ($\alpha$\par
$\beta$) = (e1 e2)TT-1 ($\alpha$\par
$\beta$) = (e1 e2) ($\alpha$\par
$\beta$)\par
广「漂亮地，所有结果都聚齐了。 」\par
千奈「哇啊，好漂亮!结觉得好感的啊……!」\par
广 P 「对于矩子A, B，有 AB=BTAT成立的理由是，教，虑矩子的积时一样，把矩子\par
切开来想，思路会好一点。 矩子的积里， 左侧用行切， 右侧用列切， 所以如果直接交换，\par
就会变成不同的组块， 。就不一样了。 但如果转置后再交换， 因为转置能把行与列互换，\par
所以即使左右交换，组块也能持实不变。结果就是 AB=BTAT了。有兴趣的人，自己证\par
明一下，会成为很好的锻炼。 」\par
广「这下，转置的谜团是解开了，但我通过这个谜团想说的是，反变分量与共变分\par
量的关系性。 这个式子里的所有项， 表示的向量都是同一个定向量(x, y) = ($\alpha$, $\beta$)。 不管\par
什么坐标，它们都是用反变之物与共变之物的积来表示的。也就是说， 反变之物与共变\par
之物的积所表示的量，是在坐标变换前后都不变的量。 」\par
千奈「来来如此!的确，可以这么说呢!」\par
广「…。至于为什么能说坐标变换前后不变，那是因为，正如上面式变形所看到的\par
那样，反变之物所具有的、与坐标变换相反的作用，与共变之物所具有的、与坐标变换\par
相同的作用，互相抵消了。 」\par
广「结之，反变与共变的积，在坐标变换前后是不变的。这最重要，是我今天最想\par
传达的，希望觉能把这个记在脑子里。 」\par
千奈「我明白了!」\par

% --- PDF page 152 ---
151\par
广「今天的讨论，虽然全都是 2 维的，但这个内容，对3 维，还有更多维也是成立\par
的。所以，今后觉可以放心的处理反变和共变，没问题的。 」\par
千奈「了解啦!」\par
广「那么，今天就到这里。辛苦了。 」\par
Day2 相对论是什么\par
广「制作人，起我说。 」\par
一进入视频通话，\par
泽同学又在和制作人说话。 今天仓本同学似乎也已经，进来了。\par
广「我，课题的舞蹈，已经能跳得基本不卡壳了。。 」\par
广 P「那很厉害啊。 」\par
广「结算，把两周前的课题完成了。啊啊……」\par
千奈「为什么觉在沮丧啊……?」\par
广 P「…，不必每进步一点就情…低落一次。 」\par
不管怎么说， 距离上次研讨会已经过的大概一周了。 这段时间正好是演唱会多的时\par
期，制作人科的人日程也排得很满。\par
广 P「差不多，可能是时候再提高一点训练负荷了吧。 」\par
广「…。 」\par
广 P 「那么， 我把下个舞蹈的课题视频再发给觉， 这次请在一周内把这个舞蹈练好。 」\par
广「呵呵。觉没打算让我休息呢。 」\par
广 P「哪里哪里，就算\par
泽同学太拼命，我也会行行让觉休息的，请放心。 」\par
广「再找佑芽帮我做柔韧训练吧。 」\par
千奈「诶诶!?又要做那个地狱般的柔韧训练 ?」\par
广「呵呵。虽然痛，但很快乐。 」\par
千奈「这是重症患者啊……」\par
果然，\par
泽同学这种心的，我是怎么也比不上的。哎呀，不好。不小心起闲聊起入\par
如了。\par
千 P「辛苦了。 」\par
千奈「啊，老师!好慢啊～」\par
千 P「抱歉，有点事回是晚了。 」\par
广「那么，人到齐了，开始吧。 」\par
～（研讨会开始）～\par

% --- PDF page 153 ---
152\par
广「从今天起，我来一点点说明相对论是什么。 」\par
千奈「好期待啊～!!」\par
广「在讲物理之前，千奈学过微分和积分，还有速度和加速度之类的 ?」\par
千奈「是的，虽然很难，但结算理解了!」\par
广「太好了。那么，位置x对时间做一次、两次微分，分别会变成什么?」\par
千奈「…，位置的时间微分，表示在极短时间里位置移的了多少，所以一阶微分是\par
速度v。速度的时间微分，表示在极短时间里速度变化了多少，所以二阶微分应该是加\par
速度a!」\par
广「…，没错。那今天似乎也没问题。 」\par
广「作为相对论解说的前段，首，简单说明一下物体的运的。在牛顿运的定律中，\par
用下面 3 条公理来说明运的。 」\par
1. 一切物体，只要不受外力，就会一直静止，或沿直线以恒定速度实续运的。\par
2. 物体运的的变化，平行于施加在该物体上的力，且与力的大小成比例地发生。\par
3. 当两个物体互相施加力时，那些力的大小相等，方向完全相反。\par
千奈「公理……?」\par
广「所谓公理，就是虽然无法附上理由，但被当作正确的前提。也就是说，牛顿是\par
接受这 3 条为正确的，然后说明了运的，就是这么回事。 」\par
广「第 1 条，觉可能起过叫「惯性定律」，第 2 条叫「运的定律」，第 3 条叫「作\par
用$\cdot$反作用定律」。我们来逐条仔细看看。 」\par
广「首，是第 1 条，惯性定律。比如，搬行李时，广着台车前进，要让台车加速或\par
减速，当然就需要力。那么，让台车以恒定速度跑的时候，会怎么样呢。 」\par
千奈「虽然为了持实前进方向笔直会加力，但基本上，是不会又广又拉的……」\par
广「对。速度不变的时候，不需要给台车加力。也就是说，对台车加力，仅在速度\par
改变的时候。这件事，也可以说成，物体对速度的变化有抵抗的性质。这个性质，正是\par
惯性。所以，叫惯性定律。 」\par
广 「那么， 下一个问题。 有点突然， 千奈坐电车坐在座位上的时候， 身体摇来过 。 」\par
千奈「是的，比如发车的时候，或者停车的时候，会来呢。 」\par
广「来来如此。那么，电车行驶期间呢?」\par
千奈 「虽然会稍微横向来来……但说起来， 那种发车或停车时的冲击， 不太有啊。 」\par

% --- PDF page 154 ---
153\par
广「就是那样。那么，还有一个问题。就算电车发车或停车，只要不是特别猛烈的\par
突然启的或急刹车， 人的身体就不会被甩飞。 这就是说， 电车里面的人也教电车一起加\par
速或减速。根据惯性定律，人不受到力，是无法加速或减速的。这个力，是从哪里来的\par
呢。 」\par
千奈「…，这么一说，我就不知道了……」\par
广「平时我们不太在意吧。但是，如果坐着，就应该受到来自座面的静摩擦力，站\par
着就该受到吊环的拉力和来自地板的静摩擦力。也就是说，靠着这些力，人才能承受和\par
电车一样的加速。 」\par
广 「那么， 电车在加减速时， 要是把棒应之类的东西往自己正上方投， 会怎么样呢。 」\par
千奈「这个棒应，没从任何地方受到力!也就是说，应该会在电车里移的啊。 」\par
广「…，会那样。便一，说得非常细的话，教电车一起被密闭的空气是在的的，所\par
以应会受到空气阻力。不过那不会很大，所以应还是会在电车里移的。 」\par
广「那么，加速时和减速时，从电车里面的人看，应具体会怎么的，知道 ?」\par
千奈 「首，， 根据惯性定律， 应想要持实离开手那瞬间的电车速度。 加速的情况下，\par
因为电车的速度变快，应就会被电车抛在后面。这个，如果从电车里面的人看，电车越\par
来越快，教应的速度差越来越大，所以应该会看起来是朝前进方向的反方向加速了!」\par

% --- PDF page 155 ---
154\par
千奈「另一方面，电车减速的情况下，电车会被应抛在后面。这种情况下，电车也\par
是越来越慢， 教应的速度差也渐渐变大， 所以应该会看起来是朝前进方向的同方向加速\par
了!」\par
广「完全正确。像这样，从固定在加减速的东西里面在的的人来看，看起来就像是\par
有一个教加减速相反方向的力在作用。这个力，并不是真实在作用的吧。实际上，才才\par
的例子，应从电车里面的人看来，像是加速跑的了哪里，但实际却是遵循惯性定律，什\par
么力也没受到。像这样，从加减速的人看时，所。测到的表。力，就叫惯性力。 」\par
广 「那么， 话说回来， 再继续，虑惯性。 千奈， 惯性的行度， 是因物体而异的 。 」\par
千奈「…——，结之就是，抵抗的行度，是具因物体而异，对吧。 」\par
广「对。例如，时速150km 的棒应，和时速 150km 的特急列车，哪个更难停下来\par
呢。 」\par
千奈「当然是特急列车那一方啊!」\par
广「是呢。棒应用捕手的手套就能简单停下，但特急列车就算全力刹车，也没法马\par
上停住。也就是说，越重的东西，越难改变速度。这个就叫惯性质量。 」\par

% --- PDF page 156 ---
155\par
广 「接着， 来看看第2 定律。 才才， 我好像说过， 速度改变时会加上力， 反过来说，\par
力会改变速度。这个加速或减速的程度，才才叫加速度了。加减速，也有往哪个方向加\par
减速之分， 所以加速度是向量。 这条公理是在说， 与力平行发生速度的变化。 也就是说，\par
用某个常数k，就会变成下面这样。这里，请注意，力F和加速度a都是向量。 」\par
F = ka\par
广 「那么， 关于这个k， 回想一下才才的惯性话题， 似乎不同物质不会是同样的k。\par
具体地， 就算是同样的力， 如果惯性质量大， 就越难被加速， 如果小， 就越容易被加速。\par
于是，牛顿认为惯性质量与加速度成反比，用力F、惯性质量m，如下表示了。 」\par
F = ma\par
广「这个，正是牛顿的运的方程。 」\par
千奈「啊，我好像见过这个!」\par
广「另外，牛顿还这么想了。既然惯性质量和速度定定那个物体的运的，那用下面\par
这样的向量p，不就能记述运的的特性了 。 」\par
p = mv\par
广「这个p，就叫的量。千奈，把的量用时间微分一下看看。 」\par
千奈「…，m就那样，v变成a……啊!变成了力F!」\par
dp\par
dt = d\par
dt (mv) = m dv\par
dt = ma = F\par
广「…，会那样。……但是，其实，之后就会明白，这个思，方式有一个不太妙的\par
地方。 」\par
千奈「诶!?」\par
广「嘛，那个就留作之后的乐趣吧。现在，这样就行。 」\par
广「那么，最后是第 3 定律。现在，设设物体 A 和物体 B 在互相施加力。这条公\par
理想说的，用算式来说，就是，设设物体A 对物体 B 施加力FA，物体 B 对物体 A 施加\par
力FB，那么在这两个之间，FB = -FA成立。 」\par
广 「这里， 把物体A、B 的的量各自设为pA, pB。 此时， 由才才的讨论， 下式成立。 」\par
dpA\par
dt = FA, dpB\par
dt = FB\par
广「那么，把pA + pB用时间微分的话，会怎样呢。 」\par
千奈「我算算看……咦咦咦?」\par

% --- PDF page 157 ---
156\par
d\par
dt (pA + pB) = dpA\par
dt + dpB\par
dt = FA + FB = 0\par
广「对，变成零向量。也就是，的量的和，不着时间而变。像这样，不着时间而变\par
的东西，也就是时间微分后变成0 的东西，叫做守恒量。这次的情况下，的量的结和就\par
是守恒量。这也叫，的量的和被守恒。 」\par
广「到此为止，力学基上部分的牛顿运的定律就说明完了，没问题 ?」\par
千奈「没问题!」\par
广「太好了。那么，就渐渐进入相对论的话题吧。首，从坐标的话题开始。 」\par
广「才才的牛顿运的定律的第1 条，惯性定律成立的坐标系，叫做惯性坐标系。也\par
就是说，在惯性坐标系里，不管把静止的物体放在哪里，它都会持实静止，就算给那个\par
物体速度，它也会一直以那个速度实续运的。 」\par
广「那么，提问千奈。要是，虑地应上的坐标，那是惯性坐标系 。 」\par
千奈「不可能让应在空中静止着浮住不的呢。它必定会便着重力掉下来啊。 」\par
广「对。也就是说，地应上的坐标是非惯性坐标系。更准确地说，在重力起作用的\par
地方取坐标的话，全都会变成非惯性坐标系。严格来说，不管在宙空空间的哪里，都会\par
从某处的天体感受到重力，所以，其实现实世界里，完全的惯性坐标系是不存在的。 」\par
千奈「…，但结觉得有点疙瘩啊。如果对于某个坐标系，承认有重力的话，我感觉\par
它似乎能满足惯性定律啊。 」\par
广「来来如此。那么，来想想为什么有重力的话，就不能被承认为惯性坐标系吧。\par
现在， ，设设我们能虚假地的到一个没有重力的场所，在那个无重力的空间里，千奈乘\par
坐着火箭。 首，， 当火箭停住的时候， 不管怎么想， 千奈的坐标都是惯性坐标系。 那么，\par
当火箭以恒定速度运的时，和加减速时，各自的千奈坐标会怎样呢。 」\par

% --- PDF page 158 ---
157\par
千奈「以恒定速度运的时，我和我周围的物体，都会以和火箭相同的速度运的吧。\par
就算有和我速度不同的物体，因为力不起作用，我和那个物体的速度差也会被持实吧。\par
也就是说，如果火箭是恒定速度，就是惯性坐标系啊。 」\par
广「…。会那样。 」\par
千奈「接着，加减速时，像才才的电车例子那样，物体明明没受实际的力，却会有\par
惯性力作用啊。也就是说，这是非惯性坐标系。 」\par
广「就是那样。那么，设设载着千奈的火箭，朝着千奈绪的正上方的方向，一边加\par
速一边飞的。这时候，千奈周围的物体会怎么行的呢。 」\par
千奈 「因为会有和加速方向相反的惯性力作用， 所以物体会受到朝脚底方向的力呢。 」\par
广「…。这个，不觉得像什么 。 」\par
千奈「啊!教重力好像!」\par
广「对。设设，我们调节火箭加速度的大小，使惯性力变得和地应上的重力大小相\par
同。这时候，如果千奈一手拿着相机，另一手拿着小物体，将在空中抓着物体的手放开\par
时的物体样子录了像。这时候的影像， 是千奈在地应上展示的火箭模型里面，在静止态\par
的下做的实验呢，还是在加速中的真正火箭里面做的实验呢，将无法分辨。像这样，惯\par
性力和重力，通过加速运的，是无法区分的。这个就叫等效来理。 」\par
千奈「来来如此，的确是这样呢!」\par

% --- PDF page 159 ---
158\par
广 「那么， 下一个问题。 千奈， 有一天觉成了宙航员偶像， 要在宙空空间站里停留。\par
宙空空间站里面，当然，是无重力态的。这是为什么呢。 」\par
千奈「的确，仔细想想，不知道是为什么呢……」\par
广「那么，来解的一下这个吧。千奈，坐过咖啡杯 ?」\par
千奈「坐过!但是，转得太厉害的话，感觉要被甩飞，好可的……」\par
广「才才，觉说了什么?」\par
千奈「诶?转得太厉害的话，好可的……」\par
广「为什么会觉得可的呢。」\par
千奈「那、那当然，是因为感觉要被甩飞……」\par
广「对。也就是说，滴溜溜一转，就会在远离转转中心的方向上感觉到力。这个，\par
叫离心力。这个离心力，其实不是实际的力，而是和惯性力一样的表。力。 」\par
千奈「来来如此。但是，转转的运的，就算速度一样，运的的方向也会变，所以根\par
据惯性定律，要是没有力作用就奇怪了!」\par
广「。，真敏锐。确实，在各点，虑速度向量的话，就算向量的大小一样，方向也\par
不同，所以似乎确实有力在作用。于是，，虑像下面这样的例子。 」\par

% --- PDF page 160 ---
159\par
广「例如，，虑用水桶装了水，用肩膀抡着甩的时候。这时，水桶里的水不会掉下\par
来，但这件事经常被说成来因是离心力。但是，正如才才所说，离心力不是实际的力，\par
是表。的力。 」\par
广「那么提问千奈。如果放开在抡水桶的手，会怎么样呢。 」\par
千奈「那个，当然是水桶会飞到别处的啊。 」\par
广「是呢。也就是说，让水桶做圆周运的的，可以说是紧紧握着水桶提梁的握力。\par
水桶虽然想的外面， 但因为手指把水桶朝手臂的方向拉近， 所以水桶才能画出和手臂相\par
同的轨迹。 」\par
千奈「的确，是那样呢!」\par
广「手臂的方向，也就是说，是朝向圆周运的中心的方向，所以是这个方向的力在\par
作用， 构成了圆周运的。 这个， 朝向中心的力， 叫向心力。 向心力， 是实际在作用的力。 」\par
千奈 「来来如此， 我明白了!可是， 这时候水桶里的水， 并没有和手臂等连接在一起，\par
为什么却会和水桶一起运的呢?」\par
广「那是因为，它从水桶受到着法向力。 」\par
千奈「法向……是什么呢?」\par

% --- PDF page 161 ---
160\par
广「法向力。例如，我们在地应上结是受着重力，但就算站着或坐着，也不会沉下\par
的。这，只要认为是从地板或椅子，受到了教重力大小相等、方向相反的力，就能说得\par
通。这个力，叫支实力。这类支实力，会垂直于接触着的面作用。因此，从接触面感觉\par
到的支实力，也特别地叫做法向力。 」\par
广 P 「严格来讲，支实力有摩擦力和法向力两种。摩擦力是平行于面方向作用的支\par
实力，法向力是垂直于面方向作用的支实力。法向力，可以看作一种弹性力。弹性力，\par
是物体变形时，为了恢复来态而作用的力。就算是对坚硬物体施加小力的情况，在来子\par
层级看也是变形了，为了恢复来态，弹性力在作用。这样想，就能说明法向力了。 」\par
广「回到水桶的话题，现在，从水桶底部作用的法向力，应该和手臂的方向一『。\par
也就是说，这个法向力，在引起水的圆周运的。 」\par

% --- PDF page 162 ---
161\par
广「对于处于和水同样立场的人来说，由于离心力作用，在桶底的方向上受到力，\par
所以看起来，就仿佛在向桶底落下一般。另一方面，一旦接触桶底，就会从桶受到法向\par
力。这个法向力和离心力，是平衡的。这，教在地面上站着的人没有区别。在地应上，\par
重力作用在人身上，站着的时候，从地面受到同样大小的法向力，由此静止。所以，水\par
也看起来像是自立在桶底一样。结算是形成印象了 。 」\par
广「那么，基于此前的讨论，绕着地应周围回转的宙空空间站内部，为什么会变成\par
无重力，觉明白了 。 」\par
千奈 「是的!宙空空间站以地应为中心， 正以重力为向心力做圆周运的。 在宙空空间\par
站内部，从和空间站一起绕地应上回转的人来看，空间站内，除了重力，还有表。的力\par
离心力在作用。 因为这些平衡了， 所以从内部看， 空间站内看起来就像无重力一样呢!」\par
广「就是那样。其实，宙空空间站的人们，也还是受着重力的。 」\par
广「像这样，就算存在重力的情况下，通过用惯性力或离心力这种，因坐标系而产\par
生的表。力， 也能够，虑如同不存在重力的坐标系。这样构成的坐标系，就叫局域惯性\par
系。 」\par

% --- PDF page 163 ---
162\par
广 P 「便一，因为离心力也是因坐标系而产生的表。力，所以从这点来看，在无重\par
力空间里，虑转转坐标系的情况，这也会成为非惯性系。 经常被误解的是， 既然匀速运\par
的就是惯性系，那匀速转转运的也是惯性系，但这完全是错的，请注意。 」\par
广「那么，大『明白惯性坐标系是什么东西了，接下来，讲讲伽利略变换。 」\par
千奈「Galilei，是伽利略$\cdot$伽利雷 ?」\par
广「对。伽利略认为，一切物理定律，都必须能用对一切。测者相同形式的算式来\par
写。 另外反过来， 他也认为， 不管从哪一。测者看， 都能用同样形式的算式来写的定律，\par
才是物理定律。这个，叫相对性来理。 」\par
广「进而，他认为运的终究只是对。测者而言的运的，并不存在运的的绝对基准，\par
也就是绝对运的。 」\par
千奈「绝对运的……?」\par
广「例如，才才的惯性系里，设设千奈乘坐着火箭，火箭外面，除了一个小应，什\par
么都没有，是一片计白的世界。这时候，就算火箭外面的应从千奈看来在的，也没法判\par
别哪一方是停着的，哪一方在的，对吧。搞不好，两者都在的的可能性也有。 」\par
千奈「可是，至少我自己在的这点，结该能知道吧。 」\par
广「来来如此。但是，有能证明自己在的的证据 。 」\par
千奈「那、那么一说，确实……」\par
广「周围的风景在的，不能成为自己在的的证据。因为，有可能是风景自身在的。\par
只要除此之外没有自己在的的证据，就连自己是不是在的，都搞不清楚。所以，运的这\par
东西，终究只能讨论从那个人看来如何如何的话题了。 」\par
千奈「来来如此……虽然有点难，但结算是大『形成印象了。 」\par
广 「现在， 设设除了千奈， 我和佑芽也在惯性坐标系里。 现在， 千奈和佑芽两个人，\par
乘坐着分别的火箭，从我看来，两人的火箭各自在以v千奈, v佑芽前进。那么，这时候千\par
奈看佑芽的火箭，速度会怎样呢。 」\par

% --- PDF page 164 ---
163\par
千奈「…，只要求从我看花海同学的相对速度就可以了吧。这样一来，从对方减的\par
自己就行，所以从我看来，似乎是以v佑芽 - v千奈在前进啊。 」\par
广「…，答对了。那么，设v := v佑芽 - v千奈 ，把x, y, z 的成分各自设为vx, vy, vz 。\par
对千奈而言的坐标系(x, y, z, t)和对佑芽而言的坐标系(x$\prime$, y$\prime$, z$\prime$, t$\prime$)之间的变换， 会怎样呢。\par
这里， 设(x, y, z)和(x$\prime$, y$\prime$, z$\prime$)都是直角坐标系。 便一， 这里设设两人从同一场所发射了火\par
箭， 坐标以两人出发的时间点为来点。也就是说， 设(x, y, z, t) = (x$\prime$, y$\prime$, z$\prime$, t$\prime$) = (0, 0, 0, 0)\par
成立。此外，设x和x$\prime$轴、y和y$\prime$轴、z和z$\prime$轴各自朝向同一方向。 」\par
千奈「t，是时间 ?」\par
广「对。t是对千奈而言的时间，t$\prime$是对佑芽而言的时间。时间来点，是指出发的瞬\par
间，这里，是设t = t$\prime$ = 0的瞬间为同时出发的瞬间，这么个意思。 」\par
千奈「为什么时间突然冒出来了……?」\par
广「那是因为，两人的位置因时间而变。不用时间的变数，就没法表示坐标。 」\par
千奈「来来如此，是这么回事，才突然冒出了时间啊。时间我觉得，应该不会受到\par
变换……」\par
t$\prime$ = t\par
千奈 「然后， 空间的坐标，t 秒后花海同学前进了vt， 所以对花海同学而言的来点，\par
是在对我而言的vt位置上。也就是说……」\par
(\par
x$\prime$\par
y$\prime$\par
z$\prime$\par
) = (\par
x\par
y\par
z\par
) - vt = (\par
x - vxt\par
y - vyt\par
z - vzt\par
)\par
广「是那样呢。 结结起来， 千奈的坐标系(x, y, z, t)， 相对于它以速度v = (vx, vy, vz)\par
运的的佑芽的坐标系的关系，是这样的。 」\par
(\par
x$\prime$\par
y$\prime$\par
z$\prime$\par
t$\prime$\par
) = (\par
x - vxt\par
y - vyt\par
z - vzt\par
t\par
) = (\par
1 0 0 -vx\par
0 1 0 -vy\par
0 0 1 -vz\par
0 0 0 1\par
) (\par
x\par
y\par
z\par
t\par
)\par

% --- PDF page 165 ---
164\par
广「这个变换，就是伽利略变换。那么，来看看这个变换的性质吧。这次，我又乘\par
上教两人不同的火箭， 从和两人一起的出发地点， 我也发射自己的火箭。 设设从佑芽看，\par
我在以速度w运的。 设我的坐标系为(x$\prime$$\prime$, y$\prime$$\prime$, z$\prime$$\prime$, t$\prime$$\prime$)， 按才才的要领， 对齐坐标轴和来点。\par
这时候，从千奈看，我必须是以速度v + w在运的才行，但真的是那样 。 」\par
千奈「首，，花海同学和\par
泽同学的坐标关系，是这样呢……」\par
(\par
x$\prime$$\prime$\par
y$\prime$$\prime$\par
z$\prime$$\prime$\par
t$\prime$$\prime$\par
) = (\par
1 0 0 -wx\par
0 1 0 -wy\par
0 0 1 -wz\par
0 0 0 1\par
) (\par
x$\prime$\par
y$\prime$\par
z$\prime$\par
t$\prime$\par
)\par
千奈「把才才的式子代入这个……」\par
(\par
x$\prime$$\prime$\par
y$\prime$$\prime$\par
z$\prime$$\prime$\par
t$\prime$$\prime$\par
) = (\par
1 0 0 -wx\par
0 1 0 -wy\par
0 0 1 -wz\par
0 0 0 1\par
) (\par
1 0 0 -vx\par
0 1 0 -vy\par
0 0 1 -vz\par
0 0 0 1\par
) (\par
x\par
y\par
z\par
t\par
) = (\par
1 0 0 -vx - wx\par
0 1 0 -vy - wy\par
0 0 1 -vz - wz\par
0 0 0 1\par
) (\par
x\par
y\par
z\par
t\par
)\par
千奈 「啊!这和设设从我看\par
泽同学以v + w运的时的伽利略变换的式子， 是一『的\par
啊!」\par
广「就是那样。由此，明白了伽利略变换能很好地成立。 」\par
千奈「那么，就是可喜可贺，皆大欢喜了啊!」\par
广「……不，其实，事的没那么简单。 」\par
千奈「……诶?」\par

% --- PDF page 166 ---
165\par
广「那么，小测验。某一天，佑芽说『一起朝着夕阳跑吧 !』 ，于是我和佑芽一起，\par
朝着太阳，沿着海岸，往同方向跑起来。我马上就累趴了，于是乘上了路边就有的电车\par
的第一节车厢，佑芽则就那样继续沿路跑下的。当然，阳光都有照到两人，但从两人来\par
看，各自太阳的光速，哪边更快呢。 」\par
千奈「诶，诶诶?那个，当然应该是乘着电车的\par
泽同学那一方不是 ……?」\par
广「……其实，会变得完全相同。 」\par
千奈「……诶、诶诶诶诶诶诶诶诶!?」\par
广「更进一步说，在这个计白的惯性坐标系空间里，站着不的看太阳的人，和乘着\par
火箭朝太阳以极猛的速度靠近太阳的人，无论对哪一个人，光都以同样的速度前进。像\par
这样， 光不管从哪一惯性坐标系看， 都以某个恒定的速度前进。 这个， 叫光速不变来理。\par
说是「来理」，和公理是同个意思，是把它当作正确来接受的这么个东西。 」\par
千奈「可、可是，那种事怎么知道的?」\par
广 「其实， 过的有做过实验的人。 是叫迈克尔逊-莫雷实验，是作为那个实验结果，\par
而知道的事情。实验内容深挖起来很难，所以这次只采用结果，结之，就是知道了不管\par
从谁看，光速都是恒定的。 」\par
广「但是，这样的话，伽利略变换就没法用了。为什么呢，因为不同的人会。测到\par
光的不同速度。所以，结果我们必须，虑新的变换。必须，虑什么样的变换呢，就是将\par
光速持实不变的变换。 」\par
千奈「好像很难啊……」\par
广「接下来我就想求这个变换，洛伦兹变换。但这个变换，要真正理解，必须扎实\par
理解一个叫电学学的领域的基上方程组，麦克斯韦方程组。」\par
千奈「电、电学……?基上……?麦克斯……?呜哇，已经全是起不懂的词了～!!」\par

% --- PDF page 167 ---
166\par
广「放心。接下来要展示的，是不需要麦克斯韦方程组的方法的广致。虽然需要一\par
点发散想象力，但我想现在的千奈也能理解。 」\par
千奈「太、太好了……」\par
广 「那么， 把光速不变来理放在心上， 再次求一下才才态况下， 千奈的坐标(x, y, z, t)\par
和佑芽的坐标(x$\prime$, y$\prime$, z$\prime$, t$\prime$)的关系。 坐标轴也取平行， 设t = t$\prime$ = 0为发射的瞬间， 出发位\par
置也对齐。只是这次，直接对v处理似乎很难。于是，最初就把x, x$\prime$轴的正方向取在v的\par
方向上。这样一来，v就成了只有x分量的向量，变得简单了。 」\par
千奈「从现在起，必须要解这个4 $\times$ 4矩子 ……呜，好难……」\par
(\par
x$\prime$\par
y$\prime$\par
z$\prime$\par
t$\prime$\par
) = (\par
a11 a12 a13 a14\par
a21 a22 a23 a24\par
a31 a32 a33 a34\par
a41 a42 a43 a44\par
) (\par
x\par
y\par
z\par
t\par
)\par
广「这个很难呢。所以，在计算之前，，想想有没有什么能定下来的东西。 」\par
广 「首，， 第1 行。 由定义，x$\prime$ = 0的位置， 从千奈看来， 在千奈的时刻过了t之后，\par
必须移的到了x = vt才行。所以，设k为常数，则必须取x$\prime$ = k(x - vt)的形式。在这一\par
点上，如果y, z参与进来了， 似乎就会出问题。因为如果他们参与进来，x$\prime$ = 0的位置就\par
会因y, z的。而改变了。由此可知a12 = a13 = 0。」\par

% --- PDF page 168 ---
167\par
广「接着，第 2 行。首，，如果y依赖于t，彼此的来点就会在y, y$\prime$方向移的，这违\par
反沿x, x$\prime$轴做匀速运的的设定。所以a24 = 0。于是，x, y, z的系数会持留下来。如果x, z\par
的系数留下来，就意味着y轴与y$\prime$轴之间会与时间无关地发生倾斜， 所以a21 = a23 = 0。\par
最后，y的系数取方向一『时必须是正的。另外，如果系数不为 1，y轴与y$\prime$轴的倍率就\par
会改变。 这是因为， 从佑芽。测千奈时， 千奈应该以-v运的， 但此时的变换也会着之改\par
变倍率，致『对于千奈来说，在具有a22倍刻度的佑芽的坐标轴上测量千奈的坐标轴时\par
会得到a22的长度，这就会变得奇怪。因此可知a22 = 1。第三行同理，可得a31 = a32 =\par
a34 = 0，a33 = 1。」\par
广「最后，只有第 4 行，很遗憾，这个还不太清楚，所以持留来样。 」\par
千奈「诶!?时间，难道不是不受变换 !?」\par
广「不。 时间， 可能是因人而异的。 这次能用的， 只有相对性来理和光速不变来理。\par
这些里面，并不包含「时间对任何人都是一样的」。有一瞬间，觉可能会觉得，从相对\par
性来理似乎可以主张， 时间应当对任何。测者都以同样方式流逝。 但相对性来理所说的，\par
终究只是物理定律应当能以同样形式的算式来写， 仅此而已， 而时间是为了记述物理定\par
律的东西，并非物理定律本身，所以，因人而异也是可以的。 」\par
千奈「可、可是，比如，设设在某处的墙上挂着时钟，那个时钟的指针，不管从谁\par
看，都必须同样地的不是 ?」\par
广「那种情况下的相对性来理， 是希望记述指针运的的算式，对各个。测者也都是\par
同样形式，仅此而已。绝不是，希望时钟的指针对。测者以同样速度的，或希望看起来\par
是同样的运的。 」\par
千奈「来来如此……?」\par
广「嘛，就当被我骗了，试试变换看看。 」\par
千奈「虽然难以接受，但我明白了……」\par
广「那么，到这里已经变得相当简单了，再重写一次矩子。 」\par
(\par
x$\prime$\par
y$\prime$\par
z$\prime$\par
t$\prime$\par
) = (\par
k 0 0 -kv\par
0 1 0 0\par
0 0 1 0\par
a41 a42 a43 a44\par
) (\par
x\par
y\par
z\par
t\par
)\par
广「那么。到这里了，接下来使用光速不变来理。千奈，应面的方程式，已经学过\par
了 。 」\par
千奈「是的 !中心是 (x, y, z) = (a, b, c) ，半径为r 的应面的方程式，是 (x - a)2 +\par
(y - b)2 + (z - c)2 = r2对吧!」\par

% --- PDF page 169 ---
168\par
广「就是那样。这个应面方程式很重要。现在，设设千奈和佑芽各自在手边拿着灯\par
泡。灯泡的光，设光的速度，即光速为c ，那么在千奈的时间过了t 之后，就会前进ct 。\par
由此，千奈手边灯泡的光，在时刻t到达了哪里呢?，虑一下，就成了下式这样，以千奈\par
的位置(x, y, z) = (0,0,0)为中心，半径ct的应面。 」\par
x2 + y2 + z2 = (ct)2\par
广「同样，要是也以在佑芽手边的灯泡来，虑，下式成立。这里，请注意，从千奈\par
看来，佑芽或许在的，但从佑芽看来，自己没在的，是千奈在的的样子，所以是和千奈\par
完全相同的态况。 」\par
x$\prime$2 + y$\prime$2 + z$\prime$2 = (ct$\prime$)2\par
广「由此，(ct)2 - x2 - y2 - z2 = (ct$\prime$)2 - x$\prime$2 - y$\prime$2 - z$\prime$2 = 0成立， 所以代入y$\prime$ =\par
y, z$\prime$ = z的话，下式就成立。 」\par
(ct)2 - x2 = (ct$\prime$)2 - x$\prime$2\par
广 「进而， 代入x$\prime$ = kx - kvt, t$\prime$ = a41x + a42y + a43z + a44t的话， 就成了这样。 」\par
(ct)2 - x2 = c2(a41x + a42y + a43z + a44t)2 - (kx - kvt)2\par
广「比好左边和右边，因为右边才有的y, z在左边没有，所以a42 = a43 = 0。代入\par
这个，继续广进计算，就成了下面这样。 」\par
(ct)2 - x2 = c2a41\par
2 x2 + 2c2a41a44xt + c2a44\par
2 t2 - k2x2 + 2k2vxt - k2v2t2\par
广「我希望这个对任何x, t都成立，所以比好左边和右边的系数，就成了这样。 」\par
\{\par
c2a41\par
2 - k2 = -1\par
2c2a41a44 + 2k2v = 0\par
c2a44\par
2 - k2v2 = c2\par
广「使用第 2 个式子的话，就成了这样。 」\par
k2 = - c2\par
v a41a44\par
广「将这个代入第 1 和第 3 个式子，整理后，就成了这样。 」\par
\{va41\par
2 + a41a44 = -v/c2\par
a44\par
2 + va41a44 = 1\par
广「将这两个式子变形的话，就成了这样。 」\par

% --- PDF page 170 ---
169\par
\{a41(a44 + va41) = -v/c2\par
a44(a44 + va41) = 1\par
广「由此可知a41, a44都不为 0，故下式成立。 」\par
a41 = - v\par
c2 a44\par
广「使用这个式子，a41\par
2 , a44\par
2 就如下求得。 」\par
\{\par
a41 (- c2\par
v a41 + va41) = -v/c2 $\Leftrightarrow$ a41\par
2 = v2/c4\par
1 - (v2/c2)\par
a44 (a44 - v2\par
c2 a44) = 1 $\Leftrightarrow$ a44\par
2 = 1\par
1 - (v2/c2)\par
广「这里，当v = 0时，对千奈而言的x, t与对佑芽而言的x$\prime$, t$\prime$必须一『。由此，必\par
须存在k > 0, a44 > 0，所以，虑这个，k, a41, a44就像下面这样求得。 」\par
k = a44 = 1\par
$\sqrt{}$1 - (v2/c2)\par
, a41 = - v/c2\par
$\sqrt{}$1 - (v2/c2)\par
广「这样，洛伦兹变换就求出来了。辛苦了。 」\par
(\par
x$\prime$\par
y$\prime$\par
z$\prime$\par
t$\prime$\par
) =\par
(\par
1\par
$\sqrt{}$1 - (v2/c2)\par
0 0 - v\par
$\sqrt{}$1 - (v2/c2)\par
0 1 0 0\par
0 0 1 0\par
- v/c2\par
$\sqrt{}$1 - (v2/c2)\par
0 0 1\par
$\sqrt{}$1 - (v2/c2) )\par
(\par
x\par
y\par
z\par
t\par
)\par
千奈「好长啊……」\par
广「因为知道了y, z, y$\prime$, z$\prime$完全无关变换，所以如果选择不无示这些分量，就成了这\par
样。 」\par
(x$\prime$\par
t$\prime$)=(\par
1\par
$\sqrt{}$1-(v2/c2) -\par
v\par
$\sqrt{}$1-(v2/c2)\par
-\par
v/c2\par
$\sqrt{}$1-(v2/c2)\par
1\par
$\sqrt{}$1-(v2/c2)\par
) (x\par
t)\par
广「进而，设$\beta$ := v/c, $\gamma$ := 1/$\sqrt{}$(1 - $\beta$2)来改写，就成了这样。 」\par

% --- PDF page 171 ---
170\par
(x$\prime$\par
t$\prime$)=( $\gamma$ -$\gamma$v\par
-$\gamma$$\beta$/c $\gamma$ ) (x\par
t)\par
广「进而，把t, t$\prime$分别作ct, ct$\prime$的话，就成了这样。 」\par
( x$\prime$\par
ct$\prime$)=( $\gamma$ -$\gamma$$\beta$\par
-$\gamma$$\beta$ $\gamma$ ) ( x\par
ct)\par
广「这样，外。大『变好起来了。 」\par
千奈「变得超级简单了啊～!」\par
广 P 「虽然这里没有介绍，但使用麦克斯韦方程组也能致出洛伦兹变换。便一，学\par
电学学下的，我想会出现洛伦兹条件啦洛伦兹规范之类的东西。请留意不要把这个\par
Lorenz 写成 Lorentz\par
4\par
，是不同的人。另外，其实伽利略变换的具定，也可以通过对麦克\par
斯韦方程组做伽利略变换来致出。在相对性来理中，要求从任何。测者的角度来看，物\par
理定律的表达形式都应持实一『， 但通过麦克斯韦方程组的伽利略变换， 我们可以看出\par
这一点并不成立。」\par
广「最后，来，虑一下一般以v = (vx, vy, vz)表示时的洛伦兹变换。以下，设v的分\par
量中有 2 个以上不为 0，设v = |v| = $\sqrt{}$vx2 + vy2 + vz2。千奈和佑芽的坐标系设定相同，\par
从千奈看来佑芽的速度是偏离了x, x$\prime$轴的感觉。」\par
千奈「最初的条件，「2 个以上不为 0」是……」\par
广「首，，如果有 2 个为 0，v就会和某处的轴的方向一『。这种情况，只要把才\par
才的变换，从x换成那个轴的分量来用就行，很无趣。另外，如果 3 个都为 0，这，就\par
只是千奈和佑芽一直在同一地方静止不的而已， 这也无趣……不如说， 没有，虑的意义。\par
所以，存在这个设定。 」\par
千奈「来来如此。在各方向都有速度的时候，该怎么做才好呢……」\par
广「只要x, x$\prime$轴朝着速度向量的方向，就能用才才求出的洛伦兹变换的式子吧。于\par
是，首，，从千奈看到v = (vx, vy, vz)的来来设定的坐标系，，虑要怎么移的，才能使x\par
轴和v的方向一『。这，是不改变来点，改变各轴方向，也就是摆弄基底的操作，所以\par
似乎能用3 $\times$ 3矩子表示。 」\par
千奈「的确……」\par
4\par
译者注：译者这里全部译为了中文人名，此段前文洛伦兹条件和洛伦兹规范均是 Lorenz，除此之外\par
都是 Lorentz\par

% --- PDF page 172 ---
171\par
广 「那么， 现在， 将千奈来来的坐标系设为(x, y, z)， 将通过转转使x轴与v方向一『\par
而重新设定的坐标系设为(X, Y, Z)。 稍微想想， 只要X轴教v同方向就行，Y, Z轴， 只要作\par
为直角坐标系能设定， 怎么设定都可以呢。 就是说， 残留着以X轴为中心的转转自由度。\par
把这个放在心上，，虑基底变换的矩子的话，会怎么样呢。 」\par
千奈「首，，(\par
1\par
0\par
0\par
)会移到\par
1\par
v (\par
vx\par
vy\par
vz\par
)吧。Y, Z轴的基，该怎么法呢……」\par
广「，，适当定定其中某一个吧。…，比如说，虽然有点行行……」\par
(\par
vx\par
vy\par
vz\par
) $\cdot$ (\par
vy\par
-vx\par
0\par
) = vxvy - vyvx + 0 = 0\par
广「这个能用作Y轴的基。 」\par
千奈「来来如此，那么，(\par
0\par
1\par
0\par
)就移到\par
1\par
$\sqrt{}$vx2+vy2\par
(\par
vy\par
-vx\par
0\par
)好了。那么，Z 轴的基看来能定\par
下来……但是，我不知道该怎么计算才好～!」\par
广「是呢。于是，「外积」登场。这个，作为符号，使用$\times$，像下面这样计算。 」\par
(\par
a1\par
a2\par
a3\par
) $\times$ (\par
b1\par
b2\par
b3\par
) = (\par
a2b3 - a3b2\par
a3b1 - a1b3\par
a1b2 - a2b1\par
)\par
广 「作为这个外积， 在左边放X轴的向量， 右边放Y轴的向量来计算的话， 得到的向\par
量，就是能用于Z轴基的向量。 」\par
千奈「嘿……虽然不知道为啥，但我明白该做什么了。 」\par

% --- PDF page 173 ---
172\par
广 P 「对外积的性质做更一点说明的话，一般，对3 维实向量u, v，u $\times$ v会成为与\par
u和v所张平面垂直的法线向量的方向。 这里， 法线向量一定会张在右手系的方向。 怎么\par
判别这个方向呢，很简单，用右手，像弗莱明右手定则那样，把拇指、食指、中指摆成\par
互相直角，把拇指对向u的方向，食指对向v的方向，中指所指的方向就是u $\times$ v的方向\par
了。 」\par
广 P 「这里，如果把拇指便着v，食指便着u，中指所指的方向就和才才正相反。其\par
实，外积里也有v $\times$ u = -u $\times$ v。。另外，不u, v为线性相关，外积就为0。这个，各自\par
计算一下就会明白的。另外，在数学上，坐标系的右手系、左手系，指的是各坐标轴的\par
规范正交基向量\{ex, ey, ez\}满足下面这样的关系，也一并掌握吧。 」\par
ex $\times$ ey = ez, ey $\times$ ez = ex, ez $\times$ ex = ey→右手系\par
ex $\times$ ey = -ez, ey $\times$ ez = -ex, ez $\times$ ex = -ey→左手系\par
千奈「那么，计算才才向量的外积……」\par
(\par
vx\par
vy\par
vz\par
) $\times$ (\par
vy\par
-vx\par
0\par
) = (\par
0 + vzvx\par
vzvy - 0\par
-vx2 - vy2\par
) = (\par
vzvx\par
vzvy\par
-vx2 - vy2\par
)\par
千奈「…，好奇怪的式子啊。 」\par
广「嘛，没法法。那么，用这个当Z轴的基。 」\par
千奈「…，那么，首，才才向量的长度是……」\par
$\sqrt{}$(vzvx)2 + (vzvy)\par
2\par
+ (-vx2 - vy2)\par
2\par
= $\sqrt{}$vz2(vx2 + vy2) + (vx2 + vy2)\par
2\par
千奈「…，想弄得更干净些……」\par
广「因为(vx2 + vy2)\par
2\par
= (vx2 + vy2)(vx2 + vy2)，所以能逆用分配律。。 」\par
千奈「真的。那么……」\par

% --- PDF page 174 ---
173\par
$\sqrt{}$vz2(vx2 + vy2) + (vx2 + vy2)\par
2\par
= $\sqrt{}$(vx2 + vy2 + vz2)(vx2 + vy2) = v$\sqrt{}$vx2 + vy2\par
千奈「某种程度变干净了。由此，(\par
0\par
0\par
1\par
)会移到\par
1\par
v$\sqrt{}$vx2+vy2\par
(\par
vzvx\par
vzvy\par
-vx2 - vy2\par
)呢。那么，基的\par
变换矩子是……」\par
(\par
vx\par
v\par
vy\par
$\sqrt{}$vx2 + vy2\par
vzvx\par
v$\sqrt{}$vx2 + vy2\par
vy\par
v - vx\par
$\sqrt{}$vx2 + vy2\par
vzvy\par
v$\sqrt{}$vx2 + vy2\par
vz\par
v 0 - vx2 + vy2\par
v$\sqrt{}$vx2 + vy2\par
)\par
千奈「呜哇，奇怪的形态!」\par
广「这样，基的变换矩子就求出来了。但是，想要的是，坐标的变换的关系。 」\par
千奈「……啊，有非常不好的言感……」\par
广「坐标是和基底相反的变化。也就是说，需要求这个的逆矩子。 」\par
千奈「……」\par
广「没关系。其实，不把规范正交基设为\{e1, e2, e3\}，就成了下面这样。 」\par
(\par
e1T\par
e2T\par
e3T\par
) (e1 e2 e3) = (\par
e1Te1 e1Te2 e1Te3\par
e2Te1 e2Te2 e2Te3\par
e3Te1 e3Te2 e3Te3\par
) = (\par
1 0 0\par
0 1 0\par
0 0 1\par
)\par
广「进而，下式成立。 」\par
(e1 e2 e3)T = (\par
e1T\par
e2T\par
e3T\par
)\par
广「才才求出来的，是规范正交基，所以能用这个。也就是说，逆矩子一瞬间就能\par
求出。 」\par
千奈「太好了……」\par

% --- PDF page 175 ---
174\par
(\par
vx\par
v\par
vy\par
$\sqrt{}$vx2 + vy2\par
vzvx\par
v$\sqrt{}$vx2 + vy2\par
vy\par
v - vx\par
$\sqrt{}$vx2 + vy2\par
vzvy\par
v$\sqrt{}$vx2 + vy2\par
vz\par
v 0 - vx2 + vy2\par
v$\sqrt{}$vx2 + vy2\par
)\par
-1\par
=\par
(\par
vx\par
v\par
vy\par
v\par
vz\par
vvy\par
$\sqrt{}$vx2 + vy2\par
- vx\par
$\sqrt{}$vx2 + vy2\par
0\par
vzvx\par
v$\sqrt{}$vx2 + vy2\par
vzvy\par
v$\sqrt{}$vx2 + vy2\par
- vx2 + vy2\par
v$\sqrt{}$vx2 + vy2\par
)\par
千奈「做出来了……咦，但是，说起来，时间会怎么样呢?」\par
广「只是改变了轴的取法而已，所以对千奈而言的时间，不会被变换。所以，时间\par
就那样，没问题。因为逐一写长式子很烦了，所以用v-z = $\sqrt{}$vx2 + vy2来改写掉吧。 」\par
(\par
X\par
Y\par
Z\par
ct\par
) =\par
(\par
vx v⁄ vy v⁄ vz v⁄ 0\par
vy v-z⁄ - vx v-z⁄ 0 0\par
vzvx (vv-z)⁄ vzvy (vv-z)⁄ - v-z v⁄ 0\par
0 0 0 1)\par
(\par
x\par
y\par
z\par
ct\par
)\par
广「这样，想求的坐标变换就求出来了。现在，只是重新取了千奈世界里的坐标而\par
已，所以接着，来看看千奈的坐标与佑芽的坐标的关系。 」\par
广 「首，， 设佑芽的坐标为(X$\prime$, Y$\prime$, Z$\prime$, ct$\prime$)， 取坐标轴与千奈的X, Y, Z轴平行。 这些坐\par
标的关系是怎样的呢。 」\par
千奈「这个是，才才做过的那个对吧!」\par
(\par
X$\prime$\par
Y$\prime$\par
Z$\prime$\par
ct$\prime$\par
) = (\par
$\gamma$ 0 0 -$\gamma$$\beta$\par
0 1 0 0\par
0 0 1 0\par
-$\gamma$$\beta$ 0 0 $\gamma$\par
) (\par
X\par
Y\par
Z\par
ct\par
)\par
广「…，成了那样。进而，取佑芽的坐标x$\prime$, y$\prime$, z$\prime$的坐标轴与千奈的x, y, z轴平行的\par
话，从x$\prime$, y$\prime$, z$\prime$到X$\prime$, Y$\prime$, Z$\prime$的变换，就与从x, y, z到X, Y, Z的变换相同。这个关系是怎样的\par
呢。 」\par
千奈「应该是这样!」\par
(\par
x$\prime$\par
y$\prime$\par
z$\prime$\par
ct$\prime$\par
) = (\par
vx v⁄ vy v-z⁄ vzvx (vv-z)⁄ 0\par
vy v⁄ - vx v-z⁄ vzvy (vv-z)⁄ 0\par
vz v⁄ 0 - v-z v⁄ 0\par
0 0 0 1\par
) (\par
X$\prime$\par
Y$\prime$\par
Z$\prime$\par
ct$\prime$\par
)\par

% --- PDF page 176 ---
175\par
广「那么，把至今的式子全部结的员，求求看(x, y, z, t)和(x$\prime$, y$\prime$, z$\prime$, t$\prime$)之间成立的洛\par
伦兹变换吧。 」\par
千奈「我明白了!」\par
(\par
x$\prime$\par
y$\prime$\par
z$\prime$\par
ct$\prime$\par
) =\par
(\par
vx v⁄ vy v-z⁄ vzvx (vv-z)⁄ 0\par
vy v⁄ - vx v-z⁄ vzvy (vv-z)⁄ 0\par
vz v⁄ 0 - v-z v⁄ 0\par
0 0 0 1)\par
(\par
X$\prime$\par
Y$\prime$\par
Z$\prime$\par
ct$\prime$\par
)\par
=\par
(\par
vx v⁄ vy v-z⁄ vzvx (vv-z)⁄ 0\par
vy v⁄ - vx v-z⁄ vzvy (vv-z)⁄ 0\par
vz v⁄ 0 - v-z v⁄ 0\par
0 0 0 1)\par
(\par
$\gamma$ 0 0 -$\gamma$$\beta$\par
0 1 0 0\par
0 0 1 0\par
-$\gamma$$\beta$ 0 0 $\gamma$\par
) (\par
X\par
Y\par
Z\par
ct\par
)\par
=\par
(\par
$\gamma$vx v⁄ vy v-z⁄ vzvx (vv-z)⁄ -$\gamma$$\beta$vx v⁄\par
$\gamma$vy v⁄ - vx v-z⁄ vzvy (vv-z)⁄ -$\gamma$$\beta$vy v⁄\par
$\gamma$vz v⁄ 0 - v-z v⁄ -$\gamma$$\beta$vz v⁄\par
-$\gamma$$\beta$ 0 0 $\gamma$ )\par
(\par
X\par
Y\par
Z\par
ct\par
)\par
=\par
(\par
$\gamma$vx\par
v\par
vy\par
v-z\par
vzvx\par
(vv-z)\par
-$\gamma$$\beta$vx\par
v\par
$\gamma$vy\par
v - vx\par
v-z\par
vzvy\par
(vv-z)\par
-$\gamma$$\beta$vy\par
v\par
$\gamma$vz\par
v 0 - v-z\par
v\par
-$\gamma$$\beta$vz\par
v\par
-$\gamma$$\beta$ 0 0 $\gamma$ )\par
(\par
vx\par
v\par
vy\par
v\par
vz\par
v 0\par
vy\par
v-z\par
- vx\par
v-z\par
0 0\par
vzvx\par
(vv-z)\par
vzvy\par
(vv-z) - v-z\par
v 0\par
0 0 0 1)\par
(\par
x\par
y\par
z\par
ct\par
)\par
=\par
(\par
1 + ($\gamma$ - 1)vx2\par
v2\par
($\gamma$ - 1)vxvy\par
v2\par
($\gamma$ - 1)vzvx\par
v2 - $\gamma$$\beta$vx\par
v\par
($\gamma$ - 1)vxvy\par
v2 1 +\par
($\gamma$ - 1)vy2\par
v2\par
($\gamma$ - 1)vyvz\par
v2 - $\gamma$$\beta$vy\par
v\par
($\gamma$ - 1)vzvx\par
v2\par
($\gamma$ - 1)vyvz\par
v2 1 + ($\gamma$ - 1)vz2\par
v2 - $\gamma$$\beta$vz\par
v\par
- $\gamma$$\beta$vx\par
v - $\gamma$$\beta$vy\par
v - $\gamma$$\beta$vz\par
v $\gamma$ )\par
(\par
x\par
y\par
z\par
ct\par
)\par
广「这样，想求的洛伦兹变换就求出来了。辛苦了。 」\par
千奈「结、结算做完了……」\par
广「最后，设v = (v, 0, 0)，确认一下和才才的洛伦兹变换是具一『吧。 」\par
千奈「…，含有vy, vz的项都变成 0 了，所以……」\par

% --- PDF page 177 ---
176\par
(\par
x$\prime$\par
y$\prime$\par
z$\prime$\par
ct$\prime$\par
) = (\par
1 + $\gamma$ - 1 0 0 -$\gamma$$\beta$\par
0 1 0 0\par
0 0 1 0\par
-$\gamma$$\beta$ 0 0 $\gamma$\par
) (\par
x\par
y\par
z\par
ct\par
) = (\par
$\gamma$ 0 0 -$\gamma$$\beta$\par
0 1 0 0\par
0 0 1 0\par
-$\gamma$$\beta$ 0 0 $\gamma$\par
) (\par
x\par
y\par
z\par
ct\par
)\par
千奈「一『了啊!」\par
广「恭喜。今天就到这里吧。计算很多，真的辛苦了。 」\par
千奈「好累啊……且时不想再看到矩子了……」\par
Day3 狭义相对论\par
今天是休息日， 但不巧外面台风来袭， 天气极为恶劣， 所以定定白天集合开研讨会。\par
进入视频通话时，\par
泽同学似乎在吃着什么。\par
广 P「请稍微吃点看起来好吃的东西啊。 」\par
广「…。虽然看起来是那样，但还挺好吃的。。 」\par
一看，正把某种鲜粉色的东西往嘴里送。\par
千 P「辛苦了。直接问一下，那是什么?」\par
广「啊，千奈的制作人。辛苦了。 」\par
泽同学一边腮帮鼓鼓的，一边用手掩着嘴回答。\par
广「这是，牛肉饭。 」\par
千 P「怎么看都不像是牛肉饭啊。 」\par
大概在粉色下面藏着牛肉饭吧， 但结之，那鲜粉色的糊态太过引人注意，什么都起\par
不进的了。\par
广「…。这个，说是把山药泥和红姜放搅拌机里打成糊做的。 」\par
千 P「 「说是」……这，不是\par
泽同学做的 。 」\par
广「…。是第 2 名的人给我做的。 」\par
广 P「果然犯人是咲季小姐啊。 」\par
问了食材后松了口气，但即一如此，仍是让人毫无食欲的颜色和外。。 之后\par
泽同\par
学继续吃着咲季的手作一当，吃完时正好仓本同学进来了。\par
广「那么，开始吧。 」\par
～（研讨会开始）～\par
广「今天来处理狭义相对论。 」\par
千奈「今天也请多关照!」\par

% --- PDF page 178 ---
177\par
广 「在进入正题之前， 必须，说， 相对论有两种理论。 狭义相对论， 和广义相对论。 」\par
千奈「狭义，和，广义 ?」\par
广「…。从名字大概能想象，与其说是两种理论，不如说，狭义相对论是广义相对\par
论的特殊情况，这种感觉。千奈，还记得上次研讨会的惯性坐标系话题 。 」\par
千奈「是的!所谓惯性坐标系，就是惯性定律成立的坐标系对吧。 」\par
广「对。那么，具体来说怎样的坐标系会成为惯性坐标系，记得 ?」\par
千奈「…，在完全无重力的地方，静止，或做匀速运的的情况。 」\par
广「是啊。相反，非惯性坐标系，就是不是那样的坐标系的坐标系。」\par
广 「狭义相对论， 说到 「狭义」是指什么， 就是只能在惯性坐标系之间使用的理论。\par
那么，说到广义相对论，就如想象，是连在非惯性坐标系里也能用的理论。所以，广义\par
相对论比狭义相对论，对事物的扩展性更高。 」\par
千奈「来来如此……也就是说， 今后的话题， 全都是仅在惯性坐标系之间才成立的\par
话题，对 。 」\par
广「…，就是这回事。把这个放在心上，起今后的内容，会很有帮助。 」\par
广「那么，在这里突然问千奈一个问题。佑芽同学，在t = 0时，从x = 0处的是出\par
门， 以每分钟60 m的速度朝车站的了。 佑芽同学的姐姐咲季同学， 注意到佑芽同学忘带\par
了一当，在佑芽同学出发 6 分钟后，骑自行车以每分钟240 m追佑芽同学。但是，不论\par
怎么追咲季同学都找不到佑芽同学， 咲季同学骑自行车7 分钟到达车站后， 认为佑芽同\par
学不可能走这么快，马上沿来路折返了。请问，咲季同学能把一当送到佑芽同学手上，\par
是在咲季同学出发后几分钟?」\par
千奈「条件好多啊……，画图来想吧……」\par

% --- PDF page 179 ---
178\par
千奈「做出来了!是 10 分钟后。 」\par
广「…，正确。对了，才才画了什么样的图呢。 」\par
千奈「诶?横轴取t、纵轴取x，只是这样画了图而已……」\par
广「其实，在相对论里那个图很重要。 无示空间坐标位置如何着时间变化而变化的\par
线，叫世界线。这里，算是画了佑芽和咲季同学两个人的世界线。便一，当然，世界线\par
和运的轨迹是不同的，请注意。 世界线终究只是记述物体如何着时间移的， 并没有无示\par
物体实际走了怎样的轨迹。 」\par
千奈「结之，因为t 轴表示时间的经过，所以要把世界线，教只用空间坐标讨论的\par
轨迹区别开，对吧!」\par
广「就是这么回事。还有，在这里，虑一下在各坐标系里，换算时间单位和空间单\par
位。 在相对论里， 承认光速不变来理。 这也就是说， 对各。测者， 光速是一样的， 所以，\par
不论什么坐标系，似乎都能用光速当尺子吧。 」\par
千奈「速度当尺子……?」\par

% --- PDF page 180 ---
179\par
广「具体来说，距离用 「光走a秒」这样的表达法就好。时间的 「b秒」部分，因人\par
而流逝方式不同，但这正对应于在坐标变换那里说过的，对每个人各自的「1」的长度\par
不同。结之在这里，掌握住能用光速把距离换算成时间来表示就好。 」\par
千奈「来来如此!」\par
广「那么，这样一来，就能把目前为止用(x, y, z, t)表示的坐标，用(x/c, y/c, z/c, t)\par
表示， 把各轴的量纲统一为时间了。 可是，3 个分量里都除以 c， 稍微有点烦了。 于是，\par
反过来在 t 上乘以光速 c，变成(x, y, z, ct)。这样算是把量纲统一成了长度， 但因为说的\par
是和才才相同的事情，所以今后定定使用这个。 」\par
千奈「用ct而不是t，来来是这么个来因啊……」\par
广 「对。 那么， 把t换成ct， 这次把运的方向增加到2 维来，虑世界线吧。 ，虑世界\par
线时， 需要空间坐标和时间轴， 但这次普通，虑xy平面， 从那个平面垂直伸出t轴就行。\par
那么， 设设千奈， 正坐着转转木马玩， 在中心(x, y) = (0,0)、 半径r 的圆上滴溜溜地转。\par
这时候，千奈的世界线是怎样的呢。 」\par
千奈「…，会变成这样 ……?」\par
广「…，应该会形成那种螺转。这里，如果沿着t 轴方向投射世界线的影子，那个\par
影子就成了轨迹，能明白 。」\par
千奈「确实，会变成圆呢。 」\par

% --- PDF page 181 ---
180\par
广「还有，如果想用垂直于t轴的某个断面，也就是表示为t = t0的平面来切，世界\par
线和那个断面，必定相交于 1 点。这，表示在时刻 t0 时千奈的位置，也就是，会变成\par
像是「那个瞬间拍摄的照片」一样的东西，这么回事。 」\par
千奈「来来如此，结算是形成印象了!」\par
广「太好了。那么接下来，把t 轴换成ct 轴，和上次一样，，虑千奈拿着灯泡的情\par
况。 」\par
广 「现在， 设设千奈在时刻t = 0点亮了手边的灯泡。 这一瞬间发出的光的世界线，\par
会怎样呢。 」\par
千奈 「…——， 在时刻t的瞬间， 光应该位于ct的位置。 因为光应该以自己为中心呈\par
圆形前进……所以应该是这样!」\par

% --- PDF page 182 ---
181\par
广 「就是那样。 这次因为光是向四面下方飞的， 所以变成了面。 这个会成为圆锥面，\par
但这个面无示了光的世界线，所以也叫光锥。 」\par
广「接着，对t < 0也扩展世界线的概念。t < 0因为是。测者的现在t = 0之前，所\par
以是对应。测者过的的部分。这里也能，虑光的世界线。 具体来说，只要，虑到达千奈\par
的光，是通过怎样的世界线过来的就好。 」\par
千奈「例如，在时间t0 > 0 期间光前进ct0 ，所以如果我和光的距离在t = -t0 时为\par
ct0，则在t = 0时光会到达。所以，会变成这样!」\par
广「就是那样。便一，不用平行于ct轴的平面的切这个圆锥面，不论用怎样的平面\par
切，都会得到下面这样的曲线。 」这个叫双曲线。之后会处理双曲线，但请记住，双曲\par
线和圆锥面是有关联的。 」\par

% --- PDF page 183 ---
182\par
千奈「咦，如果是包含ct轴的平面，就不是曲线，而是两条直线了啊。 」\par
广「确实如此。但是，这个也可以看作双曲线之一。 」\par
千奈「为什么呢……?」\par
广「设设用距离ct 轴为r 的某个平面的切这个圆锥面，得到了双曲线。从垂直于那\par
个平面的方向的投射ct轴，相对于那个轴的影子，在t = 0的位置正交地取s轴吧。方向\par
朝哪边都行。 」\par
广「于是，这个双曲线，可以表示为(ct)2 - s2 = r2。现在，变量是ct, s两个，r是\par
按各平面定定的常数。那么，如果是包含ct轴的平面，会怎样呢。 」\par
千奈「…，这种情况r = 0，所以代入这个……」\par
(ct)2 - s2 = 0 $\Leftrightarrow$ s = $\pm$ct\par
千奈「这是什么……?」\par
广 「这个， 如果把ct轴和s轴各自换读成x, y， 就表示y = $\pm$x， 也就是y = x和y = -x\par
两条直线。因为这是光锥的断面，所以必须与光的世界线一『，但这表示在ct = ct0，\par
也就是t = t0时，光沿着s 轴在$\pm$ct0的地方，这么回事。这可以看作是沿着s轴的光的行\par
为， 所以似乎没问题。 因此， 这个式子， 确实对应着用包含ct轴的平面切时的两条直线。 」\par
千奈「来来如此!!」\par
广「…。所以，这个也能看作双曲线之一。能接受了 。 」\par
千奈「接受了!」\par
广「太好了。那么，在这里更详细地看看双曲线。千奈，学过三角函数了 。 」\par

% --- PDF page 184 ---
183\par
千奈「是的，学过了。 」\par
广「那么，(x, y) = (cos $\theta$ , sin $\theta$)在什么位置呢。 」\par
千奈「是从点(1,0)开始，以来点为中心，转转$\theta$后的地方。 」\par
广「是啊。那么，这个单位圆的式子，是怎样的呢。 」\par
千奈「是x2 + y2 = 1。」\par
广 「就是那样。(x, y) = (cos $\theta$ , sin $\theta$)是这上面的点， 所以cos2 $\theta$ + sin2 $\theta$ = 1成立。\par
那么，连接来点与这个点的直线的斜率，是怎样的呢。 」\par
千奈「是tan $\theta$!」\par
广「…，是啊。tan $\theta$是如下定义的函数。 」\par
tan $\theta$ = sin $\theta$\par
cos $\theta$\par
广 「对了， 回想才才双曲线的式子，x2 - y2 = 1也表示双曲线。 这个教单位圆的式\par
子非常像。在这里，就定定把这个式子所表示的双曲线，称为单位双曲线吧。 」\par
广「那么，像三角函数sin $\theta$ , cos $\theta$那样，能不能用函数表示双曲线x2 - y2 = 1上的\par
点呢?来，虑一下。这里，试着将双曲函数sinh u , cosh u , tanh u定义如下吧。 」\par
sinh u = eu - e-u\par
2 , cosh u = eu + e-u\par
2 , tanh u = sinh $\theta$\par
cosh $\theta$\par
广 「其实， 使用这个函数， 就能表示为(x, y) = (cosh $\theta$ , sinh $\theta$)这样。 来实际计算确\par
认一下看看。 」\par
千奈「我明白了!」\par

% --- PDF page 185 ---
184\par
cosh2 u - sinh2 u = e2u + 2 + e-2u\par
4 - e2u - 2 + e-2u\par
4 = 2 + 2\par
4 = 1\par
千奈「确实，似乎不论什么 u 进的都会成立呢。 」\par
广「…。这里，把u想成是相当于三角函数里$\theta$的一种角度。因为$\theta$是角度，所以就\par
把u叫双曲角度吧。 」\par
千奈「………的确， 这样定义双曲函数的话， 就能漂亮地像三角函数一样表示双曲\par
线上的点，这点我明白了。可是，突然间啪一下看到这函数，为什么是这种形式，实在\par
不太能接受啊……」\par
广「这个函数以此形式出现，其实有物理学的经过的。拿着绳子的两端，不把两端\par
拉紧而让绳子垂下，就会变成类似抛物线y = x2 形态的曲线。例如，用身边的例子说，\par
想象电线下垂的形态或许就好。 但这个其实， 是教抛物线不同的曲线， 叫悬链线的曲线。\par
这个悬链线会变成怎样，用力学来，虑的话， cosh x 就出现了。虽然也可以广致这个，\par
但我想现在稍微有点难， 所以想请觉们看看双曲函数的性质，感受一下 「教三角函数好\par
像啊」。」\par

% --- PDF page 186 ---
185\par
广 P「更数学地，靠复数来连接。在复函数论中，由欧拉公式eiz = cos z + i sin z，\par
成为cos z =\par
eiz+e-iz\par
2 , sin z =\par
eiz-e-iz\par
2i ，能看出，在这里出现了双曲函数的形式。 」\par
广「……那么，千奈，学过三角函数的微分了 。 」\par
千奈「是的!应该是这样!」\par
d\par
d$\theta$ sin $\theta$ = cos $\theta$ , d\par
d$\theta$ cos $\theta$ = - sin $\theta$ , d\par
d$\theta$ tan $\theta$ = 1\par
cos2 $\theta$\par
广「那么同样地，试把双曲函数对u微分一下吧。 」\par
千奈「我明白了!」\par
d\par
du sinh u = eu + e-u\par
2 = cosh u , d\par
du cosh u = eu - e-u\par
2 = sinh u\par
d\par
du tanh u = cosh u $\cdot$ cosh u - sinh u $\cdot$ sinh u\par
cosh2 u = 1\par
cosh2 u\par
千奈「确实，教三角函数很像!」\par
广 「接着， 来看看角度取反号时的变化吧。 三角函数是怎样来着，就是这种感觉。 」\par
cos(-$\theta$) = cos $\theta$ , sin(-$\theta$) = - sin $\theta$ , tan(-$\theta$) = - sin $\theta$\par
cos $\theta$ = - tan $\theta$\par
广「双曲函数里会怎样呢。 」\par
千奈「来做做看!」\par

% --- PDF page 187 ---
186\par
cosh(-u) = e-u + eu\par
2 = cosh u , sinh(-u) = e-u - eu\par
2 = - eu - e-u\par
2 = - sinh u,\par
tanh(-u) = - sinh u\par
cosh u = - tanh u\par
千奈「确实是同样的感觉!」\par
广「…。最后，加法公式。这个稍微有点复杂，但有这样的对应关系。关于双曲函\par
数，计算右边，确认它变成左边就好了。 」\par
sin($\alpha$ + $\beta$) = sin $\alpha$ cos $\beta$ + cos $\alpha$ sin $\beta$\par
对应\par
$\leftrightarrow$ sinh(u + v) = sinh u cosh v + cosh u sinh v\par
cos($\alpha$ + $\beta$) = cos $\alpha$ cos $\beta$ - sin $\alpha$ sin $\beta$\par
对应\par
$\leftrightarrow$ cosh(u + v) = cosh u cosh v + sinh u sinh v\par
tan($\alpha$ + $\beta$) = tan $\alpha$ + tan $\beta$\par
1 - tan $\alpha$ tan $\beta$\par
对应\par
$\leftrightarrow$ tanh(u + v) = tanh u + tanh v\par
1 + tanh u tanh v\par
广「这样就大『能接受了吧。那么，就用这个双曲函数，进行更深入的相对论讨论\par
吧。首，，还记得洛伦兹变换 。 」\par
千奈「是的，是这样的对吧!」\par
(\par
x$\prime$\par
y$\prime$\par
z$\prime$\par
ct$\prime$\par
) =\par
(\par
1 + ($\gamma$ - 1)vx2\par
v2\par
($\gamma$ - 1)vxvy\par
v2\par
($\gamma$ - 1)vzvx\par
v2 - $\gamma$$\beta$vx\par
v\par
($\gamma$ - 1)vxvy\par
v2 1 +\par
($\gamma$ - 1)vy2\par
v2\par
($\gamma$ - 1)vyvz\par
v2 - $\gamma$$\beta$vy\par
v\par
($\gamma$ - 1)vzvx\par
v2\par
($\gamma$ - 1)vyvz\par
v2 1 + ($\gamma$ - 1)vz2\par
v2 - $\gamma$$\beta$vz\par
v\par
- $\gamma$$\beta$vx\par
v - $\gamma$$\beta$vy\par
v - $\gamma$$\beta$vz\par
v $\gamma$ )\par
(\par
x\par
y\par
z\par
ct\par
)\par
广「…。这里， (x, y, z, ct) 是千奈的坐标，(x$\prime$, y$\prime$, z$\prime$, ct$\prime$) 是对从千奈看以速度 v =\par
(vx, vy, vz)运的的佑芽而言的坐标。这里，$\beta$, $\gamma$各自是什么式子来着。 」\par
千奈「是这样!」\par
$\beta$ = v\par
c , $\gamma$ = 1\par
$\sqrt{}$1 - $\beta$2\par
广「是啊。这次，为让讨论单计化，把x, x$\prime$轴和速度向量的方向对齐吧。 」\par
(\par
x$\prime$\par
y$\prime$\par
z$\prime$\par
ct$\prime$\par
) = (\par
$\gamma$ 0 0 -$\gamma$$\beta$\par
0 1 0 0\par
0 0 1 0\par
-$\gamma$$\beta$ 0 0 $\gamma$\par
) (\par
x\par
y\par
z\par
ct\par
) $\Leftrightarrow$ (\par
x\par
y\par
z\par
ct\par
) = (\par
$\gamma$ 0 0 $\gamma$$\beta$\par
0 1 0 0\par
0 0 1 0\par
$\gamma$$\beta$ 0 0 $\gamma$\par
) (\par
x$\prime$\par
y$\prime$\par
z$\prime$\par
ct$\prime$\par
)\par

% --- PDF page 188 ---
187\par
广「这里，为了习惯洛伦兹变换，首，来看看由洛伦兹变换致出的重要性质。 」\par
广「首，，关于物体的速度。设设千奈和佑芽正在。改同一个物体。设对千奈而言\par
的物体位置为x = (x, y, z)， 对佑芽而言的物体位置为x$\prime$ = (x$\prime$, y$\prime$, z$\prime$)。 当然，x, x$\prime$会着时\par
间广移而变化。这里设设物体上没有力在作用。这里，速度虽然曾是 「每单位时间位置\par
变化多少」，但这次，时间也好，位置的变化程度也好，千奈和佑芽各自不同。所以，\par
设 对 千 奈 而 言 物 体 的 速 度 为w = (wx, wy, wz) ， 对 佑 芽 而 言 物 体 的 速 度 为w$\prime$ =\par
(wx$\prime$ , wy$\prime$ , wz$\prime$)，则下式成立。 」\par
w = (\par
wx\par
wy\par
wz\par
) = (\par
dx dt⁄\par
dy dt⁄\par
dz dt⁄\par
) , w$\prime$ = (\par
w$\prime$x\par
w$\prime$y\par
w$\prime$z\par
) = (\par
dx$\prime$ dt$\prime$⁄\par
dy$\prime$ dt$\prime$⁄\par
dz$\prime$ dt$\prime$⁄\par
)\par
广 「现在， 来，虑这个w, w$\prime$之间的变换。 关于微分， 有\par
dx$\prime$\par
dt$\prime$ =\par
dt\par
dt$\prime$\par
dx$\prime$\par
dt 成立。 其他分量\par
也同样变形的话，就成下面这样。 」\par
w$\prime$ = (\par
dx$\prime$ dt$\prime$⁄\par
dy$\prime$ dt$\prime$⁄\par
dz$\prime$ dt$\prime$⁄\par
) = dt\par
dt$\prime$ (\par
dx$\prime$ dt⁄\par
dy$\prime$ dt⁄\par
dz$\prime$ dt⁄\par
)\par
广「这里，使用洛伦兹变换的话，首，dt/dt$\prime$变成这样。 」\par
dt\par
dt$\prime$ = cdt\par
cdt$\prime$ = $\gamma$$\beta$dx$\prime$ + $\gamma$cdt$\prime$\par
cdt$\prime$ = $\gamma$$\beta$\par
c\par
dx$\prime$\par
dt$\prime$ + $\gamma$ = $\gamma$ (1 + $\beta$w$\prime$x\par
c )\par
广「进而，把x$\prime$对t微分则如下。 」\par
(\par
dx$\prime$ dt⁄\par
dy$\prime$ dt⁄\par
dz$\prime$ dt⁄\par
) = (\par
($\gamma$dx - $\gamma$$\beta$cdt)/dt\par
dy/dt\par
dx/dt\par
) = (\par
$\gamma$wx - $\gamma$$\beta$c\par
wy\par
wz\par
) = (\par
$\gamma$wx - $\gamma$v\par
wy\par
wz\par
)\par
广「因此，代入这个的话，就成下面这样。 」\par
w$\prime$ = $\gamma$ (1 + $\beta$w$\prime$x\par
c ) (\par
$\gamma$wx - $\gamma$v\par
wy\par
wz\par
)\par
广「看看这个x$\prime$分量，关于wx整理式子的话，就成下面这样。 」\par
w$\prime$x = $\gamma$ (1 + $\beta$w$\prime$x\par
c ) ($\gamma$wx - $\gamma$v) $\Leftrightarrow$ wx = w$\prime$x\par
$\gamma$2 (1 + $\beta$w$\prime$x\par
c )\par
+ v\par

% --- PDF page 189 ---
188\par
广 「因为$\gamma$ = 1/$\sqrt{}$1 - $\beta$2， 代入这个， 设$\beta$x =\par
wx\par
c , $\beta$x$\prime$ = wx$\prime$ /c， 两边除以c 整理的话，\par
就成下面这样。 」\par
wx = w$\prime$x(1 - $\beta$2)\par
1 + $\beta$$\beta$x$\prime$\par
+ v $\Leftrightarrow$ $\beta$x = $\beta$x$\prime$(1 - $\beta$2)\par
1 + $\beta$$\beta$x$\prime$\par
+ $\beta$ = $\beta$ + $\beta$x$\prime$\par
1 + $\beta$$\beta$x$\prime$\par
广「如果想变成求$\beta$x$\prime$的式子的话，只需把$\beta$的v换成-v就行。 」\par
$\beta$x$\prime$ = -$\beta$ + $\beta$x\par
1 - $\beta$$\beta$x\par
广「由此，可知下面这样的变换成立。 」\par
w$\prime$ = $\gamma$ (1 - $\beta$2 + $\beta$$\beta$x\par
1 - $\beta$$\beta$x\par
) (\par
$\gamma$wx - $\gamma$v\par
wy\par
wz\par
) = $\gamma$ 1 - $\beta$2\par
1 - $\beta$$\beta$x\par
(\par
$\gamma$wx - $\gamma$v\par
wy\par
wz\par
)\par
千奈「结感觉教洛伦兹变换完全不一样了呢……」\par
广「对。接着，，虑速度的合成法则。现在，设设我坐在速度为w的物体上，为简\par
单起见设w = (w, 0,0)。那么w$\prime$会怎样呢。 」\par
千奈「…，代入wx = w, wy = wz = 0……」\par
w$\prime$ = $\gamma$ 1 - $\beta$2\par
1 - $\beta$(w/c) (\par
$\gamma$w - $\gamma$v\par
0\par
0\par
)\par
千奈「wy, wz似乎会变成 0!那样的话，只，虑x$\prime$分量的话……」\par
w$\prime$ = $\gamma$2 1 - $\beta$2\par
1 - $\beta$(w c⁄ ) (w - v) = 1\par
1 - $\beta$2 $\cdot$ 1 - $\beta$2\par
1 - $\beta$(w c⁄ ) (w - v) = w - v\par
1 - $\beta$(w c⁄ )\par
千奈「做出来了!」\par
广「…。这，就是从佑芽看时的我的速度。那么，将$\beta$ = v/c重新写成$\beta$v。进而，\par
设$\beta$w: = w/c。这样一来，这个式子也能写成这样。 」\par
w$\prime$ = $\beta$w - $\beta$v\par
1 - $\beta$v$\beta$w\par
c\par
广「进而，从佑芽看千奈时的速度为-v，所以由此设$\beta$-v: = -v/c = -$\beta$v的话，就\par
成下面这样。这个式子，不觉得在哪见过 。 」\par
w$\prime$ = $\beta$w + $\beta$-v\par
1 + $\beta$-v$\beta$w\par
c\par

% --- PDF page 190 ---
189\par
千奈「啊!是tanh u的加法公式!」\par
广「就是那样。稍微整理现在的态况，上面的式子w$\prime$是从佑芽看我时的速度。为此\par
所使用着的，是从佑芽看的千奈的速度-v，和从千奈看的我的速度w。这，在伽利略变\par
换的世界里， 曾是单计的w$\prime$ = -v + w， 但在洛伦兹变换的世界里， 却成了tanh u的加法\par
公式的形式。 」\par
千奈「来来如此，洛伦兹变换和双曲函数，似乎有什么关系呢。 」\par
广「进而从上式，会出现思，相对论上重要的概念。千奈，把从佑芽看的千奈的速\par
度设为光速的k1倍， 从千奈看的我的速度设为光速的k2倍， 试计算一下w$\prime$是c的多少倍。 」\par
千奈「我明白了……」\par
$\beta$w + $\beta$-v\par
1 + $\beta$-v$\beta$w\par
= k2 + k1\par
1 + k1k2\par
千奈「是\par
k2+k1\par
1+k1k2\par
倍。 」\par
广「…。这里，如果设定不存在超过光速的物体，那么k1 $\leq$ 1, k2 $\leq$ 1。也就是说，\par
1 - k1\par
2 > 0, 1 - k2\par
2 > 0。由此，下式成立。 」\par
1 - k2 + k1\par
1 + k1k2\par
= 1 - k2 - k1 + k1k2\par
1 + k1k2\par
= (1 - k1)(1 - k2)\par
1 + k1k2\par
$\geq$ 0 $\Leftrightarrow$ k2 + k1\par
1 + k1k2\par
$\leq$ 1\par
广「也就是说，由此可知，一切物体就算进行速度的合成，在洛伦兹变换下也不会\par
超过光速 c。目前，还没有发现超过光速的物质，所以这世上的任何物体，都无法超过\par
光的速度。 」\par
千奈「就是不管多快，都没法超过光，这么回事吧!」\par
广「就是那样。由此，可以致出因果律。就某。测者的时空来，虑的话，通过来点\par
的世界线，不允许有比光速更快移的的瞬间。也就是说，不管是。测者在来点扔应，或\par
是用水管洒水，被扔出的应或洒出的水，都不得超越光。这个，不用世界线换言，就是\par
世界线不得从光锥的ct轴侧逸出。所谓逸出，意味着不设逸出瞬间的时刻为t0，，虑在\par
t0 拍摄的照片，也就是平面t = t0 处的断面的话，物体的位置就变成了超越光的位置，\par
所以就不合理了。 」\par

% --- PDF page 191 ---
190\par
广 「关于过的， 也同样可以说。 过的要对在t = 0位于来点的。测者产生影响， 世界\par
线就必须通过光锥之中。 结之， 在。测者绪顶上雨云铺开， 雨滴从那里落下， 在t = 0的\par
瞬间打到。测者身上时，这雨滴在过的，不可运的得比光更快，这么回事。 」\par
广「由此，。测者在t = 0时受到其他东西什么影响的事项，或者。测者在t = 0时\par
做什么事的影响其他东西的事项，所有这些，其世界线都必须通过光锥之中。于是，把\par
光锥内部当中，t < 0的部分叫因果过的，t > 0的部分叫因果未来。相反，无法影响到\par
。测者的过的领域， 或。测者无法施加影响的未来领域， 则成为光锥的外侧。这叫空间\par
领域。 」\par

% --- PDF page 192 ---
191\par
千奈「虽然稍微有点复杂……结之，因果过的指能对。测者施加影响的事项领域，\par
因果未来指。测者能施加影响的事项领域，对吧!」\par
广「就是那样。这件事，用算式来写，就像这样。」\par
\{x2 + y2 + z2 > (ct)2 → 光锥的内侧 = 因果未来 or 因果过的\par
x2 + y2 + z2 < (ct)2 → 光锥的外侧 = 空间领域\par
广「那么，一旦明白了因果律，就能一口气理解各种事情。 」\par
广 「例如， 设设。测者在t = 0用锤子敲了手边木棒的一端。 于是， 通过锤子传给这\par
根木棒的力的传递速度， 就不得超越光速。 为什么呢， 因为在t = 0的瞬间， 棒的另一端\par
还没进入光锥里面。因锤子敲打了成的影响，必须进入因果未来之中。从这个例子，可\par
知力传递的速度，不得超越光速。 」\par
广「进而，上次的开绪话题里，我想讲过了质量的话题，但从物体无法超越光速可\par
知，如果认为质量一直不变，那不论多小的力，着着时间流逝都会超过光速。所以，必\par
须重新思，质量。由此，得到了「惯性质量必须是速度v的函数，具则就不合理」这个\par
结论。」\par
千奈「来来如此，也就是越快的东西变得越重，对吧!」\par
广「那个话题之后再说。现在再教洛伦兹变换更熟络一点吧。那么，从目前为止的\par
话题， 因为$\beta$是表示坐标系相对速度为光速的多少倍的量， 所以可以说$\beta$小于 1。 基于此，\par
可知存在使$\beta$ = tanh u成立的u。」\par

% --- PDF page 193 ---
192\par
广 「那么， 使用这个$\beta$ = tanh u的关系， 把洛伦兹变换用u来表示的话， 会怎样呢。 」\par
( x$\prime$\par
ct$\prime$)=( $\gamma$ -$\gamma$$\beta$\par
-$\gamma$$\beta$ $\gamma$ ) ( x\par
ct)\par
千奈「…，首，计算$\gamma$……」\par
$\gamma$ = 1\par
$\sqrt{}$1 - tanh2 u\par
= 1\par
$\sqrt{}$(cosh2 u - sinh2 u)/ cosh2 u\par
= 1\par
$\sqrt{}$1/ cosh2 u\par
= cosh u\par
千奈「因为$\gamma$$\beta$ = cosh u tanh u = sinh u，所以用这个的话，就变成这样了!」\par
( x$\prime$\par
ct$\prime$)=( cosh u - sinh u\par
-sinh u cosh u ) ( x\par
ct)\par
广 「会那样。 那么， 为了思，这意味着什么，设(x, ct) = (k cosh $\phi$ , k sinh $\phi$)来试着\par
计算。 这里， 只要改变k, $\phi$的取。方式， 就能对应光锥外的任意x, t， 所以不用担心那里。\par
千奈，计算一下看看。 」\par
千奈「我明白了!」\par
( cosh u - sinh u\par
-sinh u cosh u ) (k cosh $\phi$\par
k sinh $\phi$) = k ( cosh u cosh $\phi$ - sinh u sinh $\phi$\par
-sinh u cosh $\phi$ + cosh u sinh $\phi$)\par
= k (cosh($\phi$ - u)\par
sinh($\phi$ - u)) = ( x$\prime$\par
ct$\prime$)\par
千奈「感觉，双曲角度转转了啊!」\par
广「就是那样。其实，这样我们就明白了，洛伦兹变换的真面目，是关于双曲角度\par
的转转操作。进而，试分别计算x2 - (ct)2和x$\prime$2 - (ct$\prime$)2的话，会怎样呢。 」\par
千奈「两个都会变成k2!」\par

% --- PDF page 194 ---
193\par
广 「对。 也就是说， 这个洛伦兹变换， 把x2 - (ct)2这个量持实为一定。。 这次因为\par
y = y$\prime$, z = z$\prime$成立， 所以结果x2 + y2 + z2 - (ct)2 = x$\prime$2 + y$\prime$2 + z$\prime$2 - (ct$\prime$)2成立。根据\par
这个变式， 也可以$\sqrt{}$|x2 + y2 + z2 - (ct)2|在洛伦兹变换下被持实为一定。 这与伽利略变\par
换世界里的距离相似。也就是说，从千奈看来，到以 (x, y, z) 表示的物体位置的距离\par
$\sqrt{}$x2 + y2 + z2， 在伽利略变换的世界里， 从其他任何。测者，例如从我看来， 那个物体\par
和千奈也相距$\sqrt{}$x2 + y2 + z2 ，但在洛伦兹变换里，$\sqrt{}$|x2 + y2 + z2 - (ct)2| 这个量，具\par
有和伽利略变换里的那个$\sqrt{}$x2 + y2 + z2相同的作用。 就是说， 把$\sqrt{}$|x2 + y2 + z2 - (ct)2|\par
当作洛伦兹变换里的距离般的东西来用就行。这样定义的距离，叫闵可夫斯基范数。 」\par
千奈「闵可……什么?」\par
广「闵可夫斯基。 」\par
千奈 「来来如此!就是说， 在洛伦兹变换的世界里， 不是用至今为止的$\sqrt{}$x2 + y2 + z2，\par
而是$\sqrt{}$|x2 + y2 + z2 - (ct)2|这个作为距离，这被称为闵可夫斯基范数，对吧!」\par
广「…，就是这么回事。 」\par
广 P 「在这里，来想想 光锥 外 的 任 意 点 ， 为 什 么 能 用k, $\phi$ 表 示 为(x, ct) =\par
(k cosh $\phi$ , k sinh $\phi$) 吧。如从才才的讨论，固定 k 只改变$\phi$ ，这个点就对应于表示 x2 -\par
(ct)2 = k2的任意点。这里，求这个双曲线的渐近线的话，就是x = $\pm$ct，与k无关。另\par
一方面， 改变k的话当然双曲线也变化。 也就是说，这对应于表示以x = $\pm$ct为渐近线的\par
双曲线群。不k = 0，这对应于渐近线本身，通过从此渐渐增大k，双曲线在x轴上的交\par
点从来点离开，k → $\infty$时交点行进到k → $\pm$$\infty$， 所以这些点表示光锥外的任意点。 另外，\par
在这里，因为明白了洛伦兹变换是对应于双曲角转转的操作，所以可知变换后的点，结\par
果仍为对于变换目的地的。测者坐标而言的光锥外侧的点。 」\par

% --- PDF page 195 ---
194\par
广 P 「另外，作为由此致出的性质可知，依靠洛伦兹变换，光锥内侧的点被移到光\par
锥内侧，光锥上的点被移到光锥上，光锥外侧的点被移到光锥外侧。这次就等于是展示\par
了关于光锥外侧的对应关系，但不想，虑 光锥内侧的情况时，作为 x = k sinh $\phi$ , ct =\par
k cosh $\phi$来，虑的话， 会成为(ct)2 - x2 = k2， 能表示光锥内侧的任意点， 所以试试对此\par
施行洛伦兹变换，大概就能展示了吧。 」\par
广「那么，最后，陈述一下固有时间和洛伦兹-菲茨杰拉德收缩，就结束吧。 」\par

% --- PDF page 196 ---
195\par
广「才才，我好像说了，在洛伦兹变换的世界里，用闵可夫斯基范数作为距离。那\par
么， 此前作为距离的$\sqrt{}$x2 + y2 + z2是不是就不再被使用了呢， 那倒也不是。 这次， 这个\par
会被用来测量物体的长度。那么，设设佑芽在自己的坐标系里，拿着一根棒子，其两端\par
位置以点(x$\prime$, y$\prime$, z$\prime$) = (a0, b0, c0), (a1, b1, c1) 表示。对佑芽而言这根棒子没在的，所以在\par
任意t$\prime$，两端都固定在这些位置。要求这根棒对千奈而言的长度，对千奈而言同一时刻\par
时， 两端是什么情况， 这很重要。 由洛伦兹变换， 在千奈坐标系下某时刻t = T时这些点\par
的变换目的地，将这些分别作为(x0, y0, z0), (x1, y1, z1)，各自成为这样。 」\par
(\par
x0\par
y0\par
z0\par
cT\par
) = (\par
$\gamma$ 0 0 $\gamma$$\beta$\par
0 1 0 0\par
0 0 1 0\par
$\gamma$$\beta$ 0 0 $\gamma$\par
) (\par
a0\par
b0\par
c0\par
ct$\prime$\par
) = (\par
$\gamma$a0 + $\gamma$$\beta$ct$\prime$\par
b0\par
c0\par
$\gamma$$\beta$a0 + $\gamma$ct$\prime$\par
) , (\par
x1\par
y1\par
z1\par
cT\par
) = (\par
$\gamma$a1 + $\gamma$$\beta$ct$\prime$\par
b0\par
c0\par
$\gamma$$\beta$a1 + $\gamma$ct$\prime$\par
)\par
广「这里，留意一下时间分量的话，结觉得在发生奇妙的事。 」\par
cT = $\gamma$$\beta$a0 + $\gamma$ct$\prime$ = $\gamma$$\beta$a1 + $\gamma$ct$\prime$   ???\par
广「其实，由此，在洛伦兹变换的世界里，又有一个重要的概念、「同时性的相对\par
性」被致出。 整理现在的态况， 我们在，虑对千奈而言的同一时间里， 棒的两端在哪里。\par
这里，如果对千奈而言同时的事，对佑芽而言也是同时的话，会变成怎样呢。 」\par
千奈「第 2 项和第 3 项的t$\prime$的。一样呢……」\par
广「…。但那样的话，要使等式成立，要么$\gamma$, $\beta$变成 0，要么必须a0 = a1。前者从\par
设设来看不可能，而后者则得到了事，未设设的结果，这就变得奇怪了。 」\par
千奈「诶～，已经缠在一起了啊!!」\par
广「渐渐变得复杂了呢。但是，变得奇怪的来因是什么呢。 」\par
千奈「是，把对我而言的同时，也认为对花海同学而言是同时，这一点 ……?」\par
广「…，就是那样。也就是说，对千奈而言同时的事，未必对佑芽而言也是同时，\par
就成了这么回事。因此，「t$\prime$的。一样」这个想法是错的，就是这么回事。 」\par
千奈 「诶!?但是， 那样的话，a0, b0, c0和a1, b1, c1不是也得加上什么变形才行 ……?」\par
广「冷静。因为棒对于佑芽是静止的，所以不论 t$\prime$ 是什么。，棒两端的空间坐标，\par
对佑芽而言都不变。 」\par
千奈「来来如此!那样的话，就必须区分了。 」\par
广「…。下标为 0 时作t0\par
$\prime$，为 1 时作t1\par
$\prime$的话，就成下面这样。 」\par
(\par
x0\par
y0\par
z0\par
cT\par
) = (\par
$\gamma$a0 + $\gamma$$\beta$ct0\par
$\prime$\par
b0\par
c0\par
$\gamma$$\beta$a0 + $\gamma$ct0\par
$\prime$\par
) , (\par
x1\par
y1\par
z1\par
cT\par
) = (\par
$\gamma$a1 + $\gamma$$\beta$ct1\par
$\prime$\par
b0\par
c0\par
$\gamma$$\beta$a1 + $\gamma$ct1\par
$\prime$\par
)\par

% --- PDF page 197 ---
196\par
广「便一，像这样，因为存在对自己而言是同时，但在其他坐标系里未必是同时的\par
可能性，所以，把钟放在运的的东西上，不得不在那里，虑，在那个与物体一起运的的\par
坐标系里，钟会如何前进。这叫那个物体的固有时间，以$\tau$表示。用这次的例子来说，\par
千奈在佑芽的棒两端放置时钟，各自作$\tau$0, $\tau$1，再在t = 0的瞬间使$\tau$0 = $\tau$1 = 0。在t = T\par
的时候看那些钟，则有$\tau$0 = t0\par
$\prime$ , $\tau$1 = t1\par
$\prime$。」\par
广「那么，在这个态的下，，虑左右的坐标相减的话，就成下面这样。 」\par
(\par
x1 - x0\par
y1 - y0\par
z1 - z0\par
0\par
) = (\par
$\gamma$(a1 - a0) + $\gamma$$\beta$c(t1\par
$\prime$ - t0\par
$\prime$ )\par
b1 - b0\par
c1 - c0\par
$\gamma$$\beta$(a1 - a0) + $\gamma$c(t1\par
$\prime$ - t0\par
$\prime$ )\par
)\par
广「从时间的项看，$\gamma$c(t1\par
$\prime$ - t0\par
$\prime$ ) = -$\gamma$$\beta$(a1 - a0)，所以用这个的话，关于x 分量的\par
长度，就成了下面这样。 」\par
x1 - x0 = $\gamma$(a1 - a0) - $\gamma$$\beta$2(a1 - a0) = 1 - $\beta$2\par
$\sqrt{}$1 - $\beta$2\par
(a1 - a0) = 1\par
$\gamma$ (a1 - a0)\par
广「所以，表示棒子的长度的向量，就成了下面这样。 」\par
(\par
x1 - x0\par
y1 - y0\par
z1 - z0\par
) = (\par
1 $\gamma$⁄ (a1 - a0)\par
b1 - b0\par
c1 - c0\par
)\par

% --- PDF page 198 ---
197\par
广「注意到$\beta$ < 1，所以$\gamma$ > 1，由此可知，在佑芽运的着的x方向上，棒子看起来\par
缩小到了1/$\gamma$倍。 也就是说， 物体， 在运的着的方向上， 看起来缩成了静止态的时1/$\gamma$倍\par
的长度。这个，叫洛伦兹-菲茨杰拉德收缩。 」\par
千奈「的的东西看起来是缩的啊!好不可思议啊!」\par
广「那么，狭义相对论也差不多快到高潮了。最后，来试着改写运的定律。 」\par
广 「才才， 进行了各自坐标系里物体速度w, w$\prime$之间变换的广致，当时我觉得变成了\par
一个奇怪的乱上下糟的形式。 这是因为对于。测者来说， 时间的流逝方式以及同时性发\par
生了崩坏，致『至今为止认为相同的时间变得乱上下糟了。但是，才才我们学习了将时\par
钟固定在运的物体上、 思，在该物体上的坐标系中时间如何流逝这一概念， 即固有时间\par
的概念。于是，作为速度的时间基准，不使用。测者自身的时间t ，而使用固定在物体\par
上的时钟所无示的时间$\tau$来，虑速度。这就叫做固有速度。」\par
广「把话题回到速度那时，，虑一个物体，对千奈而言，用她自身时刻t 所，虑的\par
速度为w = (wx, wy, wz)。设这个物体的固有时间为$\tau$w，固有速度为w̌。这个w̌由下式定\par
义。 」\par
w̌ : = (\par
dx/d$\tau$w\par
dy/d$\tau$w\par
dz/d$\tau$w\par
d(ct)/d$\tau$w\par
)\par
千奈「咦?为什么进了时间分量?」\par
广「这因为，是在施行洛伦兹变换时，4 维的更方一。3 维的话来本就没法变换。\par
像这样，为准备洛伦兹变换，把向量搞成 4 维的，叫四维向量。 」\par
千奈「来来如此，确实非 4 维不行呢……」\par

% --- PDF page 199 ---
198\par
广「…。那么，把w用四维向量来表示，就成下面这样。因为这是速度的四维向量\par
版本，所以就叫四维速度吧。 」\par
w = (\par
dx/dt\par
dy/dt\par
dz/dt\par
d(ct)/dt\par
) = (\par
wx\par
wy\par
wz\par
c\par
)\par
广「使用微分的性质，就能这样变形。 」\par
w̌ = dt\par
d$\tau$w\par
(\par
dx/dt\par
dy/dt\par
dz/dt\par
d(ct)/dt\par
) = dt\par
d$\tau$w\par
w\par
广「这里成问题的是，物体的固有时间与千奈的时间之间，成立着怎样的关系，觉\par
知道 。 」\par
千奈「设设物体处有。测者，由洛伦兹变换， 的，虑对那个。测者而言来点处的时\par
间就好!所以， 有\par
泽同学在，\par
泽同学的坐标是(0,0,0, $\tau$w)，从\par
泽同学看， 我正以-w\par
在移的，所以注意到$\beta$ = w/c……」\par
(\par
x\par
y\par
z\par
ct\par
) =\par
(\par
1 + ($\gamma$ - 1)wx2\par
w2\par
($\gamma$ - 1)wxwy\par
w2\par
($\gamma$ - 1)wzwx\par
w2 - $\gamma$$\beta$wx\par
w\par
($\gamma$ - 1)wxwy\par
w2 1 +\par
($\gamma$ - 1)wy2\par
w2\par
($\gamma$ - 1)wywz\par
w2 - $\gamma$$\beta$wy\par
w\par
($\gamma$ - 1)wzwx\par
w2\par
($\gamma$ - 1)wywz\par
w2 1 + ($\gamma$ - 1)wz2\par
w2 - $\gamma$$\beta$wz\par
w\par
- $\gamma$$\beta$wx\par
w - $\gamma$$\beta$wy\par
w - $\gamma$$\beta$wz\par
w $\gamma$ )\par
(\par
0\par
0\par
0\par
c$\tau$w\par
)\par
= (\par
$\gamma$wx$\tau$w\par
$\gamma$wy$\tau$w\par
$\gamma$wz$\tau$w\par
$\gamma$c$\tau$w\par
)\par
千奈「现在想要的是t与$\tau$w的关系，所以……」\par
ct = $\gamma$c$\tau$w $\Leftrightarrow$ t = $\gamma$$\tau$w\par
千奈「也就是说，dt/d$\tau$w是$\gamma$!……咦?那么复杂的矩子，竟然一下就……」\par

% --- PDF page 200 ---
199\par
广「太棒了。已经相当浸染相对论了呢，很好。如从此可知，以速度v运的的物体\par
的固有时间$\tau$v，不论其速度的方向如何，都可以用$\beta$ = v/c时的$\gamma$表示为$\tau$v = t/$\gamma$。这里\par
由于$\gamma$ > 1，$\tau$这方会比t小。也就是说，速度越大，可知时间的流逝就越慢。而且，固有\par
四维速度，可以由下面这样简单的式子表示。 」\par
w̌ = $\gamma$w\par
广 「也就是说， 只需要把普通的四维速度， 乘上由物体速度定定的$\gamma$倍就好。 那么，\par
对佑芽而言，这个物体的固有速度会怎样呢。 」\par
千奈「这，普通地做坐标变换就行呢……这次的$\gamma$, $\beta$是$\beta$ = v/c，教才才不同，所以\par
设对应于才才物体速度w的$\gamma$为$\gamma$w的话……」\par
(\par
x$\prime$\par
y$\prime$\par
z$\prime$\par
ct$\prime$\par
) = (\par
$\gamma$ 0 0 -$\gamma$$\beta$\par
0 1 0 0\par
0 0 1 0\par
-$\gamma$$\beta$ 0 0 $\gamma$\par
) (\par
$\gamma$wwx$\tau$w\par
$\gamma$wwy$\tau$w\par
$\gamma$wwz$\tau$w\par
$\gamma$wc$\tau$w\par
)\par
广「等一下。请仔细看这个矩子右边乘着的向量。 」\par
千奈「……啊!是w̌ $\tau$w!」\par
(\par
x$\prime$\par
y$\prime$\par
z$\prime$\par
ct$\prime$\par
) = (\par
$\gamma$ 0 0 -$\gamma$$\beta$\par
0 1 0 0\par
0 0 1 0\par
-$\gamma$$\beta$ 0 0 $\gamma$\par
) (\par
$\gamma$wwx$\tau$w\par
$\gamma$wwy$\tau$w\par
$\gamma$wwz$\tau$w\par
$\gamma$wc$\tau$w\par
) = $\tau$w (\par
$\gamma$ 0 0 -$\gamma$$\beta$\par
0 1 0 0\par
0 0 1 0\par
-$\gamma$$\beta$ 0 0 $\gamma$\par
) w̌\par
广「…。把这件事放在心上，广进计算看看。 」\par
千奈「我明白了。…，教才才一样，提取关于时间的项……」\par
ct$\prime$ = -$\gamma$$\beta$$\gamma$wwx$\tau$w + $\gamma$$\gamma$wc$\tau$w $\Leftrightarrow$ t$\prime$ = $\gamma$$\gamma$w (1 - $\beta$wx\par
c ) $\tau$w\par
千奈 「…因为这样， 所以使用w$\prime$的话， 对佑芽而言的固有速度w̌ $\prime$是……咦， 说起来\par
w̌ $\prime$的时间分量是……」\par
广「即使在佑芽的坐标系中表示， 根据定义也会变成c。 这次， 当用w的分量表示w$\prime$\par
时，会附带上多余的系数，所以请留意。 」\par
千奈「来来如此，那么把想求的分量设为wt$\prime$的话……」\par
$\gamma$ 1 - $\beta$2\par
1 - $\beta$$\beta$x\par
wt$\prime$ = c $\Leftrightarrow$ wt$\prime$ = $\gamma$(1 - $\beta$$\beta$x)c = -$\gamma$$\beta$wx + $\gamma$c\par

% --- PDF page 201 ---
200\par
w̌ $\prime$ = $\gamma$$\gamma$w (1 - $\beta$wx\par
c ) w$\prime$ = $\gamma$2$\gamma$w(1 - $\beta$$\beta$x) 1 - $\beta$2\par
1 - $\beta$$\beta$x\par
(\par
$\gamma$wx - $\gamma$v\par
wy\par
wz\par
-$\gamma$$\beta$wx + $\gamma$c\par
) = $\gamma$w (\par
$\gamma$wx - $\gamma$v\par
wy\par
wz\par
-$\gamma$$\beta$wx + $\gamma$c\par
)\par
千奈「……咦，这莫非是这样……?」\par
(\par
$\gamma$wx - $\gamma$v\par
wy\par
wz\par
-$\gamma$$\beta$wx + $\gamma$c\par
) = (\par
$\gamma$ 0 0 -$\gamma$$\beta$\par
0 1 0 0\par
0 0 1 0\par
-$\gamma$$\beta$ 0 0 $\gamma$\par
) (\par
wx\par
wy\par
wz\par
c\par
) = (\par
$\gamma$ 0 0 -$\gamma$$\beta$\par
0 1 0 0\par
0 0 1 0\par
-$\gamma$$\beta$ 0 0 $\gamma$\par
) w\par
广「。，注意到这点很好。那么，w̌ $\prime$会怎样呢。 」\par
千奈「代入这个的话……是这样!」\par
w̌ $\prime$ = $\gamma$w (\par
$\gamma$ 0 0 -$\gamma$$\beta$\par
0 1 0 0\par
0 0 1 0\par
-$\gamma$$\beta$ 0 0 $\gamma$\par
) w = (\par
$\gamma$ 0 0 -$\gamma$$\beta$\par
0 1 0 0\par
0 0 1 0\par
-$\gamma$$\beta$ 0 0 $\gamma$\par
) w̌\par
广「对。也就是说，要计算对佑芽而言物体的固有速度，只需把千奈的固有速度做\par
洛伦兹变换就好。因此， 这实际上就是与坐标变换相同的变换。坐标是基底的相反变化，\par
即反变， 所以这个固有速度也具有反变分量。 虽然从基底变换的角度来看确实属于反变，\par
但在物理中， 我们通常不关心基底本身，而只关注坐标之间的相对关系。由于坐标变换\par
和基底变换是同一类变换，所以称此为「共变性」。」\par
广「那么，使用这个固有速度，重新定义的量p。为把这个教此前 3 维的的量区别\par
开，就记作p̌吧。这里，$\mu$是此前的伽利略变换的世界里的m的意思。就是说，这个p̌里\par
的$\mu$是不依赖速度的常数。 」\par
p̌: = $\mu$w̌ = $\gamma$$\mu$w\par
广 「于是， 把质量重新设为m: = $\gamma$$\mu$的话， 这着着物体速度变大，$\gamma$也变大，v → c时\par
则有m → $\infty$，因此， 确实被，制在不超过光速的范围内。这个m， 就是惯性质量。 而$\mu$，\par
则对应于物体没有速度时的质量。这个，叫固有质量。 」\par
千奈「伏线不断在回收呢!」\par
广「因为这个是基于固有速度的，所以这个也具有共变性。进而，像下面这样，把\par
四维力F̌如下通过关于固有时间的微分来定义，结果也会具有共变性\par
5\par
。」\par
5\par
作者注：这里的共变性的证明，我想读到这里的各位肯定能行（本来是言定写的但因为时间关系省\par
略什么的我说不出口……） 。请加油!\par

% --- PDF page 202 ---
201\par
F̌: = dp̌\par
d$\tau$\par
广「这些与以往在伽利略变换坐标下，虑的的量、 力具有相同的基本来理， 可以同\par
样使用。」\par
千奈「这样，就完成了运的定律的改写了吧!啊，可是……」\par
广「怎么了?」\par
千奈「不，说起来，我开始在意，四维向量的时间分量在物理上意味着什么……」\par
广「…，是啊。最后，解说一下这个就结束吧。 」\par
广「虽然突然，但来，虑此前伽利略变换的世界里所用过的能量E吧。在这个世界\par
里，E 使用物体的质量$\mu$、速度w，如下定义。 」\par
E = 1\par
2 $\mu$w2\par
千奈「为什么是这个形式呢?」\par
广「这里不会深入探讨， 但其实如果对这个进行时间微分，就能看出它变成了运的\par
方程中v的表达形式。此外还有各种背景因素，致『了这样的结果。」\par
千奈「来来如此，现在我想就，接受吧。 」\par
广「把这个能量的式子，试试在伽利略变换的世界里用的量p来写的话，就成下面\par
这样。 」\par
E = ($\mu$w)2\par
2$\mu$ = p2\par
2$\mu$\par
广「那么，p̌的时间分量是$\gamma$w$\mu$c。当然，$\gamma$w是无量纲量，$\mu$有质量的量纲，c有速度\par
的量纲。另一方面，才才的E 在此之上，又被多乘了一个速度的量纲。所以，在这个分\par
量上再乘一个 c 的话，就成为$\gamma$w$\mu$c2，与能量的量纲一『。因此，就试试把这个，作为\par
这个洛伦兹变换世界里的能量Ě吧。 」\par
Ě: = $\gamma$w$\mu$c2\par
广「于是，下式就成立了。注意$\gamma$w = 1/$\sqrt{}$1 - $\beta$w2。」\par
Ě 2 = $\mu$2c4 + $\gamma$w2$\mu$2w2c2\par
广「沿着这个路径计算Ě，结果会是这样的。」\par
Ě = $\mu$c2$\sqrt{}$1 + $\gamma$w2w2\par
c2\par

% --- PDF page 203 ---
202\par
广 「这里， 当w比c非常小时会怎样呢， 把$\delta$当作非常小的量， 有(1 + $\delta$)$\alpha$ $\simeq$ 1 + $\alpha$$\delta$，\par
所以成下面这样。 」\par
Ě $\simeq$ $\mu$c2 + $\gamma$w2$\mu$w2\par
2 = $\mu$c2 + $\gamma$w2$\mu$2w2\par
2$\mu$\par
广「这里，如果用伽利略变换世界里的p，p = $\mu$w成立，所以成下面这样。 」\par
Ě $\simeq$ $\mu$c2 + $\gamma$w2p2\par
2$\mu$\par
千奈「啊!才才那个有印象的能量式子的形式出现了!」\par
广「对。也就是说，这个能量的定义，看来是对的。进而，速度w为 0 时，会怎样\par
呢。 」\par
千奈「会成这样。 」\par
Ě = $\mu$c2\par
广「对。进而 m=$\mu$也成立，所以成下面这样。 」\par
Ě = mc2\par
千奈「啊!那个有名的式子!」\par
广「这个叫静止能量，是质量成了能量，这样意思的证据式子。也就是说，物体，\par
仅凭存在，就拥有极大的能量，这么一回事。 」\par
千奈「这个式子的意义，来来是那种意思啊!」\par
广「…。那么，由此，就知道了因洛伦兹变换的形式问题而追加的时间分量，在物\par
理上表示了什么。p̌的时间分量，表示的是能量E除以光速的。，也就是E/c。」\par
千奈「来来如此!谜团解开了，清清爽爽!!」\par
广「那么，狭义相对论就到此结束吧。辛苦了。 」\par
千奈「今天也十分感谢!好开心啊～!」\par

% --- PDF page 204 ---
203\par
笔者追记\par
本来言定要到广义相对论的， 但路途变得相当遥远了， 所以这次就停留在了狭义相\par
对论。如果Vol.2 会发行的话，这之后的故事，就想在Vol.2 里执笔了!能读到这里，非\par
常感谢!\par
参考文献\par
[1]\par
(\par
),\par
(\par
)\par
,\par
, 2013, 254\par
.\par
[2] James J. Callahan (\par
)\par
(\par
).\par
.\par
, 2021,353\par
.\par
[3] FN.\par
(\par
)\par
. FN\par
.\par
: 2019/01/04. http://fnorio.com/0180covariant-\par
contravariant/covariant\_contravariant.html,\par
: 2024/10/20.\par
[4] FN."\par
".FN\par
.\par
:2018/01/19.http://fnorio.com/0171special\_general\_Lorentz\_transformation/special\_general\_Loren\par
tz\_transformation.html,\par
2024/10/20.\par
[5]\par
."\par
". EMAN\par
,\par
: 2001/12/6. https://eman-\par
physics.net/relativity/lorentz.html,\par
2024/10/20.\par
[6]\par
."E=mc2\par
".EMAN\par
.\par
:2008/1/20.https://eman-\par
physics.net/relativity/4momentum.html,\par
:2024/10/20.\par
[7]\par
,\par
(\par
),\par
,\par
(\par
).\par
=\par
(\par
6\par
.\par
,1978,468\par
.\par
使用素材\par
在撰写本文时，使用了以下网站的素材。非常感谢。\par
6\par
(1)https://sites.google.com/view/snarifont/\par
(\par
)\par
7\par
(2)https://hp.vector.co.jp/authors/VA039499/(\par
)\par
8\par
(3)https://pictogram2.com/(\par
2.0\par
topeconheroes\par
)\par
9\par
(4)https://kage-design.com/(\par
topeconheroes\par
)\par
10\par
(5)https://icooon-mono.com/(ICOOON MONO,TopeconHeroes\par
)\par
11\par
6\par
仅列出，译者并未使用。\par
7\par
Licence: https://sites.google.com/view/snarifont/ホーム/フォントのライセンス\par
8\par
Licence: https://hp.vector.co.jp/authors/VAO39499/\#KIYAKU\par
9\par
Licence: https://pictogram2.com/?page\_id=39\par
10\par
Licence: https://kage-design.com/terms-of-use/\par
11\par
Licence: https://icooon-mono.com/license/\par

% --- PDF page 205 ---
204\par
后记\par
500mL\par
Twitter: @0500mL\par
一句话： 真希望人生中能有这样的前\par
辈\par
tokio\par
Twitter: @dtokyo264196\par
一句话：\par
泽同学的眼镜，太棒了\par
Twitter: @qutum01\par
一句话：真可爱啊\par
泽。\par
Twitter: @Nina\_Shimegi\par
一句话：万圣节广抽卡大暴死。\par
latte\par
Twitter: @pdncu\par
一句话： 反而感觉对\par
泽广的理解加\par
深了\par
Twitter: @Esqrt2\par
Fediverse: @Esqrt2@misskey.design\par
一句话：因为\par
泽并非实存，所以起\par
不到\par
泽本人的解的什么的，太痛苦了，\par
救救我伏见\par
misskey.io:\par
@Noisy\_Daemon@misskey.io\par
一句话：后半段气力耗尽了。\par
Twitter: @mathviolinebook\par
一句话：被硕士论文追赶的每一天，\par
明明应该是在做最适合我的事， 却满身疮\par
痍。\par
Bluesky: @futozushi.bsky.social\par
一句话： 图老套的意外转折】 佑芽比广\par
个子更高\par
Alkyne\par
Twitter: @Alkatalynest\par
一句话：且时不想再思，相对论了，\par
不来谅爱因斯坦。\par

% --- PDF page 206 ---
\noindent\textit{[第 206 页没有可抽取文字层。]}\par

% --- PDF page 207 ---
\noindent\textit{[第 207 页没有可抽取文字层。]}\par
\endgroup
